%% The manual of GAOL v5. It follows the manual of GAOL 4 by Frederic Goualard
%% (manual/v4/gaol.tex), whose text it keeps where GAOL v5 behaves as GAOL 4,
%% and describes what GAOL v5 changes: the pages of doc/ give each change with
%% its measures. Every output shown in an example was obtained from GAOL v5.
%%
%% Copyright (c) 2026 ENSTA, France
%%
%% Created 2026-09-22 by Jordan NININ
\documentclass{manual}

\usepackage{amsmath}
\usepackage{amssymb}
\usepackage{graphicx}
\usepackage{pifont}
\usepackage{multicol}
\usepackage{array}
\usepackage{longtable}


\title{GAOL v5}
\subtitle{NOT Just Another \\Interval Arithmetic Library}
\author{Jordan Ninin\\[4pt]Fr\'ed\'eric Goualard}
\institute{ENSTA, France --- LS2N, Nantes Universit\'e, France}

\hypersetup{pdftitle={GAOL v5: NOT Just Another Interval Arithmetic Library},pdfauthor={Jordan Ninin, Frederic Goualard},pdfkeywords={constraint,interval arithmetic,IEEE 1788-2015,CORE-MATH}}

\newcommand{\itv}[2]{\ensuremath{[#1,\,#2]}}
\newcommand{\roundDn}[1]{\ensuremath{\downarrow#1\downarrow}}
\newcommand{\roundUp}[1]{\ensuremath{\uparrow#1\uparrow}}
\newcommand{\roundNearest}[1]{\ensuremath{\updownarrow#1\updownarrow}}
\newcommand{\hull}[1]{\ensuremath{\mathop{\Box} #1}}
\newcommand{\leftBound}[1]{\ensuremath{\underline{#1}}}
\newcommand{\rightBound}[1]{\ensuremath{\overline{#1}}}
\newcommand{\returns}{\hbox{\ding{229}}}
\let\leq=\leqslant
\let\geq=\geqslant
\let\emptyset=\varnothing

%% Non terminal
\newcommand*{\NT}[1]{\uppercase{\textit{#1}}}
%% Terminal
\newcommand*{\TM}[1]{\textsf{#1}}


\newcommand*{\prevfloat}[1]{\ensuremath{{#1}^-}}
\newcommand*{\nextfloat}[1]{\ensuremath{{#1}^+}}

%% What GAOL v5 changes: a mark in the text, the definitions being in
%% minipages, where no \marginpar can go
\newcommand*{\newinvfive}{\unskip\nobreak\ \textcolor{blue}{\textsf{\scriptsize[GAOL~v5]}}}
%% A name of IEEE 1788-2015
\newcommand*{\ieee}[1]{\textsf{#1}}
%% Tables of the chapter on accuracy
\newcolumntype{L}[1]{>{\raggedright\arraybackslash}p{#1}}
%% The names of functions and options, in typewriter, are not hyphenated: let
%% the lines holding them stretch rather than overflow into the margin
\setlength{\emergencystretch}{3em}

\input{gaol_version}

\begin{document}
\nobibliography*
\maketitle

\begin{copyrighttxt}
Copyright \copyright{} 2001 Swiss Federal Institute of Technology, Switzerland

Copyright \copyright{} 2002-2016 LINA UMR CNRS 6241, France

Copyright \copyright{} 2017-\number\year\ LS2N UMR CNRS 6004, France

Copyright \copyright{} 2026 ENSTA, France

All Trademarks, Copyrights and Trade Names are the property of their respective owners
even if they are not specified below.

Part of the work was done while Fr\'ed\'eric Goualard was a postdoctorate
at the \emph{Swiss Federal Institute of Technology, Lausanne, Switzerland}
supported by the \emph{European Research Consortium for Informatics and
Mathematics} fellowship programme.

This is edition \edition{} of the documentation of GAOL v5, which follows the
manual of GAOL 4 by Fr\'ed\'eric Goualard. It is consistent with version
\version{} of the \texttt{gaol} library.

Permission is granted to make and distribute verbatim copies of this manual
provided the copyright notice and this permission notice are preserved on all
copies.

\smallskip
GAOL IS PROVIDED "AS IS" WITHOUT WARRANTY OF ANY
KIND, EITHER EXPRESSED OR IMPLIED, INCLUDING, BUT NOT LIMITED TO, THE
IMPLIED WARRANTIES OF MERCHAN\-TA\-BILITY AND FITNESS FOR A PARTICULAR
PURPOSE.  THE ENTIRE RISK AS TO THE QUALITY AND PERFORMANCE OF THIS
SOFTWARE IS WITH YOU.  SHOULD THIS SOFTWARE PRO\-VE DEFECTIVE, YOU ASSUME
THE COST OF ALL NECESSARY SERVICING, REPAIR OR CORRECTION.
\end{copyrighttxt}

\newpage
\thispagestyle{empty}

\vspace*{7cm}
\marginpar{\begin{raggedright}
  \textbf{Recipriversexcluson. n.}
  \textit{A number whose existence can only be defined as being anything other than itself.} \par
\begin{flushright}
  Douglas Adams, {\fontfamily{cmr}\selectfont\oldstylenums{1982}}\\
  \textit{Life, the Universe and Everything}
\end{flushright}
\end{raggedright}}

\cleardoublepage

\tableofcontents

\chapter*{Copyright}\addcontentsline{toc}{chapter}{Copyright}
%%=================
GAOL is distributed under the GNU Library General Public License, version~2
(\pxref{chap:copying}).
The copyright for the initial version---named \code{cell}---(year 2001) is owned by the
\findex{cell}%
\emph{Swiss Federal Institute of Technology}, Lausanne, Switzerland.
The copyright for the following versions (from 2002 onward) is owned by the
\emph{Laboratoire d'Informatique de Nantes-Atlantique}, then by the
\emph{Laboratoire des Sciences du Num\'erique de Nantes}, France.

The files GAOL v5 adds carry the copyright of ENSTA, France, and are distributed
under the same license as GAOL; the files of GAOL that GAOL v5 changes keep the
copyright of their authors.

\medskip
GAOL v5 computes its elementary functions with CORE-MATH~\margincite{CoreMath:2022},
whose sources are distributed with GAOL (in \file{3rd/math-core}) and compiled
into its library. CORE-MATH is distributed under the MIT license, which is
reproduced with the license of GAOL (\pxref{chap:copying}).
\cindex{CORE-MATH}\cindex{license!of CORE-MATH}%

\mainmatter


\chapter{Introduction}
%%====================


GAOL$^\dag$\marginpar{$^\dag\:$For those readers who are not native
  English speakers, ``gaol'' should be pronounced \textsf{j\={a}l},
  like the word ``jail'', of which it is a, chiefly British, variant.}
\cindex{gaol!pronunciation}%
is a C++ library to perform arithmetic with floating-point intervals.
The development of GAOL was initiated at the \emph{Swiss Federal
  Institute of Technology}, Lausanne, Switzerland, while F.  Goualard
was a post-doctorate fellow supported by the \emph{Swiss National Science
  Foundation}.  It started as a limited version of \emph{Jail} (now
\emph{halloween}), a templated C++ interval library developed during
Goualard's PhD.

To our knowledge, a unique feature of GAOL among the C++ interval
arithmetic libraries available when it was written is the implementation of
\emph{relational arithmetic operators} required by interval constraint
arithmetic software (\pxref{sec:relational-arithmetic}).
\cindex{operator!relational}\cindex{relational operator}%
Hence, the game of the name: GAOL is \emph{not} JAIL (\emph{Just
  Another Interval Library}). The standard for interval arithmetic,
IEEE 1788-2015~\margincite{IEEE1788:2015}, has since required such operators,
as its \emph{reverse functions}.
\cindex{IEEE 1788-2015}%

\section{GAOL v5}
%%--------------

GAOL v5, written by Jordan Ninin at ENSTA, continues GAOL as Fr\'ed\'eric
Goualard developed it, from its version 4.2.2. It was written for Codac, whose
intervals are built upon GAOL, and brings the changes GAOL needs to compile and
compute right on every system it can: Linux, macOS and Windows, on x86, x86-64,
ARM, 64-bit ARM and the other processors of Debian. What it brings, and this
manual describes:
\cindex{GAOL v5}%

\begin{itemize}
\item \textbf{The tightest bounds for the elementary functions.} Every
  elementary function is bounded with CORE-MATH~\margincite{CoreMath:2022},
  which is correctly rounded in the rounding direction in effect: the value
  GAOL computes at a bound, rounding upward, \emph{is} the upper bound, and the
  double below it the lower one. The bounds are the same on every architecture
  and with every compiler, and there is no mathematical library to install,
  to find or to link along with GAOL: CORE-MATH is compiled into it.
\item \textbf{IEEE 1788-2015.} GAOL v5 follows the standard for interval
  arithmetic in the operations GAOL provides: the empty set, the infinite
  bounds, the domains of the functions, the comparisons and the interval
  literals. The few cases where it differs, mostly the \code{pow} of the
  namespace \code{gaol}, are stated where the operations are described, and
  the tightness of each operation is documented, as the standard requires
  (Chapter~\ref{chap:accuracy}).
\item \textbf{New operations}, which the standard requires or recommends:
  \code{exp2}, \code{exp10}, \code{log2}, \code{log10}, \code{fma}, the integer
  functions \code{sign}, \code{trunc}, \code{round\_ties\_to\_even} and
  \code{round\_ties\_to\_away}, fifteen recommended functions (\code{expm1},
  \code{log1p}, \code{hypot}, \code{rsqrt}, \code{sinpi}\ldots), the
  cancellative subtraction and addition, the radius, and the comparisons
  \code{less}, \code{strictly\_less}, \code{is\_entire} and
  \code{is\_common\_interval}; \code{atan2}, which GAOL declared, is
  implemented.
\item \textbf{The names of the standard.} The namespace \code{gaol\_ieee1788}
  gives each operation of the standard that GAOL provides the name and the
  argument order the standard gives it (Chapter~\ref{chap:ieee1788}).
\item \textbf{The rounding direction is the business of GAOL.} Each operation
  sets the rounding direction upward when it is not, whichever direction the
  calling code left: the user no longer has to keep it upward
  (\pxref{sec:rounding-direction}).
\item \textbf{Exact input and output.} Numbers are read exactly, and intervals
  written in decimal are rounded outward, whatever the C library; written in
  hexadecimal, an interval is read back bit for bit.
\item \textbf{Three builds that agree.} A CMake build, which is the one to use,
  and the autotools and meson builds of GAOL, which configure GAOL the same way;
  tests of the bounds GAOL computes against their exact values; a continuous
  integration on Linux, macOS and Windows, with the compilers and on the
  processors listed in Section~\ref{sec:platforms}.
\end{itemize}

\noindent
A program written for GAOL 4 needs a few changes, which
Chapter~\ref{chap:changes} lists: first of all, \file{gaol/gaol.h} no longer
opens the namespace \code{gaol}.

\section{About this manual}
%%-------------------------

This document is both a manual and a reference to use GAOL. It assumes
a prior knowledge of interval arithmetic. Refer to the books and
papers by Goldberg, Neumaier, and
others~\margincite{Goldberg:01,IEEE754:90,Moore:Book1966,Alefeld-Herzberger:83,Neumaier:Book1990}
for a basic presentation of floating-point arithmetic, interval
arithmetic and the use thereof.
\cindex{references!on interval arithmetic}

The main entry point for interval arithmetic on the Web is the \emph{Interval
Computation site} of Vladik Kreino\-vich, \url{http://www.cs.utep.edu/interval-comp/}.
\cindex{interval!computation site (web)}

The pages of the directory \file{doc} of the sources give the details of each
change GAOL v5 brings, with its measures (\pxref{chap:documentation}). The
mark \textcolor{blue}{\textsf{\scriptsize[GAOL~v5]}} follows what behaves
otherwise than in GAOL~4.

Classes, methods, functions, macros, constants and variables available in the
library but not described in this document are likely to change or to
be removed. Consequently, they should be used with caution, if at all. Among
them are the intervals of floats, \code{gaol::intervalf} and
\code{gaol::interval2f}, which are unfinished and compiled only when asked for
(\pxref{sec:cmake-options}).
\cindex{undocumented!use of \idxself{} functions}%
\cindex{function!undocumented}%

\chapter{Installation}\label{chap:installation}
%%====================

GAOL v5 is built with CMake, which is the build to use. The autotools and meson
builds of GAOL are kept for those who use them: on a given machine with a given
compiler, GAOL behaves the same whichever build configured it, the three giving
it the same macros and the same flags.

\section{Supported platforms}\label{sec:platforms}
%%---------------------------

The continuous integration of GAOL v5 builds it and runs its tests on:
\cindex{platforms}%

\begin{itemize}
\item Linux: Ubuntu 22.04, 24.04 and 26.04 on x86-64 and 64-bit ARM, with GCC and
  Clang; Debian 12 and 13 on x86-64, 64-bit ARM and 32-bit ARM (armhf), and
  Debian 12 on 32-bit x86; manylinux\_2\_28 on x86-64 and 64-bit ARM; Alpine
  (musl) on x86-64 and 64-bit ARM; Debian 13 on s390x, ppc64le and riscv64;
\item macOS 14, 15 and 26, on arm64 and x86-64, with AppleClang, the Clang of LLVM
  and GCC;
\item Windows: Visual Studio 2022 and 2026 on x86, x64 and arm64; MinGW-w64 12 to 15;
  MSYS2 (UCRT64 with GCC, CLANG64 with Clang).
\end{itemize}

\noindent
The bounds GAOL computes are the same on all of them.

\section{Getting the software}
%%----------------------------

GAOL v5 is available at \url{https://github.com/Jordan08/GAOL}.
\cindex{gaol!web page}%
The GAOL of Fr\'ed\'eric Goualard, which GAOL v5 continues, is available at
\url{https://github.com/goualard-f/GAOL}, and presented in the GAOL section
of \url{https://frederic.goualard.net}.

GAOL comes as source code, to be compiled for the computer it is used on.
The sources of CORE-MATH are part of it (\file{3rd/math-core}): there is
nothing else to download.

\section{Prerequisites}
%%---------------------

\subsection{Mandatory tools}
%%..........................
\cindex{tool!mandatory \idxself{} to compile gaol}%

\begin{itemize}
\item A C++11 and C99 compiler: GCC, Clang, AppleClang, Visual C++ (Visual
  Studio 2022 or later), or GCC and Clang for Windows (MinGW-w64, MSYS2).
\index[idx]{c++!standard}%
\item CMake 3.14 or later; or, for the autotools build, a shell and
  \cmd{make}; or, for the meson build, meson and ninja.
\end{itemize}

\noindent
No mathematical library is needed: CORE-MATH is compiled into GAOL.

\subsection{Compilers and options refused}\label{sec:refused}
%%........................................
\cindex{compiler!refused}%

Each of the following gave bounds not enclosing the exact results. The three
builds refuse the compilers when configuring, with a message naming what to use
instead, and keep the options away by giving GAOL the flags of interval
arithmetic (\pxref{sec:compiling}); \file{gaol/gaol\_config.h} refuses them
again when compiling, for the code including the headers of GAOL too, all
but the second one, which no macro of the compiler shows:

\begin{itemize}
\item Clang for 32-bit ARM processors, which does not honour the rounding
  direction there (GCC does);
\item a compiler saying of \option{-frounding-math} ``overriding currently
  unsupported rounding mode on this target'', as Clang 14 for 64-bit ARM
  (Clang 18 honours the rounding direction there);
\item a MinGW-w64 whose \file{<fenv.h>} answers \code{fegetround()} from a
  state of its own, before mingw-w64 12 on x64 and before 11 on 32-bit x86:
  GAOL sets the rounding direction by writing the registers of the processor,
  and CORE-MATH reads \code{fegetround()} to know it;
\item \option{-ffast-math}, \option{-Ofast} and \option{/fp:fast}, with which
  the compiler rounds to nearest and drops the checks of NaN and infinities;
\item Visual C++ without \option{/fp:strict}, which then assumes rounding to
  nearest;
\item doubles computed on the x87 unit of 32-bit x86 processors, in extended
  precision (without \option{-msse2 -mfpmath=sse}, or \option{/arch:SSE2}):
  CORE-MATH assumes every operation on doubles rounded to a double.
\end{itemize}

\subsection{Tools for maintainers}
%%................................

GAOL uses code produced by GNU Flex and GNU Bison for parsing the expressions used to
initialize an interval.
\findex{flex}%
\findex{bison}%
The three builds compile the lexer and the parser committed with the sources,
\file{gaol\_interval\_lexer.cpp}, \file{gaol\_interval\_parser.cpp} and
\file{gaol\_interval\_parser.h}, generated by flex 2.6.4 and bison 3.5.1 from
\file{gaol\_interval\_lexer.lpp} and \file{gaol\_interval\_parser.ypp}.
CMake and meson never generate them again. The autotools build does where one
of them is missing, as after \cmd{make maintainer-clean}, which deletes them,
or older than the file it is generated from, with \cmd{lex} and the
\cmd{yacc} that \cmd{configure} finds (\cmd{bison -y}, else \cmd{byacc} or
\cmd{yacc}). After changing one of the two files, the lexer and the parser are
generated again in \file{gaol}, and committed:

\begin{cmdshell}
% bison -y -d -p gaol_ gaol_interval_parser.ypp
% mv y.tab.c gaol_interval_parser.cpp; mv y.tab.h gaol_interval_parser.h
% flex -Pgaol_ gaol_interval_lexer.lpp
% mv lex.gaol_.c gaol_interval_lexer.cpp
\end{cmdshell}

\subsection{Optional tools}
%%.........................

\begin{itemize}
\item \cmd{doxygen}, and \cmd{dot} of GraphViz, to generate the HTML reference of
  the code (\cmd{make -C manual html} with autotools);
\index[idx]{doxygen}\findex{Graphviz}%
\item CppUnit, for the check programs of GAOL 4 (\cmd{make check} with autotools,
  \option{-Dwith-test=true} with meson); the tests of GAOL v5, built by CMake,
  need nothing (\pxref{sec:tests});
\index[idx]{CppUnit}%
\item \cmd{pdflatex}, \cmd{bibtex} and \cmd{makeindex}, to build this manual.
\end{itemize}

\section{Building and installing with CMake}\label{sec:cmake}
%%------------------------------------------
\cindex{CMake}%

In the root directory of the sources:

\begin{cmdshell}
% cmake -S . -B build -DCMAKE_INSTALL_PREFIX=@var^prefix~ \
        -DGAOL_BUILD_TESTS=ON
% cmake --build build --config Release
% ctest --test-dir build -C Release
% cmake --install build --config Release
\end{cmdshell}

\noindent
The build compiles the thirty-six sources of CORE-MATH GAOL uses into
\file{libgaol}, which is a static library: there is nothing else to build and
nothing else to install. The build type is \code{Release} unless another is
given; brought in by a project that gives none (\code{add\_subdirectory},
\code{FetchContent}), GAOL and CORE-MATH are compiled with \option{-O3} all
the same (\option{/O2} with Visual C++).

\subsection{Options}\label{sec:cmake-options}
%%..................

The options are given to the first command, as \option{-D\var{option}=\var{value}}:

\begin{itemize}
\item \option{CMAKE\_BUILD\_TYPE}. \code{Release} by default; \code{Debug}
  builds GAOL without optimization, with debugging information;
\item \option{CMAKE\_INSTALL\_PREFIX}. Where \cmd{cmake --install} puts GAOL;
\item \option{CMAKE\_POSITION\_INDEPENDENT\_CODE}. \code{ON} by default: GAOL is
  compiled as position-independent code (\option{-fPIC}), so that it can be
  linked into a shared library, as Python bindings; a project building GAOL
  with \code{FetchContent} that sets it is followed;
\item \option{GAOL\_BUILD\_TESTS}. \code{OFF} by default: builds the tests of
  \file{tests}, which \cmd{ctest} runs;
\item \option{GAOL\_SIMD}. \code{ON} by default: the intervals are computed
  with SSE2 instructions on x86 processors (\option{-msse2 -msse3}), their two
  bounds in one register; not with Visual C++ nor on 32-bit Windows;
\cindex{SSE2}%
\item \option{GAOL\_FMA}. \code{ON} by default: GAOL and CORE-MATH are compiled
  with the fused multiply-add instructions of the processor, where the compiler
  has a flag for them and the machine building runs a program compiled with it:
  \option{-mfma} on x86 (not with GCC for Windows), \option{/arch:AVX2} with
  Visual C++ for x64, \option{-mfpu=neon-vfpv4 -mfloat-abi=hard} on 32-bit ARM;
  64-bit ARM, POWER, s390x and RISC-V have them without a flag. The library then
  needs a processor with them (on x86, Intel Haswell and AMD Piledriver,
  2012--2013, and later), and the code using GAOL is given the flag too;
  \code{OFF} builds GAOL for any processor of the architecture. The elementary
  functions are much faster with them: \code{sin} and \code{cos} 2.4 times,
  \code{log} 1.6 times (Intel i7-1185G7, Clang~18.1);
\cindex{fused multiply-add}%
\item \option{GAOL\_ASM}. \code{ON} by default: GAOL's assembly code, where it
  has some (\code{GAOL\_USING\_ASM});
\item \option{GAOL\_VERBOSE\_MODE}. \code{OFF} by default: writes a line on the
  standard error when GAOL initializes and cleans up;
\cindex{verbose mode}%
\item \option{GAOL\_PRESERVE\_ROUNDING}. \code{OFF} by default: each operation
  restores the rounding direction it found, rather than leaving it upward
  (\pxref{sec:rounding-direction});
\cindex{rounding!preserving}%
\item \option{GAOL\_FLOAT\_INTERVALS}. \code{OFF} by default: compiles the
  intervals of floats, \code{gaol::intervalf}, and \code{gaol::interval2f}
  where \option{GAOL\_SIMD} gives SSE3. Both are unfinished, and left
  undocumented;
\item \option{GAOL\_COVERAGE}. \code{OFF} by default: compiles GAOL and its
  tests with the counters of gcov, and adds the target \code{coverage}, which
  runs the tests and writes the report of their coverage (GCC and Clang, with
  \code{Debug});
\item \option{GAOL\_U128\_EMULATION}. \code{OFF} by default: computes the
  128-bit integer of the accurate phases of CORE-MATH with two 64-bit halves,
  as where the compiler has no 128-bit type (Visual C++, and GCC for 32-bit
  targets).
\end{itemize}

\noindent
The CMake build always reports errors by exceptions (\pxref{sec:exceptions}).

Interprocedural optimization is not offered, and \code{gaol::gaol} carries no
such flag: with GCC, it moves floating-point operations across the changes of
the rounding direction, and GAOL then gives bounds that do not enclose the
exact results.
\cindex{optimization!interprocedural}%

\section{Building and installing with autotools}\label{sec:configuration}
%%---------------------------------------------
\findex{configure}%

In the root directory of the sources, or in another directory for a build
outside the sources:

\begin{cmdshell}
% ./configure --prefix=@var^prefix~
% make
% make install
\end{cmdshell}

\noindent
\cmd{configure} compiles the sources of CORE-MATH into \file{libgaol}, with the
flags it gives to the C code of GAOL. It is committed, generated by autoconf
2.69 from \file{configure.ac}: \cmd{make} does not regenerate it as long as
the files keep the dates of the checkout. It accepts the following options:

\begin{itemize}
\item \loption{help}. Displays a list of all options. Note that only
those described hereunder are supported;
\cindex{help!on configuration}%
\item \loption{prefix=\var{prefix-dir}}. The root directory where the library
will be installed. It defaults to \file{/usr/local};
\item \loption{libdir=\var{lib-dir}}. The directory where to put the
libraries. It defaults to \file{\var{prefix-dir}/lib};
\item \loption{includedir=\var{include-dir}}. The directory where to
put the header files. It defaults to \file{\var{prefix-dir}/include};
\item \loption{enable-optimize[=\var{yes/no}]}. Compiles GAOL and CORE-MATH
  with \option{-O3} and \code{NDEBUG}, and with the flags
  \option{-funroll-loops}, \option{-fomit-frame-pointer} and
  \option{-fexpensive-optimizations} where the compiler takes them; with
  \option{-O} otherwise. This option defaults to \code{yes};
\cindex{optimization!configuration option}%
\item \loption{enable-debug[=\var{yes/no}]}. Adds debugging information to the
  library, and enables the debugging macros (\pxref{sec:debugging-facilities}).
  This option defaults to \code{no};
\cindex{debugging!option}%
\item \loption{enable-simd[=\var{yes/no}]}. The SSE2 intervals on x86
  processors, as \option{GAOL\_SIMD} (\pxref{sec:cmake-options}). This option
  defaults to \code{yes};
\item \loption{enable-fma[=\var{yes/no}]}. The fused multiply-add instructions
  of the processor, as \option{GAOL\_FMA}. This option defaults to \code{yes};
\item \loption{enable-float-intervals[=\var{yes/no}]}. The intervals of floats,
  as \option{GAOL\_FLOAT\_INTERVALS}. This option defaults to \code{no};
\item \loption{enable-asm[=\var{yes/no}]}. GAOL's assembly code, as
  \option{GAOL\_ASM}. This option defaults to \code{yes};
\cindex{assembler}%
\item \loption{enable-verbose-mode[=\var{yes/no}]}. Writes a line on the
  standard error when GAOL initializes and cleans up. This option defaults to
  \code{no};
\cindex{verbose mode}%
\cindex{messages!avoiding printing of}%
\item \loption{enable-preserve-rounding[=\var{yes/no}]}. Each operation
  restores the rounding direction it found (\pxref{sec:rounding-direction}).
  This option defaults to \code{no};
\cindex{rounding!preserving}%
\item \loption{enable-exceptions[=\var{yes/no}]}. If enabled, errors
  are reported by throwing an exception (\pxref{sec:exceptions}); otherwise,
  by calling \code{gaol\_error()}, which prints a message to the standard error
  channel, and aborting. This option defaults to \code{yes};
\cindex{error!raising an exception}%
\item \loption{with-cppunit-include=\var{dir}}. The directory of the headers of
  CppUnit, for the check programs of \cmd{make check};
\cindex{cppunit!include path}%
\item \loption{with-cppunit-lib=\var{dir}}. The directory of the library of
  CppUnit.
\cindex{cppunit!library path}%
\end{itemize}

Note that, as usual, \loption{disable-\var{xxx}} is
equivalent to \loption{enable-\var{xxx}=no}. Moreover,
\loption{enable-\var{xxx}} is equivalent to
\loption{enable-\var{xxx}=yes}.

The targets of the \file{Makefile} in the root directory are the standard ones:
\option{all}, \option{check}, which compiles and runs the check programs of
GAOL~4 (with CppUnit), \option{install}, \option{clean}, \option{distclean} and
\option{maintainer-clean}. \cmd{make -C manual pdf} builds this manual and
the manual of GAOL~4, and \cmd{make -C manual html} the HTML reference of the
code, with doxygen.

\section{Building and installing with meson}
%%------------------------------------------
\cindex{meson}%

\begin{cmdshell}
% meson setup build --prefix=@var^prefix~
% meson compile -C build
% meson install -C build
\end{cmdshell}

\noindent
meson compiles the sources of CORE-MATH into \file{libgaol}, as CMake and
\cmd{configure} do, and the lexer and the parser committed with the sources:
neither flex nor bison is needed. \cmd{meson compile} needs meson 0.54; with an older one,
\cmd{ninja -C build} builds GAOL as well. The options are given to
\cmd{meson setup}, as \option{-D\var{option}=\var{value}}; all but the last two
are those of \cmd{configure}:

\begin{center}\renewcommand{\arraystretch}{1.1}\small
  \begin{tabular}{lL{6.2cm}}
    \hline
    \multicolumn{1}{c}{Option} & \multicolumn{1}{c}{Default}\\
    \hline\hline
    \code{buildtype} & \code{release}: \option{-O3} and \code{NDEBUG}\\
    \code{enable-optimize} & \code{true}\\
    \code{enable-debug} & \code{false}\\
    \code{enable-simd} & \code{true}\\
    \code{enable-fma} & \code{true}\\
    \code{enable-float-intervals} & \code{false}\\
    \code{enable-asm} & \code{true}\\
    \code{enable-verbose-mode} & \code{false}\\
    \code{enable-preserve-rounding} & \code{false}\\
    \code{enable-exception} & \code{true}\\
    \code{with-test} & \code{false}; builds the check programs of GAOL~4\\
    \code{with-doc} & \code{false}; the target \code{pdf} then builds this manual and the manual of GAOL~4\\
    \hline
  \end{tabular}
\end{center}

\section{Tests}\label{sec:tests}
%%-------------
\cindex{tests}%

Built with \option{GAOL\_BUILD\_TESTS}, the programs of \file{tests} compare the
bounds GAOL computes with the exact results of the operations, independently of
GAOL and of the floating-point environment: the exact results are computed with
integers, or were computed with 2000 bits of precision. They check, among
others, that the basic operations are the tightest enclosures, that the
elementary functions are the tightest or within one double of them, that every
operation gives the same result whatever the rounding direction the calling
code left, that numbers are read exactly and intervals written rounded outward,
that the reverse functions are accurate in the sense of IEEE 1788-2015, and
that each name of \code{gaol\_ieee1788} is the operation of the standard it
names. \cmd{ctest} runs them. The page \file{doc/tests.md} says what each
program checks.

\file{tests/find\_package} builds the same tests with an installed GAOL, and
\file{tests/fetch\_content} with a GAOL brought into a project by
\code{FetchContent}. \code{gaol\_performance} (\file{tests/performance.cpp})
measures the time per operation of GAOL's arithmetic and elementary functions.

\section{What is installed}
%%-------------------------

Each build installs:
\begin{itemize}
\item the headers, in \file{\var{prefix}/include/gaol};
\item the library, \file{libgaol}, which holds CORE-MATH: static with CMake;
  static and shared with autotools and meson;
\item \file{gaol.pc}, for pkg-config, in the \file{pkgconfig} directory of the
  library directory (except the CMake build with Visual C++);
\item with CMake, the CMake package of GAOL, which defines the target
  \code{gaol::gaol}.
\end{itemize}

\noindent
Both \file{gaol.pc} and \code{gaol::gaol} carry the flags that the code
including the headers of GAOL has to be compiled with (\pxref{sec:compiling}).


\chapter{An overview of GAOL}\label{chap:overview}
%%===========================

 In this chapter, we will assume that GAOL has already been
 properly installed, and that the libraries and header files are
 accessible to your compiler.

 Let us consider the following program to compute the
 range of the function
 \begin{align*}
 f(x,y)=&\bigl(1+(x+y)^2(19-14x+3x^2-14y+6xy+3y^2)\bigr)\\
 &\bigl(30+(2x-3y)^2(18-32x+12x^2+48y-36xy+27y^2)\bigr)
 \end{align*}
 for $x\in\itv{-2}{2}$ and $y\in \itv{-2}{2}$.

 \begin{example}
 #include <iostream>
 #include <gaol/gaol.h>
 using namespace gaol;

 int main(void)
 {
   interval
     x(-2,2),
     y(-2,2), z;

   z=(1+sqr(x+y)*(19-14*x+3*sqr(x)-14*y+6*x*y+3*sqr(y)))*
     (30+sqr(2*x-3*y)*(18-32*x+12*sqr(x)+
        48*y-36*x*y+27*sqr(y)));

   std::cout << "z = " << z << std::endl;
   return 0;
 }
 \end{example}

 First, note that we have to include the \file{gaol/gaol.h} header file in order to
 use all the facilities provided by GAOL. The type \code{interval} and the
 functions of GAOL are in the namespace \code{gaol\_core}, which a program
 names through the namespace \code{gaol}: the \code{using} directive of
 Line~3 imports it. The \file{gaol/gaol.h} header of GAOL~4 imported the
 namespace itself; that of GAOL v5 no longer does, and a program adds the
 directive, or names \code{gaol::interval} and \code{gaol::sqr}
 (\pxref{sec:namespaces}).\newinvfive
\cindex{gaol!namespace}%
\cindex{namespace!gaol}%
\cindex{header!and namespace|see{namespace}}%

 Nothing has to be called before using GAOL: it initializes itself when the
 program starts (\pxref{chap:init}), and each of its operations sets the
 rounding direction it needs (\pxref{sec:rounding-direction}).

 The \code{sqr(\var{x})} function stands for \emph{square of \var{x}}
 and is equivalent to \code{pow(\var{x},2)}.

 Let \file{f.cpp} be the name of the file containing the program
 above. To compile it, with the flags pkg-config gives
 (\pxref{sec:compiling}), we type the following command:
\cindex{compiling!a file using gaol}%

 \begin{cmdshell}
 % c++ -std=c++11 -O2 $(pkg-config --cflags gaol) -o f f.cpp \
       $(pkg-config --libs gaol)
 \end{cmdshell}

 \noindent
 GAOL is the only library to link, with the C mathematical library: CORE-MATH,
 which computes its elementary functions, is compiled into it.

 Executing \file{f}, we obtain:

 \begin{onscreen}
  z = [-56254330, 94177270]
 \end{onscreen}

 We then know that $f(x,y)$ ranges over
 \itv{-56254330}{94177270} when $x$ and $y$ range over \itv{-2}{2} independently.

\section{Compiling a program using GAOL}\label{sec:compiling}
%%--------------------------------------
\cindex{flags of interval arithmetic}%
\cindex{compiling!a file using gaol}%

The operations of GAOL on intervals are inline: the code that includes the
headers of GAOL has to be compiled with the \dfn{flags of interval arithmetic},
not only GAOL itself.\newinvfive{} With GCC and Clang, each where the compiler
takes it:

\begin{itemize}
\item \option{-frounding-math -fno-fast-math -ffp-contract=off}, so that the
  compiler rounds each operation in the direction in effect, as written, and
  neither evaluates it at compile time in the default rounding nor contracts a
  multiplication and an addition into a fused one;
\item \option{-msse2 -mfpmath=sse} on 32-bit x86, so that doubles are computed
  in double precision rather than on the x87 unit;
\item \option{-msse2 -msse3} on x86 processors with the SSE2 intervals
  (\option{GAOL\_SIMD});
\item \option{-ffloat-store} where doubles are still computed on the x87 unit;
\item \option{-mfma} on x86, and \option{-mfpu=neon-vfpv4 -mfloat-abi=hard} on
  32-bit ARM, where GAOL is compiled with the fused multiply-add instructions of
  the processor (\option{GAOL\_FMA}): the code using GAOL then runs only on a
  processor with them, as GAOL itself.
\end{itemize}

\noindent
With Visual C++, \option{/fp:strict}, and \option{/arch:AVX2} for x64 with
\option{GAOL\_FMA}. Each build installs these flags with GAOL, in
\code{gaol::gaol} and \file{gaol.pc}. \file{gaol/gaol\_config.h} refuses code
compiled by Visual C++ without \option{/fp:strict}. GCC and Clang do not tell
the code whether \option{-frounding-math} and \option{-ffp-contract=off} were
given: there, it only refuses what contradicts them, \option{-ffast-math} and
doubles computed on the x87 unit (\pxref{sec:refused}).

\subsection{From CMake}
%%.....................

\begin{example}
find_package(gaol REQUIRED)
target_link_libraries(my_target PRIVATE gaol::gaol)
\end{example}

\noindent
\code{gaol::gaol} carries the include directory, the flags above and, for Visual
C++, the definition \code{\_\_GAOL\_PUBLIC\_\_=}, GAOL being a static library.
There is no other library to link. A library whose headers include those of
GAOL links \code{gaol::gaol} \code{PUBLIC}, so that its own users get the
flags, and its CMake package finds GAOL again (\code{find\_dependency(gaol)}).
\cindex{CMake!using gaol from}%

A project can also build GAOL for itself, with the options it wants:

\begin{example}
include(FetchContent)
FetchContent_Declare(gaol GIT_REPOSITORY
    https://github.com/Jordan08/GAOL.git GIT_TAG master)
FetchContent_MakeAvailable(gaol)
target_link_libraries(my_target PUBLIC gaol::gaol)
\end{example}

\noindent
GAOL is then a target of the project, built with it, and \cmd{cmake --install}
of the project installs it with it; nothing is downloaded beyond the sources of
GAOL. GAOL is compiled as position-independent code, so that
\code{gaol::gaol} can be linked into a shared library.

\subsection{From pkg-config}
%%..........................

Each build installs \file{gaol.pc} in the \file{pkgconfig} directory of its
library directory, except the CMake build with Visual C++. Its \code{Cflags}
carries the flags above with the include directory, and its \code{Libs} GAOL
itself with the C mathematical library:
\cindex{pkg-config}%

\begin{cmdshell}
% export PKG_CONFIG_PATH=@var^prefix~/lib/pkgconfig
% c++ -std=c++11 -O2 $(pkg-config --cflags gaol) program.cpp \
      $(pkg-config --libs gaol)
\end{cmdshell}

\noindent
Built by CMake with its default options on x86-64 Linux, GAOL gives, on one
line:

\begin{cmdshell}
% pkg-config --cflags --libs gaol
-frounding-math -fno-fast-math -ffp-contract=off
-msse2 -msse3 -mfma -I@var^prefix~/include -L@var^prefix~/lib -lgaol -lm
\end{cmdshell}

\noindent
In a meson project, \code{dependency('gaol')} finds it.

\section{The rounding direction}\label{sec:rounding-direction}
%%------------------------------
\cindex{rounding!direction}%

 Floating-point interval arithmetic requires \dfn{outward rounding} in
 order to fulfill the \dfn{containment property}: for example, to add
 intervals \itv{a}{b} and \itv{c}{d}, we compute
 \itv{\roundDn{a+c}}{\roundUp{b+d}}, where \roundDn{r}
 and \roundUp{r} return the greatest (resp. smallest)
 floating-point number smaller (resp. greater) than the real result of
 $r$. These two operations are performed by switching the rounding
 direction of the FPU towards, respectively, $-\infty$ and
 $+\infty$.
 \marginpar{%
   \begin{definition}[Round down/up]
     \cindex{rounding!downward|hyperpagebf}\cindex{rounding!upward|hyperpagebf}%
     \cindex{F@\FSet}\cindex{R@\RSet}%
     Given \RSet\ the set of real numbers and \FSet\ the set of
     floating-point numbers (\texttt{double}), we have, for all $x\in\RSet$:
     \begin{align*}
       \roundDn{x} &= \max\{y\in\FSet\mid y\leq x\}\\
       \roundUp{x} &= \min\{y\in\FSet\mid y\geq x\}
     \end{align*}
   \end{definition}
 }
\cindex{containement property}%
\cindex{rounding!outward|see{outward rounding}}%
\cindex{downward rounding ($\roundDn{\:}$)}%
\cindex{upward rounding ($\roundUp{\:}$)}%
\cindex{outward rounding}%
\index[idx]{aaa@$\roundDn{\:}$|see{downward rounding}}%
\index[idx]{aab@$\roundUp{\:}$|see{upward rounding}}%

On most platforms, switching the rounding direction is costly.
However, it is possible to cut down the number of switches
by relying on the property that
$\roundDn{-r}=-\roundUp{r}$. Consequently, one can
replace nearly all downward rounding operations by upward rounding ones
by negating appropriately twice the operations performed. GAOL stores the
lower bound of an interval negated, and computes both bounds rounding upward:
the rounding direction can then be set towards $+\infty$ once and for all.

GAOL 4 put on the user the burden to ensure that the rounding direction be
always set towards $+\infty$ before any computation involving intervals, which
it called the \emph{trust rounding mode}. GAOL v5 no longer does:\newinvfive{}
\emph{each operation sets the rounding direction upward when it is not}, and
leaves it upward, whichever way GAOL is built. The bounds are then right
whatever rounding direction the calling code left. The check is an addition,
$1+2^{-60}$, which is above 1 only when rounded upward, rather than a reading
of the rounding direction, which cost far more on some platforms: under
Rosetta~2, reading it made an addition of intervals take 800~ns rather than
2.8~ns. Each function of intervals checks the direction once, and the
functions it calls on doubles take it to be upward already.
\cindex{trust rounding mode}%

Code that needs its own rounding direction after the operations of GAOL
builds GAOL with the option \option{GAOL\_PRESERVE\_ROUNDING} of CMake, which is
\loption{enable-preserve-rounding} with autotools and
\option{-Denable-preserve-rounding=true} with meson: each operation then also
restores the rounding direction it found, which makes the arithmetic
operations several times slower. It is off by default in the three builds.
\cindex{rounding!preserving}%

\paragraph{Note.} The rounding direction GAOL leaves applies to the rest of the
program too. With the GNU C library, \code{printf()} and the streams of C++
round the decimal digits they write of a double in the direction in effect:
after an operation of GAOL, \code{std::cout~<{}<~pi\_dn} writes
\code{3.141592653589794} with 16 digits, and \code{3.141592653589793} once
\code{round\_nearest()} has been called. The intervals themselves are written
by GAOL rounded outward, whatever the direction (\pxref{sec:writing}).

\section{The namespaces}\label{sec:namespaces}
%%----------------------
\cindex{namespace!gaol}\cindex{namespace!gaol\_core}\cindex{namespace!gaol\_ieee1788}%

The type \code{interval} of GAOL and its functions are in the namespace
\code{gaol\_core}. A program names them through one of two namespaces, which
\file{gaol/gaol.h} declares and does not open:\newinvfive

\begin{itemize}
\item \code{gaol}, GAOL under its own names, which this manual uses;
\item \code{gaol\_ieee1788}, GAOL under the names of IEEE 1788-2015
  (Chapter~\ref{chap:ieee1788}).
\end{itemize}

\noindent
The two namespaces hold \code{interval}, and the functions that have the same
name and the same meaning in both (\code{sin}, \code{exp}, \code{sqrt},
\code{min}\ldots) are the same functions, those of \code{gaol\_core}.
\code{pow} is the one whose meaning differs: each namespace has its own, and
none is in \code{gaol\_core}, the namespace where argument-dependent lookup
looks for a call \code{pow(x,~y)} on an interval, so that each namespace finds
its own only. A program opens one of the two namespaces, not both: with both
open, \code{pow(x,~y)} is ambiguous.

The \code{pow} of \code{gaol} takes the integer power for an integer exponent,
a negative base included, and the \code{pow} of IEEE 1788-2015 for any other
exponent (\pxref{sec:powers}). In \code{gaol\_core}, its versions are named
apart: \code{gaol\_pown()}, \code{gaol\_uipow()}, \code{gaol\_pow\_real()},
\code{gaol\_pow\_hybrid()}.

\section{Common errors}
%%---------------------

In this section, we will review common errors made when using GAOL in
the---forlorn?---hope that it will help prevent users from making
them.

\subsection{Floating-point arithmetic and rounding}
%%.................................................

Programming with floating-point numbers is one of the few activities
where one must always consider ones compiler defiantly. For example, let
us consider the following piece of code:

\begin{example}
#include <gaol/gaol>  @rem^// We do not import the gaol namespace~

using gaol::interval;

int main(void) {
  interval one_tenth(0.1); @rem^BEWARE: wrong !~

  @emph^[some code using one_tenth]~
}
\end{example}

\marginpar{%
\begin{definition}[Round to nearest]
  \cindex{Rounding!to nearest|hyperpagebf}
     Given \RSet\ the set of real numbers and \FSet\ the set of
     floating-point numbers (\texttt{double}), $\roundNearest{x}$ is, for
     $x\in\RSet$, the $y\in\FSet$ that minimizes $|x-y|$:
     \begin{equation*}
       |x-\roundNearest{x}| = \min\{|x-y|\mid y\in\FSet\}
     \end{equation*}
     the one whose significand is even when two are equally near.
\end{definition}
}
Though a rational perfectly representable in decimal, \code{0.1} is
not representable in binary (at least, not with a finite number of
bits) and thus requires rounding. Obviously, the purpose of the user
was to define an interval containing this value.  However, \code{0.1}
will be rounded \emph{at compile time}, most certainly to the nearest
representable floating-point number \roundNearest{0.1}.  As a
consequence, \code{one\_tenth} will be a degenerate interval
containing only \roundNearest{0.1}, and the \dfn{containment property}
will be violated.

The right way to deal with rational constants that might not be perfectly
representable as floating-point numbers is to stringify them, such
that they can be correctly rounded downward and upward at \emph{runtime}:

\begin{example}
#include <gaol/gaol>

using gaol::interval;

int main(void)
{
  interval one_tenth("0.1"), @rem^// OK: this is the right way~
           one_tenth2("1/10"); @rem^// Another possible way~

  @emph^[some code using one_tenth]~
  return 0;
}
\end{example}

\marginpar{%
  \begin{definition}[canonical interval]
    \cindex{interval!canonical|see{canonical interval}}%
    \cindex{canonical interval|hyperpagebf}%
    A non-empty interval $I=\itv{a}{b}$ is \emph{canonical} if and only if
    $a\geq\prevfloat{b}$.
  \end{definition}
}
Now, \code{one\_tenth} will be the smallest floating-point interval
enclosing \code{0.1}. An interval like this one, containing at most two
consecutive floating-point numbers, is called a \dfn{canonical interval}.
\cindex{canonical interval}%
Written in hexadecimal (\pxref{sec:output-format}), which shows the bounds
exactly, the two intervals are:

\begin{onscreen}
interval(0.1)   = [0x1.999999999999ap-4, 0x1.999999999999ap-4]
interval("0.1") = [0x1.9999999999999p-4, 0x1.999999999999ap-4]
\end{onscreen}

\subsection{Zero}
%%...............

\code{interval~x(0)} does not compile: the literal \code{0} is also a null
pointer, which the constructor from a string takes as well as the constructor
from a \code{double}, and the call is ambiguous. Write \code{interval~x(0.0)},
or \code{interval::zero()}.
\cindex{interval!zero}%

\subsection{Both namespaces}
%%..........................

A program opens the namespace \code{gaol} or the namespace
\code{gaol\_ieee1788}, not both (\pxref{sec:namespaces}):

\begin{example}
#include <gaol/gaol.h>
using namespace gaol;
using namespace gaol_ieee1788;

gaol::interval x(1,2), y = pow(x,x); @rem^// Does not compile~
\end{example}

\noindent
The compiler reports that the call of \code{pow} is ambiguous.

\subsection{The flags of interval arithmetic}
%%...........................................

Compiled without the flags of interval arithmetic (\pxref{sec:compiling}), the
inline operations of GAOL may be evaluated by the compiler at compile time,
rounded to nearest, or contracted into fused multiply-adds, and give bounds
that do not enclose the exact results. \file{gaol/gaol\_config.h} stops the
compilation where it can tell:

\begin{onscreen}
% c++ -ffast-math -c program.cpp
gaol/gaol_config.h: error: #error "GAOL cannot be
compiled with -ffast-math (nor -Ofast): the bounds it
computes would not enclose the exact results"
\end{onscreen}

\noindent
Taking the flags from \code{gaol::gaol} or \file{gaol.pc} avoids the question.

\chapter{Initialization and cleanup}\label{chap:init}
%%==================================

GAOL initializes itself when the program starts, and cleans up when it ends:
the functions below need not be called. Since Release 1.0 of GAOL, there is an
automatic initialization and cleanup that ensures that no problem will arise if
the user forgets to call them; since GAOL v5, each operation also sets the
rounding direction upward itself (\pxref{sec:rounding-direction}), so that
\code{init()} no longer has to be called again after the calling code changed
it.\newinvfive

\begin{deffun}{bool}{init}{(int \var{dbg\_lvl} = 0)}
Initializes the variable \code{debug\_level} (\pxref{sec:debugging-facilities}) to the
value of \var{dbg\_lvl}.

Unless GAOL was built to preserve the rounding direction, the first call sets the
floating-point environment to its default one, then the rounding direction
towards $+\infty$. In addition, it sets the number of digits to display for
interval bounds to 16.

Returns \code{true} if the library was not already initialized and \code{false}
otherwise. After its first call, the function sets the debugging level to its
new value and, unless the rounding direction is preserved, the rounding
direction upward again: the C runtime of Windows resets it when the program
starts, after the library initialized itself.
\end{deffun}

\begin{deffun}{bool}{cleanup}{(void)}
Frees what the initialization allocated. It does not change the floating-point
environment.

Returns \code{true} the first time it is called and \code{false} afterwards.
\end{deffun}


\begin{example}
#include <gaol/gaol.h>

int main(void)
{
  gaol::init(1); @rem^// First level of debugging requested~

  @emph^[Some code using interval arithmetic]~

  gaol::cleanup();
  return 0;
}
\end{example}

\chapter{Interval creation and assignment}
%%========================================

The methods for creating an interval and assigning a new value to an already existing one
are described in the following.

\section{Constructors}\label{sec:constructors}
%%--------------------

One can create an interval in six different ways:

\begin{itemize}
\item by providing its left and right bounds:
  \begin{example}
    interval x(@var^l~,@var^r~);
  \end{example}
  where \var{l} and \var{r} are \code{double}s or of a type that is
  castable into a \code{double}. The result is the empty set when
  \itv{l}{r} is not an interval: for a left bound $+\infty$, a right bound
  $-\infty$, bounds in the wrong order ($l>r$) and NaN bounds;\newinvfive
\cindex{empty!interval: constructing}%

\item by providing only one bound, for degenerate point intervals:
  \begin{example}
    interval x(@var^v~);
  \end{example}
  This is equivalent to: \code{interval x(\var{v},\var{v});}, and the empty set
  for an infinite or NaN \var{v};
\item without providing any bound:
  \begin{example}
    interval x;
  \end{example}
  This is equivalent to:
  \begin{example}
    interval x(-GAOL_INFINITY,+GAOL_INFINITY);
  \end{example}
  \noindent where \code{GAOL\_INFINITY} represents the infinity value of the \code{double} format
  (\pxref{sec:floating-point-constants});
\item by copying an already existing interval (\emph{copy constructor}):
  \begin{example}
    interval x(-12,12), y=x, z(x);
  \end{example}
\item by using a string representing an interval in the same format as
  the one used for input (\pxref{sec:input-format})

  \begin{example}
    interval x("[-23, inf]"),
             y("[5*0.1+dmin, 89*sinh(2.1)]");
  \end{example}

  If the input string does not comply with the expected format, the exception
  \code{gaol::input\_format\_error} is thrown if the library was compiled with
  exception support, which the CMake build always is; in the absence of exception
  support, the \code{errno} variable is set to \code{-1}, and the
  \code{gaol\_error()} function prints an error message before the program
  aborts (\pxref{chap:errors}). \code{textToInterval(\var{s})} of the namespace
  \code{gaol} is the same constructor, under the name IEEE 1788-2015 gives
  it;\newinvfive{} the \code{textToInterval()} of \code{gaol\_ieee1788} reads
  the names of the functions of the standard, and returns the empty set
  instead (Chapter~\ref{chap:ieee1788});
\item by using a string representing an interval in the same format as the one used for input
  (\pxref{sec:input-format}) for each bound: the left bound of the first one
  and the right bound of the second one are taken.

  \begin{example}
 @rem^// constructs x=[-4,4]~
 interval x("[-5,4]+1","[4,6]-[2,3]");
  \end{example}
  \code{textToInterval(\var{sl}, \var{sr})} of the namespace \code{gaol} is
  this constructor.
\end{itemize}

\paragraph{Caution.} you have to be very careful when creating an interval
from floating-point constants. Remember that a rational number that is
perfectly representable in the decimal base may require rounding in the
binary base. For example, if you write the following statement:

\begin{example}
  interval x(0.1);
\end{example}

\noindent
you will \emph{not} have created an interval containing $0.1$ since
this number has an infinite expansion in the binary base (i.e. it is
impossible to represent it perfectly whatever the size of the mantissa
may be). As a consequence, the constant $0.1$ has very likely
been rounded to the nearest floating-point number at compile-time.
In such a case, you have to use a string instead:

\begin{example}
interval x("0.1");
\end{example}

\paragraph{Note.}\label{par:inf_inf_note}
The degenerate intervals \itv{+\infty}{+\infty} and \itv{-\infty}{-\infty} do
not exist in GAOL v5.\newinvfive{} IEEE 1788-2015 has no such interval, and its
constructor fails there and gives the empty set (10.5.8, 12.12.7); the
constructors of GAOL give the empty set too, as those of IBEX do, for an
infinite point and for all the cases above:
\cindex{interval!with infinite bounds}%
\begin{example}
cout << interval(GAOL_INFINITY) << ' '
     << interval(-GAOL_INFINITY,-GAOL_INFINITY) << ' '
     << interval(2,1) << ' '
     << interval(1,GAOL_INFINITY) << endl;
@outputs^[empty] [empty] [empty] [1, inf]~
\end{example}
GAOL 4 built them as intervals of one infinite point, which its operators did not
support. The intervals with one infinite bound, as \itv{1}{+\infty}, are
intervals as any other.

Read from a string, the constant \element{inf}, an expression, stands for
\itv{\text{MAX}}{+\infty}, MAX being the largest double, and
\element{-inf} for \itv{-\infty}{-\text{MAX}}; the interval literals whose left
bound is $+\infty$ or whose right bound is $-\infty$ are the empty set, as
\code{numsToInterval} of IEEE 1788-2015 has it:
\begin{example}
cout << interval("inf") << ' ' << interval("[inf]") << ' '
     << interval("[-inf, -inf]") << endl;
@outputs^[1.797693134862315e+308, inf] [empty] [empty]~
\end{example}
The string \element{"<-inf, -inf>"} throws \code{input\_format\_error}, its two
bounds being different values (\pxref{sec:input-format}).

\section{Straight assignment}
%%---------------------------

It is possible to assign a new value to an already existing interval in
three different ways:

\begin{itemize}
\item by copying another interval:
  \begin{example}
    interval x(-12,12),
             y; @rem^// Here, y is @itv^-@infty~^+@infty~~
    y = x; @rem^// Now, y is @itv^-12~^12~~
  \end{example}
\item by using a string whose format follows the one expected for
input (\pxref{sec:reading-intervals}).

\begin{example}
interval x;
x = "[-inf, 123]";
\end{example}

\item by using a \code{double}:

\begin{example}
interval x;
x = 1234.5; @rem^// Now, x is @itv^1234.5~^1234.5~~
x = GAOL_INFINITY; @rem^// Now, x is empty~
\end{example}
\end{itemize}

\noindent
Both go through the constructors of the previous section, whose rules they
follow: assigning an infinity to an interval empties it.\newinvfive

Note that there is no method for modifying a bound of an
interval since intervals must be considered as an atomic concept.

\section{Assignment combined with an operation}
%%=============================================

The following assignment operators combine the value of the interval
pointed to by \code{self} and the value of the right-hand side
interval.
\marginpar{%
    In this manual, we note \code{self} the object to which
    \code{this} is a pointer in C++.
\cindex{this@\code{this}}\cindex{self@\code{self}}}

\begin{defmethod}{interval}{interval\&}{operator\&=}{(const interval\& \var{I})}
  \begin{operation}
    \op{$\code{self}\gets\code{self} \cap I$}
  \end{operation}
Assigns to \code{self} the interval resulting from the intersection of
\code{self} and \var{I}, which is the empty set when they are
disjoint.\newinvfive

\begin{example}
interval x(-12,12);
x &= interval(-6,23); @rem^// Now, @code^x~ is @itv^-6~^12~~
\end{example}

\end{defmethod}

\begin{defmethod}{interval}{interval\&}{operator|=}{(const interval\& \var{I})}
  \begin{operation}
    \op{$\code{self}\gets\hull{(\code{self} \cup I)}$}
  \end{operation}
  Assigns to \code{self} the interval hull of the union of \code{self} and \var{I}.

\begin{example}
  interval x(-12,12);
  x |= interval(-6,23); @rem^// Now, @code^x~ is @itv^-12~^23~~
\end{example}

\end{defmethod}

\defmethodx{interval}{interval\&}{operator+=}{(const interval\& \var{I})}
\begin{defmethod}{interval}{interval\&}{operator+=}{(double \var{d})}
  \begin{operation}
    \op{$\code{self} \gets \code{self} + I$}
    \op{$\code{self} \gets \code{self} + d$}
  \end{operation}

  Assigns to \code{self} the interval resulting from adding \code{self} and \var{I}
  (resp. \var{d}).

\begin{example}
interval x(-12,12);
x += interval(-6,23); @rem^// Now, @code^x~ is @itv^-18~^35~~
\end{example}
\end{defmethod}

\defmethodx{interval}{interval\&}{operator-=}{(const interval\& \var{I})}
\begin{defmethod}{interval}{interval\&}{operator-=}{(double \var{d})}
  \begin{operation}
    \op{$\code{self}\gets \code{self} - \var{I}$}
    \op{$\code{self}\gets \code{self} - \var{d}$}
  \end{operation}
  Assigns to \code{self} the interval resulting from subtracting \var{I} (resp. \var{d})
  from \code{self}.
\begin{example}
interval x(-12,12);
x -= interval(-6,23); @rem^// Now, @code^x~ is @itv^-35~^18~~
\end{example}
\end{defmethod}

\defmethodx{interval}{interval\&}{operator*=}{(const interval\& \var{I})}
\begin{defmethod}{interval}{interval\&}{operator*=}{(double \var{d})}
\begin{operation}
\op{$\code{self}\gets \code{self} \times\var{I}$}
\op{$\code{self}\gets \code{self} \times\var{d}$}
\end{operation}
Assigns to \code{self} the interval resulting from multiplying \code{self} and \var{I}
(resp. \var{d}).
\begin{example}
interval x(-12,12);
x *= interval(-6,23); @rem^// Now, @code^x~ is @itv^-276~^276~~
\end{example}
\end{defmethod}

\defmethodx{interval}{interval\&}{operator/=}{(const interval\& \var{I})}
\begin{defmethod}{interval}{interval\&}{operator/=}{(double \var{d})}
\begin{operation}
\op{$\code{self}\gets \code{self} / \var{I}$}
\op{$\code{self}\gets \code{self} / \var{d}$}
\end{operation}
Assigns to \code{self} the interval resulting from dividing \code{self} by \var{I}
(resp. \var{d}).
\begin{example}
interval x(-12,12);
x /= interval(-6,23);   @rem^// Now, @code^x~ is @itv^-@infty~^+@infty~~
x /= interval::zero();  @rem^// Now, @code^x~ is @ensuremath^@emptyset~~
\end{example}
\end{defmethod}

\defmethodx{interval}{interval\&}{operator\%=}{(const interval\& \var{I})}
\begin{defmethod}{interval}{interval\&}{operator\%= }{(double \var{d})}
Assigns to \code{self} the interval resulting from dividing \code{self} by
\var{I} (resp. \var{d}), using a \emph{relational division} (\pxref{sec:relational-arithmetic}).
\begin{example}
interval x(-12,12);
x %= interval(-6,23);   @rem^// Now, @code^x~ is @itv^-@infty~^+@infty~~
x %= interval::zero();  @rem^// Now, @code^x~ is @itv^-@infty~^+@infty~~
\end{example}
\end{defmethod}

\noindent
The operators with a double compute as those with an interval, by the sign of
the double; an infinite double gives the empty set, as \code{interval(\var{d})}
is empty.\newinvfive

\chapter{Interval constants}\label{sec:interval-constants}
%%==========================
For convenience, some useful intervals and some canonical intervals enclosing real constants
are defined as static functions of the \code{interval} class:


\begin{center}\renewcommand{\arraystretch}{1.1}
  \begin{tabular}{lc}
    \hline
    \multicolumn{1}{c}{Function} & \multicolumn{1}{c}{Value}\\
    \hline\hline
    \element{interval::emptyset()} & $\emptyset$\\
    \element{interval::half\_pi()} & \itv{\roundDn{\frac{\pi}{2}}}{\roundUp{\frac{\pi}{2}}}\\
    \element{interval::minus\_one\_plus\_one()} & \itv{-1}{1}\\
    \element{interval::negative()} & \itv{-\infty}{0}\\
    \element{interval::one()} & $1$\\
    \element{interval::one\_plus\_infinity()} & \itv{1}{+\infty}\\
    \element{interval::pi()} & \itv{\roundDn{\pi}}{\roundUp{\pi}}\\
    \element{interval::positive()} & \itv{0}{+\infty}\\
    \element{interval::two\_pi()} & \itv{\roundDn{2\pi}}{\roundUp{2\pi}}\\
    \element{interval::universe()} & \itv{-\infty}{+\infty}\\
    \element{interval::zero()} & $0$\\
    \hline
  \end{tabular}
\end{center}

\noindent
The enclosures of $\pi$, $2\pi$ and $\pi/2$ are the tightest ones. The empty set
is held as an interval whose two bounds are NaN.
\cindex{empty!interval: representation}%

\begin{example}
cout << interval::emptyset() << ' ' << interval::pi() << endl;
@outputs^[empty] [3.141592653589793, 3.141592653589794]~
\end{example}


\chapter{Interval relations}\label{sec:interval-relations}
%===========================

Interval relations may be divided into two groups. Given $I$ and
$J$ two intervals, we have:

\begin{enumerate}
\item\dfn{set relations}: intervals $I$ and $J$ are considered as sets of reals. For example:
\begin{equation*}
I=J \Leftrightarrow(\forall x\in I, \exists y\in J\colon x=y)\wedge
                   (\forall y\in J, \exists x\in I\colon x=y)
\end{equation*}
Basically, two intervals are equal in that mode if they have the same bounds;
\item\dfn{certainly relations}: the relations must be true for any tuple of
values in the intervals. For example:
\begin{equation*}
I\leq J \Leftrightarrow(\forall x\in I, \forall y\in J\colon x\leq y)
\end{equation*}
Then, $I$ is less than or equal to $J$ in that mode if its right bound is at
most the left bound of $J$.
\end{enumerate}

GAOL~4 had a third group, the \dfn{possibly relations}, true if at least one
tuple of values verifies the real relation, and a certain equality,
\code{certainly\_eq()} and \code{certainly\_neq()}. They are gone from
GAOL v5, and so are the operators \code{==} and \code{!=} on
intervals:\newinvfive{} \code{certainly\_neq()} was the negation of
\code{certainly\_eq()}, which is ``possibly not equal'', so that \itv{3}{4} was
certainly not equal to \itv{3}{4}. IEEE 1788-2015 has neither. Its equality is
\ieee{equal}, \code{set\_eq()}; that no value of $I$ equals a value of $J$ is
\ieee{disjoint}, \code{set\_disjoint()}; and ``possibly equal'' is
\code{!set\_disjoint()}.
\cindex{relation!possibly (removed)}%

The relation symbols $<$, $\leq$, $>$ and $\geq$ are the certainly relations
in every build (\pxref{sec:relational-symbols}).

The relations that IEEE 1788-2015 defines (Tables 10.3 and 10.4) follow its
definitions, for the empty set and for the infinite bounds
too:\newinvfive{} \ieee{equal}, \ieee{subset}, \ieee{interior},
\ieee{disjoint}, \ieee{precedes}, \ieee{strictPrecedes}, \ieee{less} and
\ieee{strictLess}. Each method below says which one it is; the namespace
\code{gaol\_ieee1788} gives them their names (Chapter~\ref{chap:ieee1788}).
\cindex{IEEE 1788-2015!relations}%

\section{Set relations}
%----------------------

\defmethodx{interval}{bool}{set\_contains}{(const interval\& \var{I}) const}
\begin{defmethod}{interval}{bool}{set\_contains}{(double \var{d}) const}
  \begin{operation}
    \op{$I\subseteq\code{self}$}
    \op{$d\in\code{self}$}
  \end{operation}
  Returns \code{true} if \var{I} (resp. \{\var{d}\}) is included in \code{self}:
  \ieee{subset}(\var{I}, \code{self}), and \ieee{isMember}(\var{d}, \code{self}) of
  IEEE 1788-2015 for a finite \var{d}. \var{d} shall not be a NaN.

\begin{example}
interval x(-12,34), y(-12,5);

cout << x.set_contains(y) << ' '
     << x.set_contains(interval::emptyset()) << ' '
     << y.set_contains(x) << ' '
     << interval::emptyset().set_contains(x) << ' '
     << interval::emptyset().set_contains(interval::emptyset())
     << endl;
@outputs^true true false false true~
\end{example}
\end{defmethod}

\defmethodx{interval}{bool}{set\_strictly\_contains}{(const interval\& \var{I}) const}
\begin{defmethod}{interval}{bool}{set\_strictly\_contains}{(double \var{d}) const}
  \begin{operation}
    \op{$I$ interior to \code{self}}
    \op{$d$ in the interior of \code{self}}
  \end{operation}
Returns \code{true} if \var{I} is interior to \code{self}, \ieee{interior}(\var{I},
\code{self}) of IEEE 1788-2015: each bound of \code{self} is beyond that of \var{I},
an infinite bound being beyond the same infinite bound, and the empty set is
interior to any interval, itself included.\newinvfive{} The second form returns
\code{true} if \var{d} is strictly between the bounds of \code{self}.

\begin{example}
interval x(-10,12), y(-10, 11), z, t, u(10.5);

cout << boolalpha
     << x.set_strictly_contains(y) << ' '
     << z.set_strictly_contains(t) << ' '
     << x.set_strictly_contains(u) << ' '
     << interval::emptyset().set_strictly_contains(
                                 interval::emptyset()) << ' '
     << u.set_strictly_contains(interval::emptyset())
     << endl;
@outputs^false true true true true~
\end{example}
\end{defmethod}

\begin{defmethod}{interval}{bool}{set\_disjoint}{(const interval\& \var{I}) const}
  \begin{operation}
    \op{$\code{self}\cap I=\emptyset$}
  \end{operation}
Returns \code{true} if the intersection of \code{self} and \var{I} is empty:
\ieee{disjoint} of IEEE 1788-2015.

\begin{example}
interval a(2,4), b(6,8);
cout << a.set_disjoint(b) << ' '
     << interval::emptyset().set_disjoint(interval::emptyset());
@outputs^true true~
\end{example}
\end{defmethod}

\begin{defmethod}{interval}{bool}{set\_eq }{(const interval\& \var{I}) const}
  \begin{operation}
\op{$\forall x\in \code{self} \exists y\in I\colon x=y \wedge \forall y\in I, \exists x\in \code{self}\colon y=x$}
  \end{operation}
Returns \code{true} if intervals \code{self} and \var{I} are equal
when considered as sets of reals: \ieee{equal} of IEEE 1788-2015.
\begin{example}
cout << interval(4,8).set_eq(interval(4,8)) << ' '
     << interval::emptyset().set_eq(interval::emptyset()) << endl;
@outputs^true true~
\end{example}
\end{defmethod}

\begin{defmethod}{interval}{bool}{set\_neq}{(const interval\& \var{I}) const}
    \begin{operation}
      \op{$\neg(\code{self}=I)$}
    \end{operation}
    Returns \code{true} if \code{self} and \var{I} are not equal when
    considered as sets of reals.
    \begin{example}
cout << interval::universe().set_neq(interval::emptyset());
@outputs^true~
    \end{example}
\end{defmethod}

\begin{defmethod}{interval}{bool}{set\_le}{(const interval\& \var{I}) const}
  \begin{operation}
    \op{\code{self} interior to $I$}
  \end{operation}
Returns \code{true} if \code{self} is interior to \var{I}:
\code{\var{I}.set\_strictly\_contains(self)}, \ieee{interior}(\code{self}, \var{I}) of
IEEE 1788-2015.\newinvfive
\begin{example}
cout << interval(-4.5,3).set_le(interval(-10,10)) << ' '
     << interval::emptyset().set_le(interval(5,6)) << ' '
     << interval::emptyset().set_le(interval::emptyset()) << ' '
     << interval::universe().set_le(interval::universe()) << ' '
     << interval("[2, inf]").set_le(interval("[1, inf]"));
@outputs^true true true true true~
\end{example}
\end{defmethod}

\begin{defmethod}{interval}{bool}{set\_leq}{(const interval\& \var{I}) const}
  \begin{operation}
    \op{$\code{self}\subseteq I$}
  \end{operation}
Returns \code{true} if the real set defined by \code{self} is included in \var{I}:
\ieee{subset}(\code{self}, \var{I}) of IEEE 1788-2015.
\begin{example}
cout << interval(4.5,6).set_leq(interval(4.5,6)) << ' '
     << interval(3.5,9).set_leq(interval(2,6)) << ' '
     << interval::emptyset().set_leq(interval::emptyset());
@outputs^true false true~
\end{example}
\end{defmethod}

\begin{defmethod}{interval}{bool}{set\_ge}{(const interval\& \var{I}) const}
  \begin{operation}
    \op{$I$ interior to \code{self}}
  \end{operation}
Returns \code{true} if \var{I} is interior to \code{self}: \code{set\_strictly\_contains(\var{I})}.
\begin{example}
cout << interval(-10,10).set_ge(interval(-4.5,3)) << ' '
     << interval(1,2).set_ge(interval(1,2));
@outputs^true false~
\end{example}
\end{defmethod}

\begin{defmethod}{interval}{bool}{set\_geq}{(const interval\& \var{I}) const}
  \begin{operation}
    \op{$I\subseteq\code{self}$}
  \end{operation}
Returns \code{true} if the real set defined by \code{self} contains \var{I}:
\code{set\_contains(\var{I})}.
\end{defmethod}


\section{Certainly relations}
%%---------------------------

\begin{defmethod}{interval}{bool}{certainly\_le}{(const interval\& \var{I}) const}
  \begin{operation}
    \op{$\forall x\in \code{self}, \forall y \in \var{I}\colon x<y$}
  \end{operation}
  Returns \code{true} if \code{self} is certainly strictly less than \var{I}:
  \ieee{strictPrecedes} of IEEE 1788-2015, true when either interval is
  empty.\newinvfive
  \begin{example}
cout << interval(4,5).certainly_le(interval(6,9)) << ' '
     << interval(4,5).certainly_le(interval(5,9)) << ' '
     << interval::emptyset().certainly_le(interval(4,6)) << ' '
     << interval(4,6).certainly_le(interval::emptyset());
@outputs^true false true true~
  \end{example}
\end{defmethod}

\begin{defmethod}{interval}{bool}{certainly\_leq}{(const interval\& \var{I}) const}
  \begin{operation}
    \op{$\forall x\in \code{self}, \forall y \in \var{I}\colon x\leq y$}
  \end{operation}
  Returns \code{true} if \code{self} is certainly less or equal to \var{I}:
  \ieee{precedes} of IEEE 1788-2015, true when either interval is
  empty.\newinvfive
  \begin{example}
cout << interval(4,5).certainly_leq(interval(6,9)) << ' '
     << interval(4,5).certainly_leq(interval(5,9)) << ' '
     << interval(5,9).certainly_leq(interval(4,5)) << ' '
     << interval(4,8).certainly_leq(interval(5,9)) << ' '
     << interval::emptyset().certainly_leq(interval(4,6)) << ' '
     << interval(4,6).certainly_leq(interval::emptyset());
@outputs^true true false false true true~
\end{example}
\end{defmethod}

\begin{defmethod}{interval}{bool}{certainly\_ge}{(const interval\& \var{I}) const}
  \begin{operation}
    \op{$\forall x\in \code{self}, \forall y \in \var{I}\colon x>y$}
  \end{operation}
  Returns \code{true} if \code{self} is certainly strictly greater than \var{I},
  or if either interval is empty.
  \begin{example}
cout << interval(8,10).certainly_ge(interval(4,8)) << ' '
     << interval(8,10).certainly_ge(interval(4,7)) << ' '
     << interval::emptyset().certainly_ge(interval::emptyset());
@outputs^false true true~
  \end{example}
\end{defmethod}

\begin{defmethod}{interval}{bool}{certainly\_geq}{(const interval\& \var{I}) const}
  \begin{operation}
    \op{$\forall x\in \code{self}, \forall y \in \var{I}\colon x\geq y$}
  \end{operation}
  Returns \code{true} if \code{self} is certainly greater or equal to \var{I},
  or if either interval is empty.
  \begin{example}
cout << interval(8,10).certainly_geq(interval(4,8)) << ' '
     << interval::emptyset().certainly_geq(interval::emptyset());
@outputs^true true~
  \end{example}
\end{defmethod}

\begin{defmethod}{interval}{bool}{certainly\_positive}{(void) const}
  \begin{operation}
    \op{$\code{self}\subseteq\itv{0}{+\infty}$}
  \end{operation}
Returns \code{true} if the left bound of \code{self} is greater or equal to zero,
or if \code{self} is empty.
\begin{example}
cout << interval::emptyset().certainly_positive() << ' '
     << interval(4,5).certainly_positive() << ' '
     << interval(-0.0,6).certainly_positive() << ' '
     << interval(-6,0).certainly_positive();
@outputs^true true true false~
\end{example}
\end{defmethod}

\begin{defmethod}{interval}{bool}{certainly\_strictly\_positive}{(void) const}
  \begin{operation}
    \op{$\code{self}\subseteq\left]0,+\infty\right]$}
  \end{operation}
Returns \code{true} if the left bound of \code{self} is strictly greater than zero,
or if \code{self} is empty.
\begin{example}
cout << interval::emptyset().certainly_strictly_positive() << ' '
     << interval(4,5).certainly_strictly_positive() << ' '
     << interval(-0.0,6).certainly_strictly_positive() << ' '
     << interval(-6,0).certainly_strictly_positive();
@outputs^true true false false~
\end{example}
\end{defmethod}

\begin{defmethod}{interval}{bool}{certainly\_negative}{(void) const}
  \begin{operation}
    \op{$\code{self}\subseteq\itv{-\infty}{0}$}
  \end{operation}
Returns \code{true} if the right bound of \code{self} is lower or equal to zero,
or if \code{self} is empty.
\begin{example}
cout << interval::emptyset().certainly_negative() << ' '
     << interval(4,5).certainly_negative() << ' '
     << interval(-6,0).certainly_negative() << ' '
     << interval(-6,-5).certainly_negative();
@outputs^true false true true~
\end{example}
\end{defmethod}

\begin{defmethod}{interval}{bool}{certainly\_strictly\_negative}{(void) const}
  \begin{operation}
    \op{$\code{self}\subseteq\left[-\infty,0\right[$}
  \end{operation}
Returns \code{true} if the right bound of \code{self} is strictly lower than zero,
or if \code{self} is empty.
\begin{example}
cout << interval::emptyset().certainly_strictly_negative() << ' '
     << interval(4,5).certainly_strictly_negative() << ' '
     << interval(-6,0).certainly_strictly_negative() << ' '
     << interval(-6,-5).certainly_strictly_negative();
@outputs^true false false true~
\end{example}
\end{defmethod}

\section{The order of IEEE 1788-2015}
%%-----------------------------------

\begin{defmethod}{interval}{bool}{less}{(const interval\& \var{I}) const}
  \begin{operation}
    \op{$\leftBound{\code{self}}\leq\leftBound{I}\wedge\rightBound{\code{self}}\leq\rightBound{I}$}
  \end{operation}
  Returns \code{true} if each bound of \code{self} is at most the bound of
  \var{I} on the same side: \ieee{less} of IEEE 1788-2015 (Table 10.3). Two
  empty sets are in the relation, and an empty set and a nonempty interval are
  not (Table 10.4).\newinvfive
\begin{example}
interval e = interval::emptyset();
cout << interval(1,2).less(interval(1,3)) << ' '
     << interval(1,4).less(interval(2,3)) << ' '
     << e.less(e) << ' ' << e.less(interval(1,2));
@outputs^true false true false~
\end{example}
\end{defmethod}

\begin{defmethod}{interval}{bool}{strictly\_less}{(const interval\& \var{I}) const}
  \begin{operation}
    \op{$\leftBound{\code{self}}<\leftBound{I}\wedge\rightBound{\code{self}}<\rightBound{I}$}
  \end{operation}
  Returns \code{true} if each bound of \code{self} is below the bound of \var{I}
  on the same side, an infinite bound counting as below itself:
  \ieee{strictLess} of IEEE 1788-2015. The empty sets are as for
  \code{less()}.\newinvfive
\begin{example}
cout << interval(1,2).strictly_less(interval(1,3)) << ' '
     << interval("[-inf, 1]").strictly_less(interval("[-inf, 2]"))
     << ' ' << interval(0,1).strictly_less(interval(2,3));
@outputs^false true true~
\end{example}
\end{defmethod}

\section{Relational Symbols}\label{sec:relational-symbols}
%%--------------------------

The relation symbols are the certainly relations, whichever build configured
GAOL: the option of GAOL~4 that chose between the set, certainly and possibly
relations is gone, and \cmd{configure} and meson refuse
it.\newinvfive{} \code{==} and \code{!=} are not defined on intervals: compare
them with \code{set\_eq()}, \code{set\_disjoint()} or \code{less()}.
\cindex{relation!symbols}%

\begin{deffun}{bool}{operator<}{(const interval\& \var{I1}, const interval\& \var{I2})}
  Returns \code{\var{I1}.certainly\_le(\var{I2})}: \ieee{strictPrecedes} of
  IEEE 1788-2015.
\end{deffun}

\begin{deffun}{bool}{operator<=}{(const interval\& \var{I1}, const interval\& \var{I2})}
  Returns \code{\var{I1}.certainly\_leq(\var{I2})}: \ieee{precedes} of
  IEEE 1788-2015.
\end{deffun}

\begin{deffun}{bool}{operator>}{(const interval\& \var{I1}, const interval\& \var{I2})}
  Returns \code{\var{I1}.certainly\_ge(\var{I2})}.
\end{deffun}

\begin{deffun}{bool}{operator>=}{(const interval\& \var{I1}, const interval\& \var{I2})}
  Returns \code{\var{I1}.certainly\_geq(\var{I2})}.
\begin{example}
cout << (interval(1,2) < interval(3,4)) << ' '
     << (interval(1,2) <= interval(2,4)) << ' '
     << (interval(1,3) <= interval(2,4)) << ' '
     << (interval(3,4) > interval(1,2));
@outputs^true true false true~
\end{example}
\end{deffun}


\section{Interval-specific relations}
%%-----------------------------------

\begin{defmethod}{interval}{bool}{straddles\_zero}{(void) const}
  \begin{operation}
    \op{$0 \in \code{self}$}
  \end{operation}
Returns \code{true} if \code{self} contains zero.
\begin{example}
interval x(0,4), y, z(-12,-5);

cout << boolalpha << x.straddles_zero() << ' '
     << y.straddles_zero() << ' '
     << z.straddles_zero() << endl;
@outputs^true true false~
\end{example}

\paragraph*{Note.} \code{I.straddles\_zero()} $\equiv$ \code{I.set\_contains(0.0)}
\end{defmethod}

\begin{defmethod}{interval}{bool}{strictly\_straddles\_zero}{(void) const}
  \begin{operation}
  \op{$\{0\}$ interior to \code{self}}
  \end{operation}
  Returns \code{true} if zero is included in the interior of \code{self}.
\begin{example}
interval x(0,4), y, z(-12,-5);

cout << boolalpha << x.strictly_straddles_zero() << ' '
     << y.strictly_straddles_zero() << ' '
     << z.strictly_straddles_zero() << endl;
@outputs^false true false~
\end{example}

\paragraph*{Note.} The method call to \code{I.strictly\_straddles\_zero()} is equivalent to
                   \code{I.set\_strictly\_contains(0.0)}
\end{defmethod}


\begin{defmethod}{interval}{bool}{is\_a\_double}{(void) const}
Returns \code{true} whenever the left and right bounds of the interval
are equal: \ieee{isSingleton} of IEEE 1788-2015.
\begin{example}
interval x(-12.5), y;

cout << boolalpha << x.is_a_double() << ' '
     << y.is_a_double() << ' '
     << interval::emptyset().is_a_double() << endl;
@outputs^true false false~
\end{example}
\end{defmethod}

\begin{defmethod}{interval}{bool}{is\_an\_int}{(void) const}
  Returns \code{true} whenever the left and right bounds of the interval
  are equal and castable into an integer (type \code{int}).
\begin{example}
interval x(-12.0), y;

cout << boolalpha << x.is_an_int() << ' ' << y.is_an_int() << ' '
     << interval::emptyset().is_an_int() << ' '
     << interval(3e10).is_an_int() << endl;
@outputs^true false false false~
\end{example}
\end{defmethod}


\begin{defmethod}{interval}{bool}{is\_canonical}{(void) const}
  Returns \code{true} if \code{self} is nonempty and contains at most two
  floating-point numbers.
\begin{example}
interval x(0.0), y(-14,9);

cout << boolalpha << x.is_canonical() << ' '
     << y.is_canonical() << ' '
     << interval::emptyset().is_canonical() << ' '
     << interval::pi().is_canonical();
@outputs^true false false true~
\end{example}
\end{defmethod}

\begin{defmethod}{interval}{bool}{is\_empty}{(void) const}
  \begin{operation}
    \op{$\code{self} = \emptyset$}
  \end{operation}
  Returns \code{true} if \code{self} is an empty interval: \ieee{isEmpty} of
  IEEE 1788-2015.
  \begin{example}
cout << interval(4,5).is_empty() << ' '
     << interval(5,4).is_empty() << ' '
     << interval::emptyset().is_empty();
@outputs^false true true~
  \end{example}
\end{defmethod}

\begin{defmethod}{interval}{bool}{is\_entire}{(void) const}
  \begin{operation}
    \op{$\code{self} = \itv{-\infty}{+\infty}$}
  \end{operation}
  Returns \code{true} if \code{self} is the whole real line: \ieee{isEntire} of
  IEEE 1788-2015.\newinvfive
  \begin{example}
cout << interval::universe().is_entire() << ' '
     << interval::positive().is_entire() << ' '
     << interval::emptyset().is_entire();
@outputs^true false false~
  \end{example}
\end{defmethod}

\begin{defmethod}{interval}{bool}{is\_common\_interval}{(void) const}
  Returns \code{true} if \code{self} is nonempty and bounded:
  \ieee{isCommonInterval} of IEEE 1788-2015.\newinvfive
  \begin{example}
cout << interval(1,2).is_common_interval() << ' '
     << interval::positive().is_common_interval() << ' '
     << interval::emptyset().is_common_interval();
@outputs^true false false~
  \end{example}
\end{defmethod}

\begin{defmethod}{interval}{bool}{is\_zero}{(void) const}
  \begin{operation}
    \op{$\code{self} = \itv{0}{0}$}
  \end{operation}
  Returns \code{true} if \code{self} is equal to the interval containing only $0$.
  \begin{example}
cout << interval::zero().is_zero() << ' '
     << interval(0.0,0.0).is_zero() << ' '
     << interval(-0.0,+0.0).is_zero() << ' '
     << interval(0,5).is_zero() << ' '
     << interval::emptyset().is_zero();
@outputs^true true true false false~
  \end{example}
\end{defmethod}

\begin{defmethod}{interval}{bool}{is\_symmetric}{(void) const}
  Returns \code{true} if the left bound of \code{self} is the opposite of the
  right bound. An empty interval is not symmetric.

\begin{example}
cout << interval(-5,5).is_symmetric() << ' '
     << interval::emptyset().is_symmetric() << endl;
@outputs^true false~
\end{example}
\end{defmethod}

\begin{defmethod}{interval}{bool}{is\_finite}{(void) const}
  Returns \code{true} if neither bound is infinite. The empty set, whose bounds
  are NaN, is finite; \code{is\_common\_interval()} is not true of it.

\begin{example}
cout << interval("[4,inf]").is_finite() << ' '
     << interval::emptyset().is_finite() << ' '
     << interval(5,80).is_finite();
@outputs^false true true~
\end{example}
\end{defmethod}


\chapter{Interval Arithmetic}\label{chap:arithmetic}
%%===========================

The \emph{containment principle} of (floating-point) interval arithmetic
imposes that for any operation ``$\circ$'', and any intervals \var{I} and \var{J},
the following does hold:
\begin{equation*}
I\circ J= \hull{\{i\circ j\mid i\in I, j\in J\}}
\end{equation*}
where \hull{} is a function mapping any real set to the
smallest floating-point interval containing it.

For example, if we consider the interval square root, we have:
\begin{equation*}
\sqrt{\var{I}} = \hull{\{\sqrt{i}\mid i\in \var{I}\}}
\end{equation*}

From monotonicity considerations, the square root of \itv{1}{2} is then
$\sqrt{\itv{1}{2}}=\itv{1}{\sqrt{2}}$. Now, another
interpretation of the square root function is as follows:

\begin{equation*}
\sqrt{_r\var{I}} = \hull{\{j\in\RSet\mid\exists i\in \var{I}\colon j^2=i\}}
\end{equation*}

This last definition stands for the \dfn{relational square root} and
permits obtaining both negative and positive values. Hence, we have:
\begin{equation*}
\sqrt{_r\itv{1}{2}} = \itv{-\sqrt{2}}{\sqrt{2}}
\end{equation*}

This operator arises when we consider the relation
\begin{equation*}
x^2= y
\end{equation*}
which can alternatively be written
\begin{equation*}
x=\sqrt{_r y}
\end{equation*}
Here, the \dfn{functional square root} is not suitable since it would
induce the intersection of the domain of $y$ with \itv{0}{+\infty}.

Some applications (mainly in the area of \emph{constraint programming})
require the availability of such operators. As a consequence, GAOL offers
both functional and relational versions of the main arithmetic operators.

\paragraph{Domains.} A function is computed on the part of its argument within its
domain, which is that of IEEE 1788-2015 (Tables 9.1 and 10.5): the part outside
is left out, and an argument with no point of the domain gives the empty set,
as \code{log(\itv{-4}{0})} does.\newinvfive{} Every function gives the empty set
for an empty argument. How tight each result is, for each operation, is stated
in Chapter~\ref{chap:accuracy}.
\cindex{domain of a function}%

\section{Functional Arithmetic}
%%-----------------------------

\subsection{Basic operations}
%%...........................

\defmethodx{interval}{interval}{operator+}{(void) const}
\deffunx{interval}{operator+}{(const interval\& \var{I}, double \var{d})}
\deffunx{interval}{operator+}{(double \var{d}, const interval\& \var{I})}
\begin{deffun}{interval}{operator+}{(const interval\& \var{I1}, const interval\& \var{I2})}
  Addition of two intervals, or of one interval and a double: \ieee{add} of
  IEEE 1788-2015.
\end{deffun}

\defmethodx{interval}{interval}{operator-}{(void) const}
\deffunx{interval}{operator-}{(const interval\& \var{I}, double \var{d})}
\deffunx{interval}{operator-}{(double \var{d}, const interval\& \var{I})}
\begin{deffun}{interval}{operator-}{(const interval\& \var{I1}, const interval\& \var{I2})}
  Negation (\ieee{neg}), or subtraction (\ieee{sub}) of two intervals, or of one interval and a double.
\end{deffun}


\deffunx{interval}{operator*}{(const interval\& \var{I}, double \var{d})}
\deffunx{interval}{operator*}{(double \var{d}, const interval\& \var{I})}
\begin{deffun}{interval}{operator*}{(const interval\& \var{I1}, const interval\& \var{I2})}
  Multiplication of two intervals, or of one interval and a double: \ieee{mul}
  of IEEE 1788-2015. A zero bound times an infinite bound counts as $0$, as the
  product of a real number of the interval by 0.\newinvfive
\end{deffun}

\deffunx{interval}{operator/}{(const interval\& \var{I}, double \var{d})}
\deffunx{interval}{operator/}{(double \var{d}, const interval\& \var{I})}
\begin{deffun}{interval}{operator/}{(const interval\& \var{I1}, const interval\& \var{I2})}
  Functional division of two intervals, or of one interval and a double:
  \ieee{div} of IEEE 1788-2015. The quotient by an interval containing 0 is the
  hull of the quotients by its nonzero part, and the quotient by \itv{0}{0} is
  the empty set.
\begin{example}
cout << interval(1,2) * interval(-3,4) << ' '
     << interval(0,1) * interval("[1, inf]") << ' '
     << interval(1,2) / interval(0,1) << ' '
     << interval(1,2) / interval(-1,1) << ' '
     << interval(1,2) / interval::zero();
@outputs^[-6, 8] [0, inf] [1, inf] [-inf, inf] [empty]~
\end{example}
\end{deffun}

\deffunx{interval}{operator\%}{(const interval\& \var{I}, double \var{d})}
\deffunx{interval}{operator\%}{(double \var{d}, const interval\& \var{I})}
\begin{deffun}{interval}{operator\%}{(const interval\& \var{I1}, const interval\& \var{I2})}
  Relational division of two intervals, or of one interval and a double
  (\pxref{sec:relational-arithmetic}).
\end{deffun}

\defmethodx{interval}{interval}{inverse}{(void) const}
\begin{deffun}{interval}{inverse}{(const interval\& \var{I})}
  \begin{operation}
    \op{$1/I$}
  \end{operation}
  Inverse of \var{I} (resp. \code{self}): \ieee{recip} of IEEE 1788-2015. The
  inverse of an interval containing 0 is the hull of the inverses of its
  nonzero part.
\begin{example}
cout << inverse(interval(2,4)) << ' ' << inverse(interval(-1,2));
@outputs^[0.25, 0.5] [-inf, inf]~
\end{example}
\end{deffun}

\begin{deffun}{interval}{sqr}{(const interval\& \var{I})}
  \begin{operation}
    \op{$\var{I}^2$}
  \end{operation}
  Square of \var{I}: \ieee{sqr} of IEEE 1788-2015.
\end{deffun}

\begin{deffun}{interval}{sqrt}{(const interval\& \var{I})}
  \begin{operation}
    \op{$\sqrt{I}$}
  \end{operation}
  Functional square root of \var{I}, on the part of \var{I} in \itv{0}{+\infty}:
  \ieee{sqrt} of IEEE 1788-2015. The bounds are the tightest ones, whatever the
  rounding of the square root of the C library.\newinvfive
\begin{example}
cout << sqrt(interval(-4,9)) << ' ' << sqrt(interval(-4,-1));
@outputs^[0, 3] [empty]~
\end{example}
\end{deffun}

\begin{deffun}{interval}{fma}{(const interval\& \var{X}, const interval\& \var{Y}, const interval\& \var{Z})}
  \begin{operation}
    \op{$\hull{\{xy+z\mid x\in X, y\in Y, z\in Z\}}$}
  \end{operation}
  Fused multiply-add: \ieee{fma} of IEEE 1788-2015, each bound rounded once, the
  tightest enclosure.\newinvfive
\begin{example}
cout << fma(interval(1,2), interval(3,4), interval(-1,1));
@outputs^[2, 9]~
\end{example}
\end{deffun}

\subsection{Powers and roots}\label{sec:powers}
%%...........................
\cindex{power}%

The following four functions are those of the namespace \code{gaol}; the
namespace \code{gaol\_ieee1788} has its own \code{pow} (Chapter~\ref{chap:ieee1788}).

\deffunx{interval}{pow}{(const interval\& \var{I}, int \var{n})}
\deffunx{interval}{pow}{(const interval\& \var{I}, unsigned int \var{n})}
\deffunx{interval}{pow}{(const interval\& \var{I}, double \var{p})}
\begin{deffun}{interval}{pow}{(const interval\& \var{I}, const interval\& \var{J})}
  \begin{operation}
    \op{$I^n$}
    \op{$I^p$}
    \op{$I^J$}
  \end{operation}
  Power function. \code{pow} is hybrid in GAOL v5:\newinvfive
  \begin{itemize}
  \item an integer exponent takes the \emph{integer power}, \ieee{pown} of
    IEEE 1788-2015: \var{I} to the power \var{n}, the negative part of \var{I}
    included, $x^0=1$ for every $x$, 0 included, and no value at 0 for $n<0$;
    so do a double \var{p} and a degenerate interval $\var{J}=\itv{n}{n}$ whose
    value is an integer within the \code{int}s, an integer beyond them giving
    \itv{-\infty}{+\infty};
  \item any other exponent takes the \emph{real power}, \ieee{pow} of
    IEEE 1788-2015, on the part of \var{I} in \itv{0}{+\infty}, where $0^y$ is
    $0$ for $y>0$ and has no value for $y\leq0$. An exponent \itv{+\infty}{+\infty}
    or \itv{-\infty}{-\infty}, or an infinite or NaN \var{p}, gives the empty set.
  \end{itemize}
  The result is thus not monotone for the inclusion of \var{J}:
  \code{pow(\itv{-4}{-1}, \itv{2}{2})} is \itv{1}{16}, and
  \code{pow(\itv{-4}{-1}, \itv{1.5}{2.5})} the empty set.
  GAOL~4 truncated a double exponent \var{p}, and computed the powers of a
  negative base on its magnitude.

  The integer powers are computed from exact products, and are the tightest
  enclosures, or one double beyond where the power is within $n\cdot2^{-104}$
  of a double. The real power takes the power of CORE-MATH at the corners of
  the box $I\times J$ where it is the least and the greatest: the upper bound
  is the tightest one, and so is the lower bound but where the power at the
  corner is a double, which it misses by one double.
\begin{example}
cout << pow(interval(-4,-1), 2) << ' '
     << pow(interval(-4,-1), interval(2.0)) << ' '
     << pow(interval(-4,9), interval(0.5)) << ' '
     << pow(interval(-4,-1), interval(1.5, 2.5)) << endl;
@outputs^[1, 16] [1, 16] [0, 3] [empty]~
cout << pow(interval(0.0), interval(0.0)) << ' '
     << pow(interval(-2,2), -2) << ' '
     << pow(interval(4), 0.5) << endl;
@outputs^<1, 1> [0.25, inf] [1.999999999999999, 2]~
\end{example}
\end{deffun}

\deffunx{interval}{nth\_root}{(const interval\& \var{I}, unsigned int \var{n})}
\begin{deffun}{interval}{nth\_root}{(const interval\& \var{I}, int \var{q})}
  \begin{operation}
    \op{$^{n}\sqrt{I}$}
  \end{operation}
  Computes the \var{n}th functional root of \var{I}: \ieee{rootn} of
  IEEE 1788-2015. For an odd \var{n}, the root of a negative number is the
  opposite of the root of its magnitude, and the root is defined on the whole
  real line; for an even \var{n}, only the part of \var{I} in
  \itv{0}{+\infty} is considered.\newinvfive{} A negative \var{q} gives
  $1/x^{1/|q|}$, defined on $\RSet\setminus\{0\}$ for an odd \var{q} and on
  $\left]0,+\infty\right]$ for an even one. Returns the empty set for
  $n=0$.

  The bounds are the tightest ones, or one double beyond: each is proved with
  integer powers. The cube root is that of CORE-MATH, the tightest one, and the
  root of a double that is the \var{n}th power of a double is that double.
\begin{example}
cout << nth_root(interval(-8,27), 3) << ' '
     << nth_root(interval(-4,16), 2) << ' '
     << nth_root(interval(4,16), -2) << ' '
     << nth_root(interval(-8,-1), -3);
@outputs^[-2, 3] [0, 4] [0.25, 0.5] [-1, -0.5]~
\end{example}
\end{deffun}

\subsection{Exponentials and logarithms}
%%......................................

The functions of this section, and of the next two, are bounded with CORE-MATH
at the bounds of the interval, or at the points where the function reaches its
extrema within it. CORE-MATH being correctly rounded, the value at a bound,
computed rounding upward, is the upper bound, and the double below it the
lower bound: the tightest bounds, unless the exact value is a double, which
the functions give exactly where they know it happens ($\exp(0)=1$,
$\log(1)=0$, $\sin(0)=0$, $2^3=8$\ldots).\newinvfive
\cindex{CORE-MATH}%

\begin{deffun}{interval}{exp}{(const interval\& \var{I})}
  Exponential of \var{I}.
\end{deffun}

\deffunx{interval}{exp2}{(const interval\& \var{I})}
\begin{deffun}{interval}{exp10}{(const interval\& \var{I})}
  $2^x$ and $10^x$ for $x$ in \var{I}: \ieee{exp2} and \ieee{exp10} of IEEE
  1788-2015.\newinvfive
\end{deffun}

\deffunx{interval}{expm1}{(const interval\& \var{I})}
\deffunx{interval}{exp2m1}{(const interval\& \var{I})}
\begin{deffun}{interval}{exp10m1}{(const interval\& \var{I})}
  $e^x-1$, $2^x-1$ and $10^x-1$ for $x$ in \var{I}, within \itv{-1}{+\infty}:
  \ieee{expm1}, \ieee{exp2m1} and \ieee{exp10m1} of IEEE 1788-2015
  (Table~10.5).\newinvfive
\end{deffun}

\begin{deffun}{interval}{log}{(const interval\& \var{I})}
  Natural logarithm of \var{I}, defined on $\left]0,+\infty\right]$: an
  interval holding no positive number gives the empty set, where GAOL~4 kept
  its part in \itv{0}{+\infty} and gave \itv{-\infty}{-\text{MAX}} for
  \code{log(\itv{-4}{0})}.\newinvfive
\end{deffun}

\deffunx{interval}{log2}{(const interval\& \var{I})}
\begin{deffun}{interval}{log10}{(const interval\& \var{I})}
  Logarithms in base 2 and 10 of \var{I}: \ieee{log2} and \ieee{log10} of IEEE
  1788-2015.\newinvfive
\end{deffun}

\deffunx{interval}{log1p}{(const interval\& \var{I})}
\deffunx{interval}{log2p1}{(const interval\& \var{I})}
\begin{deffun}{interval}{log10p1}{(const interval\& \var{I})}
  $\log(1+x)$, $\log_2(1+x)$ and $\log_{10}(1+x)$ for $x$ in \var{I}, defined on
  $\left]-1,+\infty\right]$, $-\infty$ being their limit at $-1$: \ieee{logp1},
  \ieee{log2p1} and \ieee{log10p1} of IEEE 1788-2015 (Table~10.5).\newinvfive
\begin{example}
cout << exp(interval(0.0)) << ' ' << log(interval(1)) << ' '
     << log(interval(-4,0.0)) << ' '
     << log(interval(-4,1)) << endl;
@outputs^<1, 1> <0, 0> [empty] [-inf, 0]~
cout << exp2(interval(3)) << ' ' << log2(interval(0.25)) << ' '
     << log10(interval(0.0,1000)) << ' '
     << expm1(interval(-1e300,0.0)) << ' '
     << log1p(interval(-2,0.0)) << endl;
@outputs^<8, 8> <-2, -2> [-inf, 3] [-1, 0] [-inf, 0]~
\end{example}
\end{deffun}

\subsection{Trigonometric functions}
%%..................................

\begin{deffun}{interval}{cos}{(const interval\& \var{I})}
  Returns the cosine of \var{I}, within \itv{-1}{1}.
\end{deffun}

\begin{deffun}{interval}{sin}{(const interval\& \var{I})}
  Returns the sine of \var{I}, within \itv{-1}{1}.
\end{deffun}

\begin{deffun}{interval}{tan}{(const interval\& \var{I})}
  Returns the tangent of \var{I}: \itv{-\infty}{+\infty} when \var{I} holds a
  pole.

  The three functions find the pieces of \var{I} on which they are monotonic
  by dividing its bounds by an enclosure of $\pi$ and, where the quotients
  cannot tell, next to an extremum or a pole and at the large magnitudes, from
  the signs of the derivative at the bounds, which CORE-MATH gives exactly:
  their bounds are within one double of the tightest ones at every
  magnitude.\newinvfive{} GAOL~4 gave \itv{-1}{1} for \code{cos(\itv{2^{60}}{2^{60}})},
  and computed the sine as a shifted cosine, which widened narrow intervals.
\begin{example}
cout << cos(interval(1,4)) << endl
     << sin(interval(1,2)) << endl
     << tan(interval(1,2)) << endl
     @rem^// 1152921504606846976 is 2 to the 60~
     << cos(interval(1152921504606846976.0)) << endl;
@outputs^[-1, 0.5403023058681398]~
@outputs^[0.8414709848078965, 1]~
@outputs^[-inf, inf]~
@outputs^[-0.5567960822766418, -0.5567960822766416]~
\end{example}
\end{deffun}

\begin{deffun}{interval}{acos}{(const interval\& \var{I})}
  Returns the arccosine of \var{I}, on the part of \var{I} in \itv{-1}{1}.
\end{deffun}

\begin{deffun}{interval}{asin}{(const interval\& \var{I})}
  Returns the arcsine of \var{I}, on the part of \var{I} in \itv{-1}{1}.
\end{deffun}

\begin{deffun}{interval}{atan}{(const interval\& \var{I})}
  Returns the arctangent of \var{I}.
\begin{example}
cout << acos(interval(1,3)) << ' ' << asin(interval(1)) << ' '
     << atan(interval(1));
@outputs^<0, 0> [1.570796326794896, 1.570796326794897]~
@outputs^[0.7853981633974482, 0.7853981633974484]~
\end{example}
\end{deffun}

\begin{deffun}{interval}{atan2}{(const interval\& \var{Y}, const interval\& \var{X})}
  Returns the hull of the angles in $(-\pi,\pi]$ of the points of ordinate in \var{Y} and abscissa
  in \var{X}, the point $(0,0)$ left out, as the \ieee{atan2} of IEEE~1788-2015:
  \code{atan2([0],[0])} is empty, and the result is $[-\pi,\pi]$ whenever the box holds points of the
  half-line $y=0$, $x<0$, where the angle is $\pi$, and points below it, where it is next to $-\pi$.
  GAOL~4 declared the function, but it raised \code{unavailable\_feature\_error}.\newinvfive
\begin{example}
cout << atan2(interval(0.0), interval(0.0)) << endl
     << atan2(interval(-1,1), interval(-1)) << endl
     << atan2(interval("[1, inf]"), interval("[1, inf]")) << endl;
@outputs^[empty]~
@outputs^[-3.141592653589794, 3.141592653589794]~
@outputs^[0, 1.570796326794897]~
\end{example}
\end{deffun}

\deffunx{interval}{sinpi}{(const interval\& \var{I})}
\deffunx{interval}{cospi}{(const interval\& \var{I})}
\begin{deffun}{interval}{tanpi}{(const interval\& \var{I})}
  $\sin(\pi x)$, $\cos(\pi x)$ and $\tan(\pi x)$ for $x$ in \var{I}:
  \ieee{sinPi}, \ieee{cosPi} and \ieee{tanPi} of IEEE 1788-2015
  (Table~10.5).\newinvfive{} Their extrema and poles are at the multiples of
  $1/2$, which are doubles, and are found exactly, at every magnitude.
  \code{tanpi} has no value at the half-integers, which are left out: a
  half-integer alone gives the empty set.
\end{deffun}

\deffunx{interval}{asinpi}{(const interval\& \var{I})}
\deffunx{interval}{acospi}{(const interval\& \var{I})}
\deffunx{interval}{atanpi}{(const interval\& \var{I})}
\begin{deffun}{interval}{atan2pi}{(const interval\& \var{Y}, const interval\& \var{X})}
  $\operatorname{asin}(x)/\pi$, $\operatorname{acos}(x)/\pi$,
  $\operatorname{atan}(x)/\pi$ and $\operatorname{atan2}(y,x)/\pi$: \ieee{asinPi},
  \ieee{acosPi}, \ieee{atanPi} and \ieee{atan2Pi} of IEEE 1788-2015
  (Table~10.5), within \itv{-1/2}{1/2}, \itv{0}{1}, \itv{-1/2}{1/2} and
  \itv{-1}{1}.\newinvfive
\begin{example}
cout << sinpi(interval(1e17)) << ' '
     << sin(interval::pi() * 1e17) << endl;
@outputs^<0, 0> [-1, 1]~
cout << tanpi(interval(0.5)) << ' '
     << tanpi(interval(0.25,0.75)) << ' '
     << asinpi(interval(1)) << ' '
     << atan2pi(interval(1), interval(-1)) << endl;
@outputs^[empty] [-inf, inf] <0.5, 0.5> <0.75, 0.75>~
\end{example}
\end{deffun}

\subsection{Hyperbolic functions}
%%...............................

\begin{deffun}{interval}{cosh}{(const interval\& \var{I})}
  Returns the hyperbolic cosine of \var{I}.
\end{deffun}

\begin{deffun}{interval}{acosh}{(const interval\& \var{I})}
  Returns the hyperbolic arccosine of \var{I}, on the part of \var{I} in \itv{1}{+\infty}.
\end{deffun}

\begin{deffun}{interval}{sinh}{(const interval\& \var{I})}
  Returns the hyperbolic sine of \var{I}.
\end{deffun}

\begin{deffun}{interval}{asinh}{(const interval\& \var{I})}
  Returns the hyperbolic arcsine of \var{I}.
\end{deffun}

\begin{deffun}{interval}{tanh}{(const interval\& \var{I})}
  Returns the hyperbolic tangent of \var{I}, within \itv{-1}{1}.
\end{deffun}

\begin{deffun}{interval}{atanh}{(const interval\& \var{I})}
  Returns the hyperbolic arctangent of \var{I}, on the part of \var{I} in \itv{-1}{1}.

  The six functions are those of CORE-MATH, where GAOL~4 took them from the
  mathematical library of the system, which is sometimes further than one
  double from the exact value.\newinvfive
\begin{example}
cout << cosh(interval(-1,2)) << ' '
     << tanh(interval(20,30)) << endl
     << acosh(interval(0,1)) << ' '
     << atanh(interval(-2,2)) << endl;
@outputs^[1, 3.762195691083632] [0.9999999999999998, 1]~
@outputs^<0, 0> [-inf, inf]~
\end{example}
\end{deffun}

\subsection{Other functions}
%%..........................

\begin{deffun}{interval}{hypot}{(const interval\& \var{X}, const interval\& \var{Y})}
  $\sqrt{x^2+y^2}$ for $x$ in \var{X} and $y$ in \var{Y}: \ieee{hypot} of
  IEEE 1788-2015 (Table~10.5), computed at the points of the box nearest to the
  origin and farthest from it.\newinvfive
\end{deffun}

\begin{deffun}{interval}{rsqrt}{(const interval\& \var{I})}
  $1/\sqrt{x}$ for $x$ in \var{I}, defined on $\left]0,+\infty\right]$,
  $+\infty$ being its limit at 0: \ieee{rSqrt} of IEEE 1788-2015
  (Table~10.5).\newinvfive
\begin{example}
cout << hypot(interval(3), interval(4)) << ' '
     << hypot(interval(-3,1), interval(4,5)) << endl
     << rsqrt(interval(4,16)) << ' '
     << rsqrt(interval(0.0,4)) << endl;
@outputs^<5, 5> [4, 5.830951894845301]~
@outputs^[0.25, 0.5] [0.5, inf]~
\end{example}
\end{deffun}

\noindent
Among the functions IEEE 1788-2015 recommends (Table~10.5), \ieee{compoundm1}
alone is not provided, CORE-MATH having no version of it for doubles.

\section{Relational Arithmetic}\label{sec:relational-arithmetic}
%%-----------------------------
\cindex{reverse functions|see{relational arithmetic}}%

The relational functions of GAOL are the \dfn{reverse functions} of
IEEE 1788-2015 (10.5.4, Table 10.1), with their arguments in another order:
each computes the hull of the $x$ of an interval \var{I} whose image by the
function is in \var{J}. Their correspondence is given at the end of this
section, and the namespace \code{gaol\_ieee1788} provides them under their
names (Chapter~\ref{chap:ieee1788}).

\defmethodx{interval}{interval}{operator\% }{(double \var{d}) const}
\defmethodx{interval}{interval}{operator\% }{(const interval\& \var{I}) const}
\begin{deffun}{interval}{operator\%}{(double \var{d}, const interval\& \var{I})}
  \begin{operation}
    \op{$\code{self}\,\%\,\var{I}=\hull{\{z\in \RSet\mid\exists x\in \code{self}, \exists y\in I\colon x=yz\}}$}
  \end{operation}
Relational division: \code{div\_rel(self, \var{I}, \itv{-\infty}{+\infty})}.
\begin{example}
cout << interval(1,2) % interval(0,1) << ' '
     << interval(1,2) % interval(-1,1);
@outputs^[1, inf] [-inf, inf]~
\end{example}
\end{deffun}

\subsection{$(n+1)$-ary relational functions}
%%...........................................

Consider the relation $y = \cos x$ where $x$ and $y$ are
interval variables. One would like to be able to express this relation
in the equivalent way: $x = \mathop{acos} y$. However, one cannot use the
$\mathop{acos}$ function because its result is always included into the
interval \itv{0}{\pi}. What we need here is a relational
version of the $\mathop{acos}$ function. But, since for any value $x$
there are infinitely many values $y$ verifying $x=\mathop{acos} y$, we
have to take into account the domain of $y$. As a consequence, we
define a new binary operator \code{acos\_rel} whose definition is as
follows:

\begin{equation*}
{\tt\mathop{acos\_rel}}(Y,X)=\hull{\{x\in X\mid\exists y\in Y\colon y=\cos x\}}
\end{equation*}

\noindent
This is to be contrasted with the previous definition of the
$\mathop{acos}$ function:

\begin{equation*}
\code{acos}(Y)=\hull{\{x\in\RSet\mid\exists y\in Y\colon x=\mathop{acos} y\}}
\end{equation*}

\noindent
Figure~\ref{fig:relational-cos} presents the different results obtained when computing
either $\code{acos}(J)$ or $\code{acos\_rel}(J,I)$. An interval \var{I} that
meets no piece of the preimage gives the empty set.\newinvfive

\begin{wrapfigure}{O}{\marginparwidth}
 \includegraphics[width=.9\marginparwidth,bb=11 262 553 470]{relation-cos}
  \caption{Relational cosine}
  \label{fig:relational-cos}
\end{wrapfigure}

\begin{deffun}{interval}{acos\_rel}{(const interval\& J, const interval\& I)}
  \begin{operation}
    \op{$\code{acos\_rel}(J,I)=\hull{\{x\in I\mid\exists y\in J\colon y=\cos x\}}$}
  \end{operation}
  Returns the \emph{relational arccosine} of \var{J} w.r.t. \var{I}.
\end{deffun}

\begin{deffun}{interval}{asin\_rel}{(const interval\& J, const interval\& I)}
  \begin{operation}
    \op{$\code{asin\_rel}(J,I)=\hull{\{x\in I\mid\exists y\in J\colon y=\sin x\}}$}
  \end{operation}
  Returns the \emph{relational arcsine} of \var{J} w.r.t. \var{I}.
\end{deffun}

\begin{deffun}{interval}{atan\_rel}{(const interval\& J, const interval\& I)}
  \begin{operation}
    \op{$\code{atan\_rel}(J,I)=\hull{\{x\in I\mid\exists y\in J\colon y=\tan x\}}$}
  \end{operation}
  Returns the \emph{relational arctangent} of \var{J} w.r.t. \var{I}.

  The three functions enclose the multiples of $\pi$ of the pieces of the
  preimage in double-double arithmetic, and add the inverse function of \var{J}
  before rounding the sum: their bounds are within one double of the accurate
  ones.\newinvfive
\begin{example}
cout << asin_rel(interval(0.0), interval(1,7)) << endl
     << atan_rel(interval(1), interval(0,5)) << endl;
@outputs^[3.141592653589793, 6.283185307179588]~
@outputs^[0.7853981633974482, 3.926990816987242]~
\end{example}
\end{deffun}

\begin{deffun}{interval}{acosh\_rel}{(const interval\& J, const interval\& I)}
  \begin{operation}
    \op{$\code{acosh\_rel}(J,I)=\hull{\{x\in I\mid\exists y\in J\colon y=\cosh x\}}$}
  \end{operation}
  Returns the \emph{relational hyperbolic arccosine} of \var{J} w.r.t. \var{I}.
\end{deffun}

\deffunx{interval}{asinh\_rel}{(const interval\& J, const interval\& I)}
\begin{deffun}{interval}{atanh\_rel}{(const interval\& J, const interval\& I)}
  \begin{operation}
    \op{$\code{asinh\_rel}(J,I)=\hull{\{x\in I\mid\exists y\in J\colon y=\sinh x\}}$}
    \op{$\code{atanh\_rel}(J,I)=\hull{\{x\in I\mid\exists y\in J\colon y=\tanh x\}}$}
  \end{operation}
  Return the \emph{relational hyperbolic arcsine} and \emph{arctangent} of
  \var{J} w.r.t. \var{I}, that is, $\operatorname{asinh}(J)\cap I$ and
  $\operatorname{atanh}(J)\cap I$, $\sinh$ and $\tanh$ being one-to-one.
\begin{example}
cout << acosh_rel(interval(1,2), interval::universe()) << endl
     << asinh_rel(interval(0,1), interval::universe()) << endl;
@outputs^[-1.316957896924817, 1.316957896924817]~
@outputs^[0, 0.8813735870195431]~
\end{example}
\end{deffun}

\begin{deffun}{interval}{sqrt\_rel}{(const interval\& J, const interval\& I)}
  \begin{operation}
    \op{$\code{sqrt\_rel}(J,I)=\hull{\{x\in I\mid\exists y\in J\colon y=x^2\}}$}
  \end{operation}
  Returns the \emph{relational square root} of \var{J} w.r.t. \var{I}.
\begin{example}
cout << sqrt_rel(interval(1,4), interval::universe()) << ' '
     << sqrt_rel(interval(1,4), interval(0,10));
@outputs^[-2, 2] [1, 2]~
\end{example}
\end{deffun}

\begin{deffun}{interval}{nth\_root\_rel}{(const interval\& \var{J}, unsigned int \var{n}, \\\hbox{}\hfill const interval\& \var{I})}
  \begin{operation}
    \op{$\code{nth\_root\_rel}(J,n,I)=\hull{\{x\in I\mid\exists y\in J\colon y=x^n\}}$}
  \end{operation}
  Returns the \emph{relational inverse \var{n}-th power} of
  \var{J} w.r.t. \var{I}, for $n\geq1$. It takes the roots of
  \code{nth\_root()}.
\begin{example}
interval J(1,16);
cout << nth_root_rel(J, 4, interval::universe()) << ' '
     << nth_root_rel(J, 4, interval(-10,0.0));
@outputs^[-2, 2] [-2, -1]~
\end{example}
\end{deffun}

\begin{deffun}{interval}{invabs\_rel}{(const interval\& J, const interval\& I)}
  \begin{operation}
    \op{$\code{invabs\_rel}(J,I)=\hull{\{x\in I\mid |x|\in J\}}$}
  \end{operation}
  Returns the \emph{relational inverse absolute value} of \var{J} w.r.t. \var{I}.
  The negative part of \var{J} holds no magnitude: GAOL~4 took it as
  magnitudes, and gave \itv{-2}{2} where the result is empty.\newinvfive
\begin{example}
cout << invabs_rel(interval(1,2), interval::universe()) << ' '
     << invabs_rel(interval(-2,-1), interval(-10,10));
@outputs^[-2, 2] [empty]~
\end{example}
\end{deffun}

\begin{deffun}{interval}{div\_rel}{(const interval\& K, const interval\& J, const interval\& I)}
  \begin{operation}
    \op{$\code{div\_rel}(K,J,I)=\hull{\{x\in I\mid\exists z\in K,\exists y\in J\colon z=xy\}}$}
  \end{operation}
  Returns the \emph{ternary relational division} of \var{K} by \var{J}
  w.r.t. \var{I}, which is empty when no $x$ of \var{I} is in the
  relation.\newinvfive
\begin{example}
cout << div_rel(interval(1,2), interval(-1,1),
                interval(0,10)) << ' '
     << div_rel(interval(5,6), interval(1,2),
                interval(10,20));
@outputs^[1, 10] [empty]~
\end{example}
\end{deffun}

\noindent
The correspondence with the reverse functions of IEEE 1788-2015, whose
last argument $x$ is \itv{-\infty}{+\infty} when it is left out:

\begin{center}\renewcommand{\arraystretch}{1.1}
  \begin{tabular}{ll}
    \hline
    \multicolumn{1}{c}{IEEE 1788-2015} & \multicolumn{1}{c}{GAOL}\\
    \hline\hline
    \ieee{sqrRev}$(c,x)$ & \code{sqrt\_rel(c, x)}\\
    \ieee{absRev}$(c,x)$ & \code{invabs\_rel(c, x)}\\
    \ieee{pownRev}$(c,x,p)$, $p\geq1$ & \code{nth\_root\_rel(c, p, x)}\\
    \ieee{sinRev}, \ieee{cosRev}, \ieee{tanRev}$(c,x)$ & \code{asin\_rel}, \code{acos\_rel}, \code{atan\_rel(c, x)}\\
    \ieee{coshRev}$(c,x)$ & \code{acosh\_rel(c, x)}\\
    \ieee{mulRev}$(b,c,x)$ & \code{div\_rel(c, b, x)}\\
    \hline
  \end{tabular}
\end{center}

\noindent
\ieee{powRev1}, \ieee{powRev2}, \ieee{atan2Rev1}, \ieee{atan2Rev2}, \ieee{pownRev}
for $p\leq0$ and \ieee{mulRevToPair} are not provided.

\section{Cancellative addition and subtraction}
%%---------------------------------------------
\cindex{cancellative operations}%

\begin{deffun}{interval}{cancel\_minus}{(const interval\& \var{X}, const interval\& \var{Y})}
  The tightest interval $Z$ such that $Y+Z$ contains $X$,
  \itv{\leftBound{X}-\leftBound{Y}}{\rightBound{X}-\rightBound{Y}}, which exists
  when \var{X} and \var{Y} are bounded and \var{X} is at least as wide as \var{Y};
  \itv{-\infty}{+\infty} otherwise, and the empty set for an empty \var{X} and a
  bounded \var{Y}: \ieee{cancelMinus} of IEEE 1788-2015 (10.5.6). Whether
  \var{X} is at least as wide as \var{Y} is decided exactly, which rounding the
  differences of the bounds cannot do.\newinvfive
\end{deffun}

\begin{deffun}{interval}{cancel\_plus}{(const interval\& \var{X}, const interval\& \var{Y})}
  \code{cancel\_minus(\var{X}, -\var{Y})}: \ieee{cancelPlus} of IEEE
  1788-2015.\newinvfive
\begin{example}
cout << cancel_minus(interval(1,5), interval(0,2)) << ' '
     << cancel_minus(interval(1,2), interval(0,3)) << ' '
     << cancel_plus(interval(1,5), interval(0,2));
@outputs^[1, 3] [-inf, inf] [3, 5]~
\end{example}
\end{deffun}


\chapter{Interval functions}\label{chap:functions}
%%==========================

The numeric functions of IEEE 1788-2015 (Table 10.2, 12.12.8) are those of this
chapter; the namespace \code{gaol\_ieee1788} names them as the standard does,
and returns $+\infty$ and $-\infty$ for the bounds of the empty set, where
\code{left()} and \code{right()} return NaN (Chapter~\ref{chap:ieee1788}).

\begin{defmethod}{interval}{double}{rad}{(void) const}
  \begin{operation}
    \op{$\code{rad}(\itv{a}{b})=\min\{r\mid [m-r,m+r]\supseteq\itv{a}{b}\}$}
  \end{operation}
Returns the radius of \code{self}, the smallest double $r$ such that
$[m-r,m+r]$ contains \code{self}, $m$ being \code{self.midpoint()}: $+\infty$
for an unbounded interval, NaN for the empty set (\ieee{rad} of IEEE
1788-2015).\newinvfive
\end{defmethod}

\begin{defmethod}{interval}{void}{mid\_rad}{(double\& \var{m}, double\& \var{r}) const}
Sets \var{m} to \code{self.midpoint()} and \var{r} to \code{self.rad()}
(\ieee{midRad} of IEEE 1788-2015).\newinvfive
\begin{example}
double m, r;
interval(1,4).mid_rad(m, r);
cout << m << ' ' << r << ' ' << interval("[1, inf]").rad();
@outputs^2.5 1.5 inf~
\end{example}
\end{defmethod}

\begin{defmethod}{interval}{double}{width}{(void) const}
  \begin{operation}
    \op{$\code{width}(\itv{a}{b})=\roundUp{b-a}$}
  \end{operation}
Returns the width of \code{self}: \ieee{wid} of IEEE 1788-2015. Returns NaN
whenever the interval is empty, as the standard has it, where GAOL~4 returned
$-1$.\newinvfive
\begin{example}
cout << interval(4,6).width() << ' '
     << (interval(1,next_float(1)).width()
         == std::numeric_limits<double>::epsilon()) << ' '
     << interval::emptyset().width();
@outputs^2 true nan~
\end{example}
\end{defmethod}

\marginnocite{Hansen:book92} % To handle the minipage in environment 'defmethod'
\begin{defmethod}{interval}{double}{mig}{(void) const}
  Returns the \dfn{mignitude} of \code{self}: \ieee{mig} of IEEE 1788-2015. See the book by
  Hansen~\margincite*[chap.~3]{Hansen:book92}. The mignitude of an
  interval \itv{a}{b} is the smallest absolute value of the numbers in
  the interval, that is: $0$ if the interval straddles \code{0}, $a$
  if the interval is strictly positive, and $-b$ otherwise.

\paragraph{Note.} The mignitude of the empty interval is a NaN.
\begin{example}
cout << interval(4,5).mig() << ' '
     << interval(-6,-3).mig() << ' '
     << interval(-3,8).mig();
@outputs^4 3 0~
\end{example}
\end{defmethod}

\marginnocite{Stahl:PhD}
\begin{defmethod}{interval}{double}{smig}{(void) const}
  Returns the \dfn{signed mignitude} of \code{self}.  See Stahl's
  thesis~\margincite*[def.~1.3.28]{Stahl:PhD}. The signed mignitude of
  an interval \itv{a}{b} is $0$ if the interval straddles \code{0},
  $a$ if the interval is strictly positive, and $b$ otherwise.

\paragraph{Note.} The signed mignitude of the empty interval is a NaN.
\begin{example}
cout << interval(4,5).smig() << ' '
     << interval(-6,-3).smig() << ' '
     << interval(-3,8).smig();
@outputs^4 -3 0~
\end{example}
\end{defmethod}

\begin{defmethod}{interval}{double}{mag}{(void) const}
  Returns the \dfn{magnitude} of \code{self}: \ieee{mag} of IEEE 1788-2015. The
  magnitude of an interval \itv{a}{b} is the greatest absolute value of the
  numbers in the interval.

\paragraph{Note.} The magnitude of the empty interval is a NaN.
\begin{example}
cout << interval(4,5).mag() << ' '
     << interval(-6,-3).mag() << ' '
     << interval(-10,5).mag();
@outputs^5 6 10~
\end{example}
\end{defmethod}

\begin{deffun}{double}{hausdorff}{(const interval\& \var{I1}, const interval\& \var{I2})}
  Returns the Hausdorff distance between the two sets defined by intervals \var{I1} and
  \var{I2}, that is:
\begin{equation*}
\code{hausdorff}(I_1,I_2)=\max(|\leftBound{I_1}-\leftBound{I_2}|,
                          |\rightBound{I_1}-\rightBound{I_2}|)
\end{equation*}
rounded upward: the tightest upper bound of the distance.\newinvfive
\begin{example}
cout << hausdorff(interval(4,8),interval(5,10));
@outputs^2~
\end{example}
\end{deffun}

\marginnocite{Goualard:TOMS2014}
\begin{defmethod}{interval}{double}{midpoint}{(void) const}
  Returns the midpoint of \code{self}, \ieee{mid} of IEEE 1788-2015. Given $a$
  and $b$ two finite floating-point numbers, we have the following cases, MAX
  being the largest positive double, \code{std::numeric\_limits<double>::max()}:
\begin{equation*}
\left\{\begin{array}{ll}
\code{midpoint}(\emptyset{})&=\text{NaN}\\
\code{midpoint}(\itv{-\infty}{+\infty})&=0\\
\code{midpoint}(\itv{-\infty}{b})&=-\text{MAX}\\
\code{midpoint}(\itv{a}{+\infty})&=\text{MAX}\\
\code{midpoint}(\itv{a}{b})&=\roundNearest{(a+b)/2} \\
\end{array}\right.
\end{equation*}
The midpoint of \itv{a}{b} is rounded to nearest, ties to even: it is the half
of $a+b$, or the sum of the halves of $a$ and $b$ when $a+b$
overflows~\margincite*{Goualard:TOMS2014}.
\end{defmethod}

\begin{defmethod}{interval}{interval}{mid}{(void) const}
  Returns the tightest interval enclosing the midpoint of \code{self}, which is
  included in \code{self}.
  With the same notations as for \code{midpoint()}, we have the cases:
\begin{equation*}
\left\{\begin{array}{ll}
\code{mid}(\emptyset{})&=\emptyset{}\\
\code{mid}(\itv{-\infty}{+\infty})&=\itv{0}{0}\\
\code{mid}(\itv{-\infty}{b})&=\itv{-\text{MAX}}{-\text{MAX}}\\
\code{mid}(\itv{a}{+\infty})&=\itv{\text{MAX}}{\text{MAX}}\\
\code{mid}(\itv{a}{b})&=\itv{\roundDn{(a+b)/2}}{\roundUp{(a+b)/2}}
\end{array}\right.
\end{equation*}
The midpoint is computed as \code{midpoint()} computes it, with the rounding
outward. GAOL~4 did not enclose the midpoint of some intervals whose bounds are
subnormal numbers.\newinvfive
\end{defmethod}

\begin{defmethod}{interval}{double}{left}{(void) const}
  \begin{operation}
    \op{$\code{left}(\itv{x}{y})=x$}
  \end{operation}
  Returns the left bound of \code{self}: NaN for the empty set, whose bounds are
  NaN. \code{gaol\_ieee1788::inf()} returns $+\infty$ there, as IEEE 1788-2015
  has it.
\end{defmethod}


\begin{defmethod}{interval}{double}{right}{(void) const}
  \begin{operation}
    \op{$\code{right}(\itv{x}{y})=y$}
  \end{operation}
  Returns the right bound of \code{self}: NaN for the empty set.
  \code{gaol\_ieee1788::sup()} returns $-\infty$ there.
\end{defmethod}

\begin{deffun}{interval}{abs}{(const interval\& \var{I})}
  \begin{operation}
    \op{$\text{abs}(I)=\{|x|\in\RSet^+\mid x\in I\}$}
  \end{operation}
  Returns the absolute value of \var{I}.
  \begin{example}
cout << abs(interval(-5,6)) << ' '
     << abs(interval(-4,-2));
@outputs^[0, 6] [2, 4]~
  \end{example}
\end{deffun}

\marginnocite{Ratschek-Rokne:1995}
\begin{deffun}{double}{chi}{(const interval \&\var{I})}
This function, introduced by Ratscheck and Rokne~\margincite*{Ratschek-Rokne:1995},
characterizes the degree of symmetry of intervals. Its definition is as follows:

\begin{equation*}
\text{For} I=\itv{a}{b}, \quad\text{chi}(I)=\left\{\begin{array}{ll}
    -1  & \text{if } I = 0\\
    a/b & \text{if } |a|\leq |b|\\
    b/a & \text{otherwise}
  \end{array}\right.
\end{equation*}
\begin{example}
cout << chi(interval(3,6)) << ' '
     << chi(interval(-6,3)) << ' '
     << chi(interval::emptyset()) << ' '
     << chi(interval::universe()) << ' '
     << chi(interval("[-5,inf]"));
@outputs^0.5 -0.5 nan 1 0~
\end{example}
\end{deffun}

\begin{deffun}{interval}{min}{(const interval \&I, const interval \&J)}
  \begin{operation}
    \op{$\min(\itv{a}{b}, \itv{c}{d}) = \itv{\min(a,c)}{\min(b,d)}$}
  \end{operation}
  Returns the minimum of two intervals.
  \begin{example}
cout << min(interval(5,6),interval(3,9)) << ' '
     << min(interval::emptyset(),interval(3,8));
@outputs^[3, 6] [empty]~
  \end{example}
\end{deffun}

\begin{deffun}{interval}{max}{(const interval \&\var{I}, const interval \&\var{J})}
  \begin{operation}
    \op{$\max(\itv{a}{b}, \itv{c}{d}) = \itv{\max(a,c)}{\max(b,d)}$}
  \end{operation}
  Returns the maximum of two intervals.
\begin{example}
cout << max(interval(5,6),interval(3,9)) << ' '
     << max(interval::emptyset(),interval(3,8));
@outputs^[5, 9] [empty]~
\end{example}
\end{deffun}

\section{Integer functions}
%%-------------------------

The functions of this section are those of IEEE 1788-2015 (Table 9.1): each is
non-decreasing, so that the bounds of the result are its values at the bounds
of the interval, and each is exact. None of them depends on the rounding
direction.

\begin{deffun}{interval}{floor}{(const interval \&\var{I})}
  \begin{operation}
    \op{$\itv{\code{floor}(I.\code{left}())}{\code{floor}(I.\code{right}())}$}
  \end{operation}

  \begin{example}
cout << floor(interval(4.5,6.5)) << ' '
     << floor(interval("[-10.4,3.5]"));
@outputs^[4, 6] [-11, 3]~
  \end{example}
\end{deffun}

\begin{deffun}{interval}{ceil}{(const interval \&\var{I})}
  \begin{operation}
    \op{\itv{\code{ceil}(I.\code{left}())}{\code{ceil}(I.\code{right}())}}
  \end{operation}
  \begin{example}
cout << ceil(interval(4.5,6.5)) << ' '
     << ceil(interval("[-10.4,3.5]"));
@outputs^[5, 7] [-10, 4]~
  \end{example}
\end{deffun}

\begin{deffun}{interval}{integer}{(const interval \&\var{I})}
  \begin{operation}
    \op{\itv{\code{ceil}(I.\code{left}())}{\code{floor}(I.\code{right}())}}
  \end{operation}
  Narrows down the bounds to the closest integers: the empty set when \var{I}
  holds no integer. Note that the
  resulting bounds are still \code{double} numbers, and may therefore
  not be representable with integral types.
  \begin{example}
cout << integer(interval(4.5,6.5)) << ' '
     << integer(interval(4.2,4.8));
@outputs^[5, 6] [empty]~
  \end{example}
\end{deffun}

\begin{deffun}{interval}{sign}{(const interval \&\var{I})}
  The signs of the elements of \var{I}: $-1$ below 0, 0 at 0, 1 above
  (\ieee{sign}).\newinvfive
\end{deffun}

\begin{deffun}{interval}{trunc}{(const interval \&\var{I})}
  The elements of \var{I} rounded toward zero (\ieee{trunc}).\newinvfive
\end{deffun}

\deffunx{interval}{round\_ties\_to\_even}{(const interval \&\var{I})}
\begin{deffun}{interval}{round\_ties\_to\_away}{(const interval \&\var{I})}
  The elements of \var{I} rounded to the nearest integer, halfway values to the
  even one, respectively away from zero (\ieee{roundTiesToEven},
  \ieee{roundTiesToAway}).\newinvfive
  \begin{example}
cout << sign(interval(-2,3)) << ' '
     << trunc(interval(-2.5,3.7)) << ' '
     << round_ties_to_even(interval(0.5,2.5)) << ' '
     << round_ties_to_away(interval(0.5,2.5));
@outputs^[-1, 1] [-2, 3] [0, 2] [1, 3]~
  \end{example}
\end{deffun}

\section{Splitting methods}
%%-------------------------

\begin{defmethod}{interval}{void}{split}{(interval\& \var{I1}, interval\& \var{I2}) const}
  Splits \code{self} into two parts using \code{midpoint()}; returns the left part
  in \var{I1} and the right part in \var{I2}.

  \var{I1} or \var{I2} may be equal to \code{self}.
  \begin{example}
interval I1a, I2a,
         I3(1.0,next_float(1.0)),
         I1b, I2b;
interval(4,5).split(I1a,I2a);
I3.split(I1b,I2b);
cout << I1a << ' ' << I2a << ' '
     << I1b.set_eq(interval(1.0)) << ' ' << I2b.set_eq(I3) << endl;
@outputs^[4, 4.5] [4.5, 5] true true~
  \end{example}
\end{defmethod}

\begin{defmethod}{interval}{interval}{split\_left}{(void) const}
  \begin{operation}
    \op{$\code{split\_left}(\itv{a}{b})=\itv{a}{m}$}
  \end{operation}
  Splits \code{self} into two parts using $m=\code{midpoint()}$ and returns the left
  part.
  \begin{example}
cout << interval(4,5).split_left();
@outputs^[4, 4.5]~
  \end{example}
\end{defmethod}

\begin{defmethod}{interval}{interval}{split\_right}{(void) const}
  \begin{operation}
    \op{$\code{split\_right}(\itv{a}{b})=\itv{m}{b}$}
  \end{operation}
  Splits \code{self} into two parts using $m=\code{midpoint()}$ and
  returns the right part.
  \begin{example}
cout << interval(4,5).split_right();
@outputs^[4.5, 5]~
  \end{example}
\end{defmethod}

\section{Union and intersection}
%%------------------------------

\begin{deffun}{interval}{operator\&}{(const interval\& I1, const interval\& I2)}
  \begin{operation}
    \op{$\var{I1}\cap\var{I2}$}
  \end{operation}
  Returns the interval resulting from the intersection of \var{I1} and
  \var{I2}: \ieee{intersection} of IEEE 1788-2015. The intersection of disjoint
  intervals is the empty set \code{interval::emptyset()}, whose bounds are NaN.
  GAOL~4 kept the greatest left bound and the smallest right bound,
  \itv{3}{2} for $\itv{1}{2}\cap\itv{3}{4}$, which the operations computing on
  the bounds did not take for empty.\newinvfive
  \begin{example}
cout << (interval(4,6) & interval(5,9)) << ' '
     << (interval(1,2) & interval(3,4));
@outputs^[5, 6] [empty]~
  \end{example}
\end{deffun}

\begin{deffun}{interval}{operator|}{(const interval\& I1, const interval\& I2)}
  \begin{operation}
    \op{$\hull{(\var{I1}\cup\var{I2})}$}
  \end{operation}
  Returns the hull of the union of \var{I1} and \var{I2}: \ieee{convexHull} of
  IEEE 1788-2015.
  \begin{example}
cout << (interval(3,6) | interval(9,12));
@outputs^[3, 12]~
  \end{example}
\end{deffun}

\chapter{Input/output}\label{chap:io}
%%====================

\section{Reading intervals}\label{sec:reading-intervals}
%%-------------------------

\begin{deffun}{istream\&}{operator>{}>}{(istream\& \var{in}, interval\& I)}
  Reads a line from the input stream \var{in}, and assigns to \var{I} the
  interval it writes. If the line is syntactically ill-formed, \var{I} is set
  to the empty set, and an \code{input\_format\_error} exception is thrown
  (\pxref{sec:exceptions}) if the library was compiled with exceptions
  enabled (\pxref{sec:configuration}); alternatively, it prints an
  error message to \code{cerr} and aborts if exceptions were disabled.
\begin{example}
std::istringstream in("[1, 2]\n[0.1, 0.2]\n");
interval a, b;
in >> a >> b;
cout << a << ' ' << b;
@outputs^[1, 2] [0.09999999999999999, 0.2000000000000001]~
\end{example}
\end{deffun}

\subsection{Input format}\label{sec:input-format}
%%.......................

A string to be translated into an interval must have the following
syntax (with terminals in lower case and non-terminals in slanted upper case):

\begin{Verbatim}[commandchars=\@\^\~,fontsize=\small]
@NT^expr~
	: @NT^number~
	| @NT^uncertain_number~ @rem^@hfill// An interval literal, never a bound~
	| @NT^literal~
	| @TM^dmin~  @rem^@hfill// Smallest positive normal floating-point number~
	| @TM^dmax~ @rem^@hfill// Largest positive floating-point number~
	| @TM^pi~
	| @TM^inf~ @rem^@hfill// Floating-point positive ``infinity''~
	| @TM^infinity~ @rem^@hfill// The same~
	| @TM^empty~ @rem^@hfill// Empty interval~
	| @NT^expr~ @TM^+~ @NT^expr~
	| @NT^expr~ @TM^-~ @NT^expr~
	| @NT^expr~ @TM^*~ @NT^expr~
	| @NT^expr~ @TM^/~ @NT^expr~
	| @TM^-~ @NT^expr~
	| @TM^+~ @NT^expr~
	| @NT^function_call~
	| @TM^(~ @NT^expr~ @TM^)~
	;

@NT^literal~
	: [ @NT^expr~ ]
	| [ @NT^expr~ @TM^,~ @NT^expr~ ]
	| [ @NT^expr~ @TM^,~ ] @rem^@hfill// Right bound +@ensuremath^@infty~~
	| [ @TM^,~ @NT^expr~ ] @rem^@hfill// Left bound -@ensuremath^@infty~~
	| [ @TM^,~ ] @rem^@hfill// Entire real line~
	| [ @TM^entire~ ] @rem^@hfill// Entire real line~
	| [ ] @rem^@hfill// Empty interval~
	| < @NT^expr~ @TM^,~ @NT^expr~ >
	;

@NT^function_call~
	: @NT^name~ @TM^(~ @NT^expr~ @TM^)~
	| @NT^name~ @TM^(~ @NT^expr~ @TM^,~ @NT^expr~ @TM^)~
	| @NT^name~ @TM^(~ @NT^expr~ @TM^,~ @NT^expr~ @TM^,~ @NT^expr~ @TM^)~
	;
\end{Verbatim}

\noindent
Every string is an expression, computed with interval arithmetic: the
literals of intervals, the numbers, the constants and the functions combine
with the operators in any order, and a literal stands wherever a number may,
in a bound too.\newinvfive{} \code{1+[1,2]}, \code{2*cos([0,1])} and
\code{[cos([0,1]), 2]} are read. GAOL~4 read a string with two grammars, one
for the intervals between brackets and one for the numbers, and refused
\code{1+[1,2]} and \code{[1,2]+1+[1,2]} while it read \code{[1,2]*2}.
\code{[empty]} is the form \code{[\NT{expr}]} of the expression \code{empty}.

Spaces, tabs and newlines are not significant except in numbers, and the
case of letters is ignored, as IEEE 1788-2015 has it for its literals:
\code{[Empty]}, \code{[1, Inf]}, \code{SIN(1)} and \code{1E3} are
read.\newinvfive{} The ``$+$'' sign before numbers and \code{inf} is optional.

\paragraph{Numbers.} A \NT{number} is written in decimal, with an optional point
and exponent (\code{12}, \code{0.1}, \code{1e-3}), or in hexadecimal, with a
binary exponent (\code{0x1.3p-1}). Numbers are read exactly, whatever their
number of digits: a number that is a double is read as that double, and any
other one is replaced by the tightest interval of doubles containing
it.\newinvfive{} GAOL~4 read them with the \code{strtod()} of the C library,
which did not enclose \code{0.1} where it rounds to nearest in every rounding
direction (the C runtime of Windows, musl on 64-bit ARM processors).
\cindex{number!reading}%

\paragraph{The uncertain form.} An \NT{uncertain\_number} of IEEE 1788-2015
(12.11.3) is a decimal number followed by \code{?}, a radius in units of its
last digit (half a unit when it is left out, and an infinite radius for a
second \code{?}), an optional direction, \code{u} or \code{d}, which keeps only
the radius above or below, and an optional exponent: \code{3.56?1} is
\itv{3.55}{3.57}, \code{3.56?1e2} is \itv{355}{357}, and \code{-10??u} is
\itv{-10}{+\infty}.\newinvfive{} It is an interval literal of its own, not a
bound: \code{"[5?1]"} and \code{"[1, 3.56?1]"} throw \code{input\_format\_error},
as IEEE 1788-2015 has it (12.11.4), while \code{"3.56?1*2"} is
\itv{7.1}{7.14}.

\paragraph{Functions.} A \NT{name} is the name of a function of GAOL, with
its arguments, whatever the case of its letters: \code{interval(const char*)},
\code{gaol::textToInterval()} and \code{operator>{}>} read the names of
GAOL, and \code{gaol\_ieee1788::textToInterval()} those of IEEE 1788-2015
(\pxref{sec:ieee1788-text}), as each namespace has its own
\code{pow}.\newinvfive{} The names of GAOL are:

\begin{center}\small
\begin{tabular}{L{11cm}}
\code{sqr}, \code{sqrt}, \code{cbrt}, \code{pow}, \code{nth\_root}, \code{exp},
\code{exp2}, \code{exp10}, \code{expm1}, \code{exp2m1}, \code{exp10m1}, \code{log},
\code{log2}, \code{log10}, \code{log1p}, \code{log2p1}, \code{log10p1}, \code{sin},
\code{cos}, \code{tan}, \code{asin}, \code{acos}, \code{atan}, \code{atan2},
\code{sinh}, \code{cosh}, \code{tanh}, \code{asinh}, \code{acosh}, \code{atanh},
\code{sinpi}, \code{cospi}, \code{tanpi}, \code{asinpi}, \code{acospi},
\code{atanpi}, \code{atan2pi}, \code{hypot}, \code{rsqrt}, \code{sign}, \code{ceil},
\code{floor}, \code{trunc}, \code{integer}, \code{round\_ties\_to\_even},
\code{round\_ties\_to\_away}, \code{abs}, \code{min}, \code{max}, \code{inverse},
\code{fma}\\
\end{tabular}
\end{center}

\noindent
Each is the function of GAOL of that name (Chapters~\ref{chap:arithmetic}
and~\ref{chap:functions}); \code{pow} is the \code{pow} of the namespace \code{gaol}
(\pxref{sec:powers}), and \code{cbrt(x)} is \code{nth\_root(x,3)}. A name that
is not in the list, or a call with a wrong number of arguments, throws
\element{input\_format\_error}. The second argument of \code{nth\_root} shall
be a point interval that can be evaluated as an integer, negative ones
included; otherwise, an \element{invalid\_action\_error} exception is thrown
(or an error is reported, depending on whether exceptions were enabled or not
at configuration time---\pxref{sec:configuration}). GAOL~4 read the
trigonometric and hyperbolic functions, \code{exp}, \code{log}, \code{pow},
\code{sqrt} and \code{nth\_root} only.

\paragraph{Bounds.} The notations "$n$" and "[$n$]" are equivalent. Expressions
in bounds are evaluated using interval arithmetic, intervals included; the left
(resp. right) bound of the result is then used, depending on the side it
appeared in: \code{"[cos([0,1]), 2]"} is \itv{\cos 1}{2}. A bound left out is
infinite. The two bounds in the string ``\element{<a, b>}'' must be expressions
that evaluate to the same double even for different rounding directions: an
\element{input\_format\_error} exception is raised otherwise, as for
\code{"<0.1, 0.1>"}.

\paragraph{Infinities.} An interval literal whose left bound is $+\infty$ or
whose right bound is $-\infty$ is the empty set, as \code{numsToInterval} of
IEEE 1788-2015 has it: \code{"[inf]"} and \code{"[-inf, -inf]"} are
empty.\newinvfive{} The expression \code{"inf"} alone is still
\itv{\text{MAX}}{+\infty}, and \code{"-inf"} \itv{-\infty}{-\text{MAX}}
(\pxref{par:inf_inf_note}).

\paragraph{Threads.} Several threads may read strings at once, in parallel:
the lexer is reentrant and the parser pure, and each reading of a string has
a state of its own.\newinvfive{} GAOL~4 kept that state in global variables,
and crashed there.

\paragraph{Errors.} A string that is not an expression throws
\element{input\_format\_error}. An error found anywhere in the string stops the
reading, and no interval is given: \code{"[nth\_root(8,1.5)]+[1,2]"} throws,
where the two grammars of GAOL~4 gave \itv{-\infty}{+\infty}.\newinvfive

\begin{example}
interval x("[4, 6*7]");
interval y("[-inf, dmax]");
interval z("[3.14,3.15]/8", "[3.14,3.15]/7");
interval t("[3.14,3.15]/[7,8]");
cout << interval("[1,]") << ' ' << interval("[ ]") << ' '
     << interval("[Entire]") << endl;
@outputs^[1, inf] [empty] [-inf, inf]~
cout << interval("[0x1.3p-1, 2/3]") << ' '
     << interval("3.56?1") << endl;
@outputs^[0.59375, 0.6666666666666668] [3.549999999999999, 3.570000000000001]~
cout << interval("1+[1,2]") << ' '
     << interval("[cos([0,1]), 2]") << endl;
@outputs^[2, 3] [0.5403023058681396, 2]~
\end{example}

\section{Writing intervals}\label{sec:writing}
%%-------------------------

Intervals may be printed into a stream like any other C++ primitive type
by using the ``\code{<{}<}'' operator.

\begin{deffun}{ostream\&}{operator<{}<}{(ostream\& \var{out}, const interval\& I)}
Prints the interval \var{I} to the output stream \var{out}. The way the intervals
are actually displayed depends on the active format (see next section).
However, whatever the format, an empty interval is always displayed as \code{[empty]}
\cindex{empty!interval: printing}%
\index[idx]{empty@\code{[empty]}}%
\end{deffun}

\noindent
In the decimal formats, each bound is rounded outward, whatever the C library:
its magnitude is written by the C library, rounded to nearest, then compared
exactly with the bound, and its last digit is moved by one when it is on the
wrong side. The text read back therefore encloses the interval
written.\newinvfive{} GAOL~4 let the C library round each bound in the
direction it set, which the C runtime of Windows and musl do not honour.

\begin{deffun}{std::string}{exact\_string}{(const interval\& \var{I})}
  Returns what \code{operator<{}<} writes in the format \code{hexa} (see next
  section), without reading nor changing the global output format:
  \code{"[empty]"} for the empty set, and otherwise the two bounds in the
  hexadecimal-significand form of IEEE 1788-2015 (13.4.1), which
  \code{interval(const char*)} reads back bit for bit. This is the exact text
  representation of the standard, its \ieee{intervalToExact}.\newinvfive
\begin{example}
cout << exact_string(interval("0.1"));
@outputs^[0x1.9999999999999p-4, 0x1.999999999999ap-4]~
\end{example}
\end{deffun}

\subsection{Converting intervals to strings}
%%..........................................

For convenience, the \code{interval} class provides a conversion operator
into the standard C++ type \code{string}, which writes the interval in the
active format.
\cindex{interval!conversion to \code{string}}%
\index[idx]{string@\code{string}!converting interval to|see{interval}}

\begin{example}
  interval I(3,4);
  string s = "test line embedding " + string(I) + " as a string";
  @rem^// Now, s is "test line embedding [3, 4] as a string"~
\end{example}

\subsection{Output format}\label{sec:output-format}
%%........................

Intervals may be displayed following five different formats:
\begin{description}
\item[\element{agreeing}.] By printing all the digits that are the same in the left and right bounds
followed by an interval containing the remaining digits:
\begin{center}
``\element{3.141\textasciitilde[5926, 6001]}'' stands for
``\element{\itv{3.1415926}{3.1416001}}''
\end{center}
Note that if the bounds do not have any agreeing digit, there will still be a tilde before the bracketed part.
The digits of each bound are kept: GAOL~4 dropped from both bounds what followed
the last character of the left one that is not a zero, and wrote
\itv{1.25}{1.2567} \code{1.25\textasciitilde[, ]}.\newinvfive
\begin{example}
 interval::format(interval_format::agreeing);
 std::cout << interval(4,6) << ' '
           << interval(1.25,1.2567) << std::endl;
 @outputs^@textasciitilde[4., 6.] 1.25@textasciitilde[0, 67]~
\end{example}

\item[\element{bounds}.] By printing the bounds between square brackets.
Degenerate intervals whose left and right bound are equal are printed with angles:
\begin{example}
 interval::format(interval_format::bounds);
 std::cout << interval(3,5) << ' ' << interval(4) << std::endl;
 @outputs^[3, 5] <4, 4>~
\end{example}
A degenerate interval whose value has more digits than those displayed is
printed with its two bounds rounded outward: \code{interval(0.1)} is printed
\code{<0.1, 0.1000000000000001>}.

\item[\element{width}.] by printing their midpoint and their width:
\begin{example}
 interval::format(interval_format::width);
 std::cout << interval(4,5) << std::endl;
 @outputs^4.5 (+/- 0.5)~
\end{example}

\item[\element{center}.] by printing their midpoint only.

\item[\element{hexa}.] by printing their left and right bounds in the
hexadecimal-significand form of IEEE 1788-2015 (13.4.1), which is exact and which
\code{interval(const char*)} reads back bit for bit, without the round-off error of a
binary-to-decimal conversion: the recovery requirement of the standard
(13.4). GAOL~4 printed the sixteen hexadecimal digits of each double, which
its parser did not read.\newinvfive
\begin{example}
 interval::format(interval_format::hexa);
 std::cout << interval("0.1") << std::endl;
 @outputs^[0x1.9999999999999p-4, 0x1.999999999999ap-4]~
\end{example}

\end{description}

\medskip
Note that \emph{the \element{bounds} and \element{hexa} formats are the ones
read back as an input} (see previous section).
\medskip

The choice of the format to use is made through the following static methods:

\defstaticmethodx{interval}{void}{format}{(interval\_format::format\_t \var{f})}
\begin{defstaticmethod}{interval}{interval\_format::format\_t}{format}{(void)}
  The first form of the method allows modifying the format to use in subsequent
  printing of intervals. The second form reports what is the current form in use.
  It returns a value of  type \code{interval\_format::format\_t}
  (see below and \ref{sec:output-example} for an example of use).
  The format is global to the program: \code{exact\_string()} gives the
  hexadecimal text without changing it.
\end{defstaticmethod}

\begin{deftype}{struct}{interval\_format}
  Structure type used to choose the output format for intervals. It has
  five possible values of type \element{interval\_format::format\_t}:
\begin{itemize}
\item \element{interval\_format::agreeing}.
\item \element{interval\_format::bounds}.
\item \element{interval\_format::width}.
\item \element{interval\_format::center}.
\item \element{interval\_format::hexa}.
\end{itemize}
\end{deftype}

\begin{example}
  interval I(interval::pi());

  interval::format(interval_format::agreeing);
  cout << I << "\n";
  @rem^// Prints 3.14159265358979@textasciitilde[3, 4]~
  @rem^// The @textasciitilde[] part is dropped if the bounds agree on all digits~
  @smallskip
  interval::format(interval_format::bounds);
  cout << I << "\n";
  @rem^// Prints [3.141592653589793, 3.141592653589794]~
  @smallskip
  interval::format(interval_format::width);
  cout << I << "\n";
  @rem^// Prints 3.141592653589793 (+/- 2.220446049250313e-16)~
  @smallskip
  interval::format(interval_format::center);
  cout << I << "\n";
  @rem^// Prints 3.141592653589793~
  @smallskip
  interval::format(interval_format::hexa);
  cout << I << "\n";
  @rem^// Prints [0x1.921fb54442d18p+1, 0x1.921fb54442d19p+1]~
\end{example}

\subsection{Choosing the number of digits to display}
%%...................................................
You can manipulate the number of digits to print by using the \code{precision()} static
methods of the interval class:

\begin{defstaticmethod}{interval}{std::streamsize}{precision}{(void)}
Returns the current number of digits used for printing bounds
of intervals, 16 by default.

See example below.
\end{defstaticmethod}

\begin{defstaticmethod}{interval}{std::streamsize}{precision}{(std::streamsize \var{n})}
Set the number of digits to use for printing bounds to \var{n}. In addition,
returns the number of digits previously used.

See example below.
\end{defstaticmethod}


\subsection{Example}\label{sec:output-example}
%%..................

\begin{example}
#include <iostream>
#include <gaol/gaol.h>
using namespace gaol;

using std::cout;
using std::endl;

int main(void)
{
  interval::precision(4);
  interval::format(interval_format::bounds);
  cout << interval::pi() << endl;

  if (interval::format() != interval_format::bounds) {
    cout << interval::pi() << endl;
  } else {
    interval::precision(16);
    interval::format(interval_format::width);
    cout << interval::pi() << endl;
  }
  return 0;
}
\end{example}

\noindent
The previous program has the following output:

\begin{onscreen}
[3.141, 3.142]
3.141592653589793 (+/- 2.220446049250313e-16)
\end{onscreen}

The first call to \code{interval::format()} is unnecessary
since the default format is \code{interval\_format::bounds}.

\paragraph*{Note.} An interval written in decimal and read back gives an
interval that encloses it, whatever the precision chosen, the bounds being
rounded outward; the fewer the digits, the wider. The format \code{hexa} gives
the same interval back, bit for bit.

The \code{interval::pi()} constant is a predefined \dfn{canonical}
interval containing $\pi$ (\pxref{sec:interval-constants}). Here,
the width of the interval is equal to the $\epsilon$ of
the format.

The \code{interval\_format::width} format may be useful whenever the number of
digits displayed is insufficient to know whether the result is a single
floating-point number or an interval whose size is very small (consider
for example the first result above), because \emph{we have the guarantee that
if the actual width of an interval is greater than zero, the width
displayed will also be different from zero}. Another indication is that a degenerate
interval is displayed between angles.


\chapter{The names of IEEE 1788-2015}\label{chap:ieee1788}
%%===================================
\cindex{namespace!gaol\_ieee1788}\cindex{IEEE 1788-2015!names}%

GAOL names its operations its own way: the reverse function
\ieee{coshRev}$(c,x)$ of IEEE 1788-2015 is \code{acosh\_rel(c, x)},
\ieee{mulRev}$(b,c,x)$ is \code{div\_rel(c, b, x)}, with its arguments in
another order, and \ieee{roundTiesToEven}$(x)$ is
\code{round\_ties\_to\_even(x)}. The namespace \code{gaol\_ieee1788}, which
\file{gaol/gaol.h} brings along, gives each operation of the standard that GAOL
provides the name and the argument order the standard gives it, so that a
reader of the standard finds it without looking for its translation. It is not
in \code{gaol}, and it holds the type \code{interval} as well: one directive
lets a program use GAOL under the names of the standard.\newinvfive

\begin{example}
#include <iostream>
#include <gaol/gaol.h>
using namespace gaol_ieee1788;

int main(void)
{
  interval x = textToInterval("[1, 2]");
  interval y = mulRev(numsToInterval(2, 2), x);   @rem^// x / 2~
  interval z = sinPi(x) + rootn(y, 3);
  bool b = strictLess(x, entire());
  std::cout << std::boolalpha
            << x << ' ' << y << ' ' << z << ' ' << b << std::endl;
  return 0;
}
@outputs^[1, 2] [0.5, 1] [-0.2062994740159004, 1] false~
\end{example}

\noindent
A program opens \code{gaol} or \code{gaol\_ieee1788}, not both
(\pxref{sec:namespaces}). A name of the program's own that one of the standard
shadows, a constant \code{inf} for instance, is to be qualified:
\code{gaol\_ieee1788::inf(x)}. So is \code{less(x, y)} next to
\code{using namespace std;}, where it is ambiguous with the class template
\code{std::less}.

\section{Where the standard and GAOL differ}
%%-----------------------------------------

The functions of GAOL that have the name and the meaning the standard gives them
(\code{sin}, \code{exp}, \code{sqrt}, \code{min}\ldots) are the same functions in
\code{gaol\_ieee1788}. Where the standard and GAOL differ, the names of
\code{gaol\_ieee1788} follow the standard:

\begin{itemize}
\item \code{pow(x, y)} is the \ieee{pow} of Table 9.1, defined for $x>0$, and
  for $x=0$ when $y>0$, where the \code{pow} of \code{gaol} takes the integer
  power for an integer exponent, a negative base included
  (\pxref{sec:powers}). \code{pow(x, p)} with a number $p$ is
  \code{pow(x, [p])}, \code{pow(x, 2)} included: the power with an integer
  exponent is \code{pown(x, n)} in the standard, whose $n$ is an integer
  rather than an interval. An integer exponent beyond the \code{int}s gives
  the power of CORE-MATH at the bounds of $x$, where the \code{pow} of
  \code{gaol} gives \itv{-\infty}{+\infty};
\item \code{inf(x)} and \code{sup(x)} are $+\infty$ and $-\infty$ for the empty
  set (Table 10.2), where \code{left()} and \code{right()} are NaN, and
  \code{inf} returns $-0$ for a left bound 0 (12.12.8);
\item \code{isMember(m, x)} is \code{false} for an infinite $m$ (10.6.3);
\item \code{textToInterval} and \code{exactToInterval} read the names of the
  functions of the standard, and return the empty set for a string that is no
  interval (12.1.3), where \code{gaol::textToInterval} and the constructor of
  GAOL read the names of GAOL and throw (\pxref{sec:ieee1788-text});
\item \code{pownRev(c, x, p)} throws \code{std::invalid\_argument} for
  $p\leq0$, which GAOL does not provide, rather than give an interval it has
  not computed.
\end{itemize}

\begin{example}
cout << pow(interval(-4,-1), interval(2.0)) << ' '
     << pown(interval(-4,-1), 2) << ' '
     << pow(interval(2,3), interval(1e10)) << endl;
@outputs^[empty] [1, 16] [1.797693134862315e+308, inf]~
cout << inf(empty()) << ' ' << sup(empty()) << ' '
     << isMember(GAOL_INFINITY, entire()) << ' '
     << textToInterval("[1, 2") << endl;
@outputs^inf -inf false [empty]~
\end{example}

\section{The operations}
%%----------------------

The last argument $x$ of the reverse functions is optional, and
\itv{-\infty}{+\infty} when it is left out (10.5.4).

\begingroup\renewcommand{\arraystretch}{1.1}\small
\begin{longtable}{L{5.2cm}L{5.8cm}}
  \hline
  \multicolumn{1}{c}{\code{gaol\_ieee1788}} & \multicolumn{1}{c}{GAOL}\\
  \hline\hline
  \endhead
  \hline
  \endfoot
  \multicolumn{2}{l}{\textit{Constants and constructors (10.5.2, 10.5.8)}}\\
  \code{empty()}, \code{entire()} & \code{interval::emptyset()}, \code{interval::universe()}\\
  \code{numsToInterval(l, u)} & \code{interval(l, u)}\\
  \code{textToInterval(s)} & \code{interval(s)} with the names of the standard, the empty set for a wrong string\\
  \hline
  \multicolumn{2}{l}{\textit{Forward elementary functions (Table 9.1)}}\\
  \code{neg}, \code{add}, \code{sub}, \code{mul}, \code{div} & \code{-x}, \code{x + y}, \code{x - y}, \code{x * y}, \code{x / y}\\
  \code{recip(x)} & \code{inverse(x)}\\
  \code{sqr}, \code{sqrt}, \code{fma} & the same\\
  \code{pown(x, p)} & \code{gaol::pow(x, p)} for an \code{int} $p$\\
  \code{pow(x, y)}, \code{pow(x, p)} & the \ieee{pow} of the standard\\
  \code{exp}, \code{exp2}, \code{exp10}, \code{log}, \code{log2}, \code{log10} & the same\\
  \code{sin}, \code{cos}, \code{tan}, \code{asin}, \code{acos}, \code{atan}, \code{atan2} & the same\\
  \code{sinh}, \code{cosh}, \code{tanh}, \code{asinh}, \code{acosh}, \code{atanh} & the same\\
  \code{sign}, \code{ceil}, \code{floor}, \code{trunc} & the same\\
  \code{roundTiesToEven}, \code{roundTiesToAway} & \code{round\_ties\_to\_even}, \code{round\_ties\_to\_away}\\
  \code{abs}, \code{min}, \code{max} & the same\\
  \hline
  \multicolumn{2}{l}{\textit{Recommended functions (Table 10.5)}}\\
  \code{rootn(x, q)} & \code{nth\_root(x, q)}\\
  \code{expm1}, \code{exp2m1}, \code{exp10m1} & the same\\
  \code{logp1}, \code{log2p1}, \code{log10p1} & \code{log1p}, \code{log2p1}, \code{log10p1}\\
  \code{hypot}, \code{rSqrt} & \code{hypot}, \code{rsqrt}\\
  \code{sinPi}, \code{cosPi}, \code{tanPi} & \code{sinpi}, \code{cospi}, \code{tanpi}\\
  \code{asinPi}, \code{acosPi}, \code{atanPi}, \code{atan2Pi} & \code{asinpi}, \code{acospi}, \code{atanpi}, \code{atan2pi}\\
  \hline
  \multicolumn{2}{l}{\textit{Reverse functions (Table 10.1)}}\\
  \code{sqrRev(c, x)}, \code{absRev(c, x)} & \code{sqrt\_rel(c, x)}, \code{invabs\_rel(c, x)}\\
  \code{pownRev(c, x, p)}, $p\geq1$ & \code{nth\_root\_rel(c, p, x)}\\
  \code{sinRev}, \code{cosRev}, \code{tanRev(c, x)} & \code{asin\_rel}, \code{acos\_rel}, \code{atan\_rel(c, x)}\\
  \code{coshRev(c, x)} & \code{acosh\_rel(c, x)}\\
  \code{sinhRev}, \code{tanhRev(c, x)} (not in the standard, named after \ieee{coshRev}) & \code{asinh\_rel}, \code{atanh\_rel(c, x)}\\
  \code{mulRev(b, c, x)} & \code{div\_rel(c, b, x)}\\
  \hline
  \multicolumn{2}{l}{\textit{Cancellative operations and set operations (10.5.6, 10.5.7)}}\\
  \code{cancelMinus}, \code{cancelPlus} & \code{cancel\_minus}, \code{cancel\_plus}\\
  \code{intersection(x, y)}, \code{convexHull(x, y)} & \code{x \& y}, \code{x | y}\\
  \hline
  \multicolumn{2}{l}{\textit{Numeric functions (Table 10.2)}}\\
  \code{inf(x)}, \code{sup(x)} & \code{x.left()}, \code{x.right()}, but for the empty set\\
  \code{mid}, \code{wid}, \code{rad}, \code{mag}, \code{mig(x)} & \code{x.midpoint()}, \code{width()}, \code{rad()}, \code{mag()}, \code{mig()}\\
  \code{midRad(x, m, r)} & \code{x.mid\_rad(m, r)}\\
  \hline
  \multicolumn{2}{l}{\textit{Boolean functions (10.5.10, Tables 10.3 and 10.4)}}\\
  \code{isEmpty(x)}, \code{isEntire(x)} & \code{x.is\_empty()}, \code{x.is\_entire()}\\
  \code{equal(a, b)}, \code{subset(a, b)} & \code{a.set\_eq(b)}, \code{a.set\_leq(b)}\\
  \code{interior(a, b)}, \code{disjoint(a, b)} & \code{a.set\_le(b)}, \code{a.set\_disjoint(b)}\\
  \code{less(a, b)}, \code{strictLess(a, b)} & \code{a.less(b)}, \code{a.strictly\_less(b)}\\
  \code{precedes(a, b)}, \code{strictPrecedes(a, b)} & \code{a.certainly\_leq(b)}, \code{a.certainly\_le(b)}\\
  \code{isCommonInterval(x)}, \code{isSingleton(x)} & \code{x.is\_common\_interval()}, \code{x.is\_a\_double()}\\
  \code{isMember(m, x)} & \code{x.set\_contains(m)}, false for an infinite $m$\\
  \hline
  \multicolumn{2}{l}{\textit{Input and output (Clause 13)}}\\
  \code{intervalToText(x)} & \code{operator<{}<} in the current format\\
  \code{intervalToExact(x)} & \code{exact\_string(x)}\\
  \code{exactToInterval(s)} & \code{textToInterval(s)}\\
  \hline
  \multicolumn{2}{l}{\textit{Expressions}}\\
  \code{pown(e, n)}, \code{pow(e1, e2)} & the expressions of \code{pown} and of the \ieee{pow} of the standard\\
\end{longtable}
\endgroup

\section{Reading with the names of the standard}\label{sec:ieee1788-text}
%%-----------------------------------------------

\code{textToInterval(\var{s})} reads a string as \code{interval(const char*)}
does (\pxref{sec:input-format}), with the names of the functions of IEEE
1788-2015: those of Tables 9.1 and 10.5 that GAOL provides, whatever the case
of their letters.\newinvfive

\begin{center}\small
\begin{tabular}{L{11cm}}
\code{neg}, \code{add}, \code{sub}, \code{mul}, \code{div}, \code{recip},
\code{sqr}, \code{sqrt}, \code{fma}, \code{pown}, \code{pow}, \code{exp},
\code{exp2}, \code{exp10}, \code{log}, \code{log2}, \code{log10}, \code{sin},
\code{cos}, \code{tan}, \code{asin}, \code{acos}, \code{atan}, \code{atan2},
\code{sinh}, \code{cosh}, \code{tanh}, \code{asinh}, \code{acosh}, \code{atanh},
\code{sign}, \code{ceil}, \code{floor}, \code{trunc}, \code{roundTiesToEven},
\code{roundTiesToAway}, \code{abs}, \code{min}, \code{max}, \code{rootn},
\code{expm1}, \code{exp2m1}, \code{exp10m1}, \code{logp1}, \code{log2p1},
\code{log10p1}, \code{hypot}, \code{rSqrt}, \code{sinPi}, \code{cosPi},
\code{tanPi}, \code{asinPi}, \code{acosPi}, \code{atanPi}, \code{atan2Pi}\\
\end{tabular}
\end{center}

\noindent
\code{pow} is the \ieee{pow} of Table 9.1, and the second argument of
\code{pown} and \code{rootn} is an integer. The names of GAOL alone
(\code{nth\_root}, \code{cbrt}, \code{inverse}, \code{integer}, \code{log1p},
\code{round\_ties\_to\_even}\ldots) are no functions there, and give the
empty set, as any string that is no interval. As \code{pow}, the function
\code{textToInterval} is in the namespace \code{gaol} too, where it reads the
names of GAOL and throws \code{input\_format\_error} for a wrong string: a
program calls the one of the namespace it opens.

\begin{example}
cout << textToInterval("pown([2,5],5)") << ' '
     << textToInterval("pow([-4,-1],2)") << ' '
     << textToInterval("nth_root(8,3)") << endl;
@outputs^[32, 3125] [empty] [empty]~
\end{example}

\noindent
With \code{using namespace gaol;}, \code{textToInterval("pow([-4,-1],2)")} is
\itv{1}{16}, and \code{textToInterval("pown([2,5],5)")} throws.

\section{What is not provided}
%%----------------------------

GAOL has no decorations: only bare intervals are provided (Clause 11). Nor does
it provide \ieee{mulRevToPair} (10.5.5), \ieee{powRev1}, \ieee{powRev2},
\ieee{atan2Rev1} and \ieee{atan2Rev2} (Table 10.1), \ieee{pownRev} for
$p\leq0$, \ieee{compoundm1} (Table 10.5), the slope functions (Table 10.6),
\ieee{overlap} (10.6.4), the reduction operations \ieee{sum}, \ieee{dot},
\ieee{sumSquare} and \ieee{sumAbs} (12.12.12), and the exact operations of
12.13.5. The end of \file{gaol/gaol\_ieee1788.h} keeps this list.

\chapter{Accuracy of the operations}\label{chap:accuracy}
%%=================================
\cindex{accuracy!of the operations}\cindex{tightness}%

IEEE 1788-2015 requires an implementation to document the tightness of each of
its interval operations (12.10.3). This chapter does so for \code{interval},
from the algorithm of each operation; the tests of \file{tests} check each
statement on random and special arguments (\pxref{sec:tests}), and the page
\file{doc/accuracy.md} gives, for each, the algorithm and the largest distance
the tests found.\newinvfive

\section{Accuracy modes}
%%----------------------

IEEE 1788-2015 defines three accuracy modes (12.10.1), for an operation $f$,
its exact value $f^*$, and an input box $x$:

\begin{itemize}
\item \dfn{tightest}: $f(x)$ is the hull, in doubles, of $f^*(x)$: each bound is
  the nearest double on its side of the exact bound;
\item \dfn{accurate}: $f(x)$ encloses $f^*(x)$, and is within
  \[\text{nextOut}(f_{\text{tightest}}(\text{nextOut}(x))),\]
  nextOut moving
  each bound one double outward: roughly, $f(x)$ is no wider than the tightest
  result for an input one double wider, plus one double on each side;
\item \dfn{valid}: $f(x)$ encloses $f^*(x)$.
\end{itemize}

\noindent
The standard requires the basic operations, the integer functions and
\code{abs}, \code{min} and \code{max} to be tightest; the other elementary
functions, the reverse functions and the recommended functions to be valid,
and recommends them to be accurate (12.10.2). Below, ``within $k$ doubles''
means that each bound is at most $k$ doubles beyond the tightest one.

GAOL bounds its elementary functions with CORE-MATH, which is correctly rounded
in the rounding direction in effect: computing upward, GAOL takes the value at
a bound as the upper bound, with nothing to add, and the double below it as the
lower one. Those are the tightest bounds, the second one unless the exact value
is a double, which the operations give exactly where it happens. Where an
algorithm evaluates the function at the bounds of the interval, the bounds of
the interval are therefore the tightest ones.

The statements hold for GAOL as the three builds make it, with the code
including its headers compiled with the flags of interval arithmetic
(\pxref{sec:compiling}), for the SSE2 and the FPU intervals, whatever the
rounding direction the calling code left, and on every architecture and with
every compiler alike. The intervals of floats are not covered.

\section{Tightness of each operation}
%%-----------------------------------

\begingroup\renewcommand{\arraystretch}{1.1}\small
\begin{longtable}{L{4.2cm}L{6.8cm}}
  \hline
  \multicolumn{1}{c}{Operations} & \multicolumn{1}{c}{Tightness}\\
  \hline\hline
  \endhead
  \hline
  \endfoot
  \multicolumn{2}{l}{\textit{Basic operations (Table 9.1): tightest required}}\\
  \code{-x}, \code{+}, \code{-}, \code{*}, \code{/}, \code{inverse}, \code{sqr} & tightest\\
  \code{sqrt} & tightest, whatever the rounding of the square root of the C library\\
  \code{fma} & tightest\\
  \hline
  \multicolumn{2}{l}{\textit{Integer and absmax functions: tightest required}}\\
  \code{floor}, \code{ceil}, \code{integer}, \code{sign}, \code{trunc}, \code{round\_ties\_to\_even}, \code{round\_ties\_to\_away}, \code{abs}, \code{min}, \code{max} & tightest, and exact, in every rounding direction\\
  \hline
  \multicolumn{2}{l}{\textit{Powers, exponentials and logarithms: valid required, accurate recommended}}\\
  \code{pow(x, n)}, integer \var{n} (\ieee{pown}) & $n=0$, 1, 2: tightest. $n\geq3$: the tightest bounds, or one double beyond where the power is within $n\cdot2^{-104}$ of a double; exact where the power is a double; within about $2(n-1)$ doubles below $2^{-968}$ and at the overflow. Negative $n$: valid, within 2 doubles where $x^n$ is normal\\
  \code{pow(x, y)} (\ieee{pow}) & tightest at the corners of the box with finite bounds, but for a lower bound one double below where the power at the corner is a double; valid with infinite bounds\\
  \code{exp}, \code{exp2}, \code{exp10}, \code{log}, \code{log2}, \code{log10} & tightest; exact where the value is a double ($\exp(0)$, $\log(1)$, $2^n$, $10^n$ for $0\leq n\leq22$\ldots)\\
  \hline
  \multicolumn{2}{l}{\textit{Trigonometric and hyperbolic functions: valid required, accurate recommended}}\\
  \code{cos}, \code{sin} & tightest where the function is evaluated at the bounds, at every magnitude; $-1$ and $1$ exactly at the extrema; \itv{-1}{1} from the width $2\underline{\pi}$ on, $\underline{\pi}$ being the double below $\pi$\\
  \code{tan} & tightest between two poles, at every magnitude; \itv{-\infty}{+\infty} exactly when a pole is within $x$, but for the widths between $\underline{\pi}$ and $\pi$\\
  \code{asin}, \code{acos}, \code{atan}, \code{atan2} & tightest; exactly 0, or the tightest enclosures of $\pm\pi/4$, $\pm\pi/2$ and $\pi$, where the value is one of them\\
  \code{sinh}, \code{cosh}, \code{tanh}, \code{asinh}, \code{acosh}, \code{atanh} & tightest\\
  \hline
  \multicolumn{2}{l}{\textit{Recommended functions (Table 10.5)}}\\
  \code{nth\_root(x, n)}, $n>0$ & $n=1$, 2, 3: tightest. Otherwise accurate: the tightest bounds or one double beyond, proved with integer powers\\
  \code{nth\_root(x, q)}, $q<0$ & accurate: the root above, then the tightest division\\
  \code{expm1}, \code{exp2m1}, \code{exp10m1}, \code{log1p}, \code{log2p1}, \code{log10p1}, \code{hypot}, \code{rsqrt} & tightest\\
  \code{sinpi}, \code{cospi}, \code{tanpi}, \code{asinpi}, \code{acospi}, \code{atanpi}, \code{atan2pi} & tightest, at every magnitude\\
  \hline
  \multicolumn{2}{l}{\textit{Reverse functions (Table 10.1): valid required, accurate recommended}}\\
  \code{sqrt\_rel}, \code{nth\_root\_rel}, \code{acosh\_rel}, \code{asinh\_rel}, \code{atanh\_rel} & accurate\\
  \code{invabs\_rel} & tightest\\
  \code{asin\_rel}, \code{acos\_rel}, \code{atan\_rel} & accurate, or one double beyond; a bound of $x$ beyond $2^{52}$ is kept\\
  \code{div\_rel}, \code{\%} & accurate; tightest where the bounds of the quotients are doubles\\
  \code{cancel\_minus}, \code{cancel\_plus} & tightest\\
  \hline
  \multicolumn{2}{l}{\textit{Set operations, constructors, numeric and boolean functions}}\\
  \code{\&}, \code{|}, \code{interval(l, u)} & tightest (exact)\\
  \code{interval(s)}, \code{operator>{}>} & tightest for the literals of IEEE 1788-2015; valid for expressions, computed with the operations above\\
  \code{midpoint()}, \code{width()}, \code{rad()}, \code{mag()}, \code{mig()} & as 12.12.8 requires\\
  \code{mid()}, \code{hausdorff()} & tightest\\
  relations & exact, as Tables 10.3 and 10.4\\
  \code{operator<{}<} & \code{hexa}: exact, read back bit for bit. Decimal formats: valid, and the tightest with the digits asked for where the C library rounds to nearest as it should\\
\end{longtable}
\endgroup

\noindent
Where a bound of \var{J} is one double from the image of a bound of the
preimage, the reverse functions keep a point of \var{I} whose image is just
outside \var{J}: \code{sqrt\_rel(\itv{1+2^{-52}}{+\infty}, \itv{-1}{1+2^{-52}})}
is \itv{-1}{1+2^{-52}}, whose tightest enclosure is \itv{1}{1+2^{-52}}. They are
accurate all the same, the widened \var{J} of nextOut holding that image.

Before GAOL v5, the elementary functions were correctly rounded to nearest and
moved one double outward, which is at most one double beyond the tightest
bound, and the hyperbolic functions were taken from the mathematical library of
the system, whose values are sometimes further than one double from the exact
ones.


\chapter{Floating-point numbers}
%%==============================

\section{Floating-point constants}\label{sec:floating-point-constants}
%%--------------------------------

In addition to the constants available through \code{numeric\_limits<double>},
GAOL defines the following \code{double} constants:

\begin{center}
  \renewcommand{\arraystretch}{1.1}
  \begin{tabular}{lc}
    \hline
    \multicolumn{1}{c}{\bfseries Constant (\texttt{double})} & \multicolumn{1}{c}{\bfseries Value}\\
    \hline\hline
    \element{two\_pi} & \roundNearest{2\pi}\\
    \element{pi} & \roundNearest{\pi}\\
    \element{half\_pi} & \roundNearest{\frac{\pi}{2}}\\
    \element{pi\_dn} & \roundDn{\pi}\\
    \element{pi\_up} & \roundUp{\pi}\\
    \element{half\_pi\_dn} & \roundDn{\frac{\pi}{2}}\\
    \element{half\_pi\_up} & \roundUp{\frac{\pi}{2}}\\
    \element{ln2\_dn} & \roundDn{\ln 2}\\
    \element{ln2\_up} & \roundUp{\ln 2}\\
    \element{two\_power\_51} & $2^{51}$\\
    \element{two\_power\_52} & $2^{52}$\\
    \element{GAOL\_NAN} & NaN (quiet)\\
    \element{GAOL\_INFINITY} & $+\infty$ \\
    \hline
  \end{tabular}
\end{center}

\section{Floating-point functions}
%%--------------------------------

\begin{deffun}{bool}{feven}{(double \var{x})}
  Returns true whenever \var{x} is even.

  This function should not be used with infinity and NaN arguments.

  \begin{example}
cout << feven(3.0) << ' ' << feven(4.0) << ' ' << feven(3.5);
@outputs^false true false~
  \end{example}
\end{deffun}

\begin{deffun}{double}{next\_float}{(double \var{x})}
  Returns the smallest \code{double} greater than \var{x}.
\end{deffun}

\begin{deffun}{double}{previous\_float}{(double \var{x})}
  Returns the greatest \code{double} smaller than \var{x}.
\end{deffun}

\begin{deffun}{bool}{is\_signed}{(double \var{x})}
Returns true whenever \var{x} is signed. No provision is made concerning the
fact that \var{x} is a NaN. If you only want to test for negative
numbers (and $-0$), you will have to test also whether \var{x} is a NaN
by using \code{std::isnan()}.
\end{deffun}

\begin{deffun}{double}{minimum}{(double \var{x}, double \var{y})}
  Returns the minimum \code{double} value of \var{x}
  and \var{y}. This function is commutative and returns $-0$ when comparing
  $-0$ and $+0$, i.e.:
\begin{equation*}
\left\{\begin{array}{lll}
    \min(x,y) &=\min(y,x), &\forall x\neq\text{NaN}, \forall y\neq\text{NaN}\\
    \min(x,\text{NaN}) &=\min(\text{NaN},x)=\text{NaN}, &\forall x\\
    \min(-0,0) &=\min(0,-0)=-0
  \end{array}\right.
\end{equation*}
\end{deffun}

\begin{deffun}{double}{maximum}{(double \var{x},double \var{y})}
  Returns the maximum \code{double} value of \var{x}
  and \var{y}. This function is commutative and returns $+0$ when comparing
  $-0$ and $+0$, i.e.:
\begin{equation*}
  \left\{\begin{array}{lll}
      \max(x,y) &= \max(y,x), &\forall x\neq\text{NaN}, \forall y\neq\text{NaN} \\
      \max(x,\text{NaN})&=\max(\text{NaN},x)=\text{NaN}, &\forall x\\
      \max(-0,0)&=\max(0,-0)=0
    \end{array}\right.
\end{equation*}
\end{deffun}

\begin{defmacro}{ULONGLONGINT}
  Macro standing for an unsigned integral data type with a size equal to 8 bytes
  (usually \code{unsigned long long int}).
\end{defmacro}

\begin{deffun}{ULONGLONGINT}{nb\_fp\_numbers}{(double \var{a}, double \var{b})}
Returns the number of floating-point numbers in the interval
\itv{\var{a}}{\var{b}}. In particular, we have:
\begin{itemize}
\item \code{nb\_fp\_numbers(a,next\_float(a)) == 2}
\item \code{nb\_fp\_numbers(a,a) == 1}
\end{itemize}

\paragraph{Note.} As a precondition, \var{a} shall be lower or equal to \var{b}.

Returns \code{numeric\_limits<ULONGLONGINT>::max()} if either
\var{a} or \var{b} is a NaN or an infinity, or if \var{a} is greater than
\var{b}. In addition, raises
an \code{invalid\_action\_error} exception (\pxref{sec:exceptions}) or calls
\code{gaol\_error} depending on the way the library was configured.
\begin{example}
cout << nb_fp_numbers(1.0, next_float(1.0)) << ' '
     << nb_fp_numbers(1.0, 2.0);
@outputs^2 4503599627370497~
\end{example}
\end{deffun}

\chapter{Manipulating the FPU}
%%============================

GAOL provides functions to set the rounding direction of the floating-point
unit. The operations of GAOL do not need them: each sets the rounding
direction upward itself (\pxref{sec:rounding-direction}).

On x86 processors, with GCC, Clang and Visual C++, they write the control
registers of the x87 and SSE units directly (\code{fnstcw}/\code{fldcw},
\code{stmxcsr}/\code{ldmxcsr}), rather than calling \code{fesetround()}, which
cost 50 to 250~ns per call with the C runtime of Windows; with Visual C++ for
x64, the SSE register only, the x87 unit being unused there.\newinvfive{}
Elsewhere (ARM, POWER, s390x, RISC-V), they call \code{fesetround()}.

\section{Rounding functions}
%%--------------------------

\begin{deffun}{void}{round\_downward}{(void)}
Sets the rounding direction mode towards $-\infty$.
\end{deffun}

\begin{deffun}{void}{round\_nearest}{(void)}
  Sets the rounding direction mode to the nearest/even.
\end{deffun}

\begin{deffun}{void}{round\_upward}{(void)}
  Sets the rounding direction mode to $+\infty$.
\end{deffun}

\section{Saving the rounding direction}
%%-------------------------------------

\begin{deffun}{unsigned short}{get\_fpu\_cw}{(void)}
  Returns the rounding direction, as \code{fegetround()} of \file{<fenv.h>}
  gives it.
\end{deffun}

\begin{deffun}{void}{reset\_fpu\_cw}{(unsigned short \var{st})}
  Sets the rounding direction \var{st}, which \code{get\_fpu\_cw()} returned,
  with \code{fesetround()}.

  GAOL~4 read and wrote the control word of the x87 unit with these functions
  on x86 Linux and macOS, which left the direction of the SSE instructions
  alone, and 16 bits of the control register on 64-bit ARM, which left out its
  rounding bits.\newinvfive
\end{deffun}

\chapter{Version information}
%%===========================

The library provides four constants to allow programs to determine at runtime with which version
they are dynamically linked with. The versioning scheme adopted is the one used by the Apache
Software Foundation described at \url{http://apr.apache.org/versioning.html}.

\begin{defconst}{unsigned int}{version\_major}
  Major version of the library: 5 for GAOL v5.
\end{defconst}

\begin{defconst}{unsigned int}{version\_minor}
  Minor version of the library.
\end{defconst}

\begin{defconst}{unsigned int}{version\_micro}
  Patch version of the library.
\end{defconst}

\begin{defconst}{const char *const}{version}
  Version of the library as a string.
\begin{example}
cout << gaol::version;
@outputs^5.0.0~
\end{example}
\end{defconst}

\chapter{Additional functions}
%%============================

The following functions are utility functions not necessarily related to
intervals or floating-point numbers.

\begin{defTfun}{template <typename T>}{bool}{odd}{(const \var{T}\& \var{x})}
  Returns \code{true} if \var{x} is odd and false otherwise. The \var{T}
  type may be any type providing the \code{\&} (``bitwise \code{and}'') operator
  with the same semantics as the one for \code{int}s.
\end{defTfun}

\begin{defTfun}{template <typename T>}{bool}{even}{(const \var{T}\& \var{x})}
  Returns \code{true} if \var{x} is even and false otherwise. The \var{T}
  type may be any type providing the \code{\&} (``bitwise \code{and}'') operator
  with the same semantics as the one for \code{int}s.
\end{defTfun}

\chapter{Error handling}\label{chap:errors}
%%======================

A program that uses GAOL may report errors in two different ways:
\begin{itemize}
\item by throwing an exception;
\item or by printing a message and aborting.
\end{itemize}

The mechanism in use depends on the way the library is configured. If
you use the option \option{-{}-enable-exceptions=yes}, the default, all errors are
reported through exception throwing; otherwise, a message is printed to the
standard error, and the program aborts. The CMake build always reports errors
by exceptions. Relying on exceptions is more in the C++ spirit,
though it may incur some overhead.

It is up to the user to comply with this mechanism when adding error reporting
code to ones program. GAOL defines the following macro to be used
whenever one wants to report an error.

\begin{defmacroA}{gaol\_ERROR}{(\var{excep},\var{msg})}
  The behavior of the macro depends on the value chosen for the option
  \option{-{}-enable-exceptions}: if exceptions are enabled, exception
  \code{gaol\_core::\var{excep}} is raised with the message \var{msg}; otherwise, the
  program aborts with message \var{msg}.

\begin{example}
interval x;

@emph^[Code manipulating x]~

if (x.is_empty()) {
   gaol_ERROR(invalid_action_error,"Emptiness of one interval");
}
\end{example}
\end{defmacroA}

The \code{gaol\_error()} function is defined as follows:

\deffunx{void}{gaol\_error}{(const char *const \var{err})}
\begin{deffun}{void}{gaol\_error}{(const char *\var{file}, int \var{line},
const char *\var{err})}
Displays a message on the standard error output. The ternary version should be
called with the \code{GAOL\_FILE\_POS} macro for the first two parameters.
\end{deffun}

The \code{GAOL\_FILE\_POS} macro is described in the next section.

\section{Exceptions}\label{sec:exceptions}
%%------------------

The library defines \code{gaol\_exception} as a class to be used as a base
class for all GAOL exceptions.
All of them provide at least the name of the file and the line number from
where the exception has been thrown. As a facility, GAOL defines the following
macro:

\begin{defmacro}{GAOL\_FILE\_POS}
Expands itself into the first two arguments of any constructor for
\code{gaol\_exception} or one of its derived classes:

\begin{example}
if (@emph^[some condition]~) {
  throw gaol_exception(GAOL_FILE_POS,
                       "No additional information");
}
\end{example}
\end{defmacro}

All GAOL exceptions can be sent to an output stream through the ``\code{<{}<}''
operator.

\subsection{The \code{gaol\_exception} exception}\label{sec:gaol_exception}
%%...............................................
The \code{gaol\_exception} class is the base class from which derive all
GAOL exceptions. It inherits from the C++ standard class \code{exception}.

Every exception class deriving from it must at least provide the name of
the file and the line where the corresponding exception was thrown. As a
consequence, the constructors for \code{gaol\_exception} are as follows:

\defmethodx{gaol\_exception}{}{gaol\_exception}{(const char* \var{f}, unsigned \var{l})}
\begin{defmethod}{gaol\_exception}{}{gaol\_exception}{(const char* \var{f}, unsigned \var{l}, \\\hbox{}\hfill const char* \var{e})}
Constructs a \code{gaol\_exception} being thrown from file \var{f} at
line \var{l}. The second form permits adding some explanatory string \var{e}.
\end{defmethod}

\noindent
The class offers the following accessors:

\begin{defmethod}{gaol\_exception}{const char*}{file}{(void) const}
  Return the name of the file from where the exception was thrown.
\end{defmethod}

\begin{defmethod}{gaol\_exception}{unsigned int}{line}{(void) const}
  Returns the line number in the file from where the exception was thrown.
\end{defmethod}

\begin{defmethod}{gaol\_exception}{const char* const}{explanation}{(void) const}
  Returns a string explaining why the exception was thrown. Returns an empty
  string if no additional information was provided.
\end{defmethod}

\subsection{The \code{input\_format\_error} exception}
%%....................................................
The \code{input\_format\_error} exception is thrown whenever one attempts to
create an interval from an invalid string. This situation may occur when
reading an interval from a stream with the \code{>{}>} operator, or when
creating an interval from a string.

This class, as all GAOL exceptions, derives from \code{gaol\_exception}
(\pxref{sec:gaol_exception}). Its constructors have the same format
than the ones for \code{gaol\_exception}, namely:

\defmethodx{input\_format\_error}{}{input\_format\_error}{(const char* \var{f}, unsigned \var{l})}
\begin{defmethod}{input\_format\_error}{}{input\_format\_error}{(const char* \var{f}, unsigned \var{l}, const char* \var{e})}
Constructs an \code{input\_format\_error} being thrown from file \var{f} at
line \var{l}. The second form permits adding some explanatory string \var{e}.
\end{defmethod}
The methods of the class are inherited from \code{gaol\_exception}
(\pxref{sec:gaol_exception}).

\begin{example}
try {
  interval x("sin(1)+");
} catch (const gaol::input_format_error& e) {
  cout << e.explanation();
}
@outputs^Syntax error in interval initialization: sin(1)+~
\end{example}

\subsection{The \code{unavailable\_feature\_error} exception}
%%---------------------------------------------------------

This exception is thrown whenever an unavailable feature is requested. No
operation of GAOL v5 throws it: \code{atan2()}, which did, is
implemented.\newinvfive

This class, as all GAOL exceptions, derives from \code{gaol\_exception}.
(\pxref{sec:gaol_exception}). Its constructors have the same format
than the ones for \code{gaol\_exception}, namely:

\defmethodx{unavailable\_feature\_error}{}{unavailable\_feature\_error}{\\\hbox{}\hfill (const char* \var{f}, unsigned \var{l})}
\begin{defmethod}{unavailable\_feature\_error}{}{unavailable\_feature\_error}{\\\hbox{}\hfill (const char* \var{f}, unsigned \var{l}, const char* \var{e})}
Constructs an \code{unavailable\_feature\_error} being thrown from file \var{f} at
line \var{l}. The second form permits adding some explanatory string \var{e}.
\end{defmethod}

The methods of the class are inherited from \code{gaol\_exception} (\pxref{sec:gaol_exception}).

\subsection{The \code{invalid\_action\_error} exception}
%%------------------------------------------------------

This exception is thrown whenever a function is called with invalid arguments
(e.g. calling \code{nb\_fp\_numbers()} with NaNs as parameters).

This class, as all GAOL exceptions, derives from \code{gaol\_exception}
(\pxref{sec:gaol_exception}). Its constructors have the same format
than the ones for \code{gaol\_exception}, namely:

\defmethodx{invalid\_action\_error}{}{invalid\_action\_error}{(const char* \var{f},\\\hbox{}\hfill  unsigned \var{l})}
\begin{defmethod}{invalid\_action\_error}{}{invalid\_action\_error}{(const char* \var{f},\\\hbox{}\hfill  unsigned \var{l}, const char* \var{e})}
Constructs an \code{invalid\_action\_error} being thrown from file \var{f} at
line \var{l}. The second form permits adding some explanatory string \var{e}.
\end{defmethod}

The methods of the class are inherited from \code{gaol\_exception} (\pxref{sec:gaol_exception}).


\section{Warnings}
%%----------------

\deffunx{void}{gaol\_warning}{(const char *\var{warn})}
\begin{deffun}{void}{gaol\_warning}{(const char *\var{file}, int \var{line}, const char *\var{warn})}
Prints the message \var{warn} on the standard error output. The second form
should be called with the \code{GAOL\_FILE\_POS} macro for the first two parameters.
\end{deffun}

\chapter{Debugging facilities}\label{sec:debugging-facilities}
%%============================

The debugging facilities described hereunder are available only if
GAOL has been configured with the debugging facilities enabled: the
\option{-{}-enable-debug} option of \cmd{configure} (\pxref{sec:configuration}),
or \option{-Denable-debug=true} with meson. The CMake build does not enable
them.

\begin{deffun}{int}{debug\_level}{}
  Global variable used to remember the current value of the debugging
  level.  This variable is set when initializing the library. The
  variable is declared in the \code{gaol\_core} namespace, which \code{gaol}
  imports.
\end{deffun}

\begin{defmacroA}{GAOL\_DEBUG}{(\var{lvl},\var{cmd})}
  Executes \var{cmd} if \var{lvl} is lower or equal to the current debugging
  level (see the variable \code{debug\_level} above).

  This macro defaults to nothing if the library was not configured with the
debugging facilities.
\end{defmacroA}

\noindent
A possible use for this macro is as follows:

\begin{example}
interval x(-10,10);

@emph^[Some code]~

GAOL_DEBUG(1,cout << "The value of x is " << x);
x += interval(3.5,4.5);
GAOL_DEBUG(2,cout << "Now the value of x is " << x);
\end{example}

The first message will be displayed whenever the debugging level is at least
1, and the second one only if it is at least 2.

\begin{defmacroA}{GAOL\_ASSERT}{(\var{cond})}
Tests whether \var{cond} holds. Aborts with an error message if it is not the case.

This macro defaults to nothing if the library was not configured with the
debugging facilities.
\end{defmacroA}

\noindent
A possible use for this macro is as follows:

\begin{example}
int x;
cout << "Give an integer no greater than 5: ";
cin >> x;
GAOL_ASSERT(x <= 5);
\end{example}

\chapter{Profiling}
%%=================

The following functions permit computing the time used for a computation. The
returned times are \emph{user} times, meaning that delays induced by
input/output operations and freezing during CPU switches in multi-programming
environments are not taken into account.

If you need to keep track of several events, consider using an object of the
\code{timepiece} class (\pxref{sec:timepiece}) instead of calling
directly the functions below.

\paragraph*{Warning.} The precision of the timing functions depends on the platform
used. What is more, despite the fact that the reported times are user
times, they may vary from an execution to another, and can get larger on
heavily loaded machines. \code{tests/performance.cpp} shows how GAOL v5 measures
the time of its operations, in nanoseconds.


\begin{deffun}{long}{get\_time}{(void)}
Returns the time in milliseconds since a certain unspecified moment.

This function should only be used to compute differences between two
calls since the starting point may vary depending on the availability of
\code{clock()} or \code{getrusage()} on the system.


\paragraph*{Warning.} if the function in use is the standard \code{clock()},
the time returned will wrap approximately every 72 minutes. Consequently, it is
not safe to use \code{get\_time()} in that case for processes requiring more than
72 minutes to execute.
\end{deffun}


\begin{deffun}{void}{reset\_time}{(void)}
  Resets the time counter. To be called just before executing some code to be
  profiled.
\end{deffun}

\begin{deffun}{long}{elapsed\_time}{(void)}
  Returns the time in milliseconds elapsed between now and the last call to
  \code{reset\_time()}.
\end{deffun}

\begin{deffun}{long}{intermediate\_elapsed\_time}{(void)}
  Returns the time in milliseconds elapsed between now and the last call to
  \code{reset\_time()} or to \code{intermediate\_elapsed\_time()}.
\end{deffun}

Here is a typical example of use of the timing functions:

\begin{example}
int main(void)
{
  reset_time();
  for (unsigned int i=0;i<1000;++i) {
     @emph^[Some time consuming operations]~
  }
  cout << "Elapsed time: " << elapsed_time() << " ms." << endl;
  return 0;
}
\end{example}

\section{The \code{timepiece} class}\label{sec:timepiece}
%%----------------------------------

A \code{timepiece} object allows to keep track of the time spent to perform
a particular task. Since the counter used is local to the object, it is possible
to monitor more than one such process.

\subsection{Methods of the \code{timepiece} class}
%%................................................

\begin{defmethod}{timepiece}{void}{start}{(void)}
  Starts the timepiece.
\end{defmethod}

\begin{defmethod}{timepiece}{void}{stop}{(void)}
  Stops the timepiece and accumulates the time spent since the last call to \code{start()}.
\end{defmethod}

\begin{defmethod}{timepiece}{void}{reset}{(void)}
  Resets to zero the counter keeping track of the total time the timepiece was running.
\end{defmethod}

\begin{defmethod}{timepiece}{long}{get\_total\_time}{(void) const}
  Returns the total amount of time the timepiece was running (time between calls to
  \code{start()} and \code{stop()}). The timepiece shall have been stopped by calling the
  \code{stop()} method before calling this one.
\end{defmethod}

\begin{defmethod}{timepiece}{long}{get\_intermediate\_time}{(void) const}
  Returns the amount of time spent since the last call to \code{start()}.
\end{defmethod}


\begin{example}
int main(void)
{
  timepiece t;
  t.start();
  for (unsigned int i=0;i<1000;++i) {
     @emph^[Some time consuming operations]~
     cout << "Intermediate time: "
          << t.get_intermediate_time() << " ms." << endl;
  }
  t.stop();
  cout << "Elapsed time: " << t.get_total_time() << " ms." << endl;
  return 0;
}
\end{example}


\chapter{Additional Documentation}\label{chap:documentation}
%%================================

\section{Documentation on GAOL}
%%-----------------------------

The primary reference is this manual. The directory \file{doc} of the sources
holds the pages that give the details of GAOL v5, with the measures behind its
choices:

\begin{itemize}
\item \file{doc/building.md}: the CMake, autotools and meson builds, and their options;
\item \file{doc/using.md}: the flags of interval arithmetic, GAOL from CMake
  and from pkg-config, the namespaces and the rounding direction;
\item \file{doc/three-builds.md}: what the three builds agree on, what the options
  are worth, and the compilers and options they refuse;
\item \file{doc/tests.md}: what the programs of \file{tests} check;
\item \file{doc/accuracy.md}: the tightness of each operation, with its algorithm
  and what the tests found;
\item \file{doc/differences.md}: each change GAOL v5 brings to GAOL 4.2.2, and
  where it comes from;
\item \file{doc/continuous-integration.md}: the systems, processors and compilers
  GAOL is built and tested on;
\item \file{doc/compare/README.md}: GAOL compared with libieeep1788, filib++,
  PROFIL/BIAS and Solaris Studio, their special cases against IEEE 1788-2015 and
  their performance;
\item \file{3rd/README.md}: the version of CORE-MATH GAOL takes, and what differs
  from it.
\end{itemize}

\noindent
There is also an HTML reference for the code itself, generated by doxygen,
which might be of interest only to developers seeking to understand or modify
GAOL. The manual of GAOL~4, by Fr\'ed\'eric Goualard, is kept in
\file{manual/v4}.

\section{References}
%%------------------
The following articles, books and standards have inspired in some way or another the
devising of the GAOL library and/or the writing of this manual.

\begin{itemize}
\item \emph{IEEE Standard for Interval Arithmetic}. IEEE Std 1788-2015~\cite{IEEE1788:2015}.
\item \emph{The CORE-MATH Project}. Alexei Sibidanov, Paul Zimmermann and
  St\'ephane Glondu~\cite{CoreMath:2022}, and \url{https://core-math.gitlabpages.inria.fr/}.
\item \emph{How do you compute the midpoint of an interval?} Fr\'ed\'eric
  Goualard~\cite{Goualard:TOMS2014}.
\item \emph{Fast and Correct SIMD Algorithms for Interval Arithmetic}. Fr\'ed\'eric
  Goualard~\cite{Goualard:PARA08}.
\item \emph{Interval Extensions of Multivalued Inverse Functions}. Fr\'ed\'eric
  Goualard~\cite{Goualard:2008}.
\item \emph{Interval Arithmetic Specification}. Dmitri Chiriaev and G.\:
William Walster. Draft revised May 1998.
\item \emph{The Extended Real Interval System}. G.\: William
Walster. April 1998.
\item \emph{C++ Interval Arithmetic Programming Reference}. Sun Microsystems,
Inc. October 2000, revision A.
\item \emph{Interval Arithmetic: From Principles to Implementation}. T.\:
Hickey, Q.\: Ju, and M.\: H.\: van Emden~\cite{Hickey-Ju-vanEmden:2001}.
\end{itemize}

\chapter{Reporting bugs}
%%======================

Bugs and suggestions for improvement of GAOL v5 shall be submitted as issues of
its repository:

\begin{center}
  \url{https://github.com/Jordan08/GAOL/issues}
\end{center}

\noindent
A report is most useful with the program that shows the problem, the interval
it gives and the one expected, the system, the processor, the compiler and the
options GAOL was built with.

\chapter{Contributors}
%%====================

The main implementor and lead designer for the GAOL library is Fr\'ed\'eric Goualard
(\url{https://frederic.goualard.net}).

GAOL v5 is written by Jordan Ninin (ENSTA), for Codac.

The \code{interval pow(const interval\&, const interval\&)} function
was designed by Marc Christie.

The code for the multiplication and the division is largely inspired
from the one presented by Tim Hickey, Qun Ju and Maarten Van Emden in
\emph{Interval Arithmetic: from Principles to Implementation}, Journal
of the ACM 48(5):1038--1068, september 2001.

GAOL v5 takes up:
\begin{itemize}
\item the patch IBEX applies to GAOL, by Gilles Chabert (initialization of the
  floating-point unit, rounding direction of \code{midpoint()} and
  \code{operator<{}<});
\item the fork of GAOL by Fabrice Le Bars, for Visual C++, MinGW and the systems
  the \cmd{configure} of GAOL does not know;
\item the CMake build of GAOL in IBEX, by Cyril Bouvier and Gilles Chabert, with
  the compilation flags of the IBEX fork of Fabrice Le Bars;
\item the rounding tests of Codac, which its tests follow.
\end{itemize}

\noindent
Its elementary functions are those of CORE-MATH, by Alexei Sibidanov, Paul
Zimmermann and the other authors of the project (Inria).

\chapter{Moving a program from GAOL 4 to GAOL v5}\label{chap:changes}
%%==============================================
\cindex{GAOL 4!moving from}%

Most programs written for GAOL~4 compile and run with GAOL v5 once the
namespace is opened. This chapter lists what a program may have to change, and
what it will see differently. The page \file{doc/differences.md} gives each
change with its reasons and its measures.

\section{Compiling and linking}
%%-----------------------------

\begin{itemize}
\item \file{gaol/gaol.h} no longer opens the namespace \code{gaol}: add
  \code{using namespace gaol;}, or name \code{gaol::interval}
  (\pxref{sec:namespaces}). The type and the functions are in
  \code{gaol\_core}, and the symbols of the library name it: a program compiled
  against the headers of GAOL~4 is compiled again.
\item GAOL is the only library to link: the mathematical library GAOL~4 was
  built with is no longer used, CORE-MATH being compiled into GAOL, and the
  options of the builds that chose it are gone.
\item The code including the headers of GAOL is compiled with the flags of
  interval arithmetic, which \code{gaol::gaol} and \file{gaol.pc} give
  (\pxref{sec:compiling}); \file{gaol/gaol\_config.h} refuses the options and the
  compilers that give wrong bounds (\pxref{sec:refused}).
\item The intervals of floats, \code{gaol::intervalf} and
  \code{gaol::interval2f}, are compiled only when asked for
  (\option{GAOL\_FLOAT\_INTERVALS}).
\item \code{gaol::init()} need not be called again after the program changed the
  rounding direction: each operation sets it upward itself
  (\pxref{sec:rounding-direction}).
\item The operators \code{==} and \code{!=} on intervals are gone, with the
  possibly relations (\code{possibly\_eq()}, \code{possibly\_le()}\ldots),
  \code{certainly\_eq()} and \code{certainly\_neq()}, and so is the option
  \loption{enable-relations}, which \cmd{configure} and meson refuse: a program
  compares intervals with \code{set\_eq()}, \code{set\_disjoint()} and
  \code{!set\_disjoint()}, and \code{<}, \code{<=}, \code{>} and \code{>=} are
  the certainly relations (\pxref{sec:interval-relations}).
\end{itemize}

\section{Results that change}
%%---------------------------

\begin{itemize}
\item \textbf{The bounds are tighter}: the elementary functions give the tightest
  bounds, rather than one double beyond, and are exact where their value is a
  double; the integer powers and the $n$th roots are computed from exact
  products; \code{sin}, \code{cos} and \code{tan} are within one double of the
  tightest bounds at every magnitude (Chapter~\ref{chap:accuracy}).
\item \textbf{The empty set.} The constructors give it for
  \itv{+\infty}{+\infty}, \itv{-\infty}{-\infty}, bounds in the wrong order and
  NaN bounds, and assigning an infinity to an interval empties it
  (\pxref{sec:constructors}); the intersection of disjoint intervals is it,
  with NaN bounds; so are \code{log} of an interval holding no positive number,
  \code{invabs\_rel} of a negative interval and \code{div\_rel} with no
  solution in \var{I}. \code{width()} of the empty set is NaN rather than $-1$.
\item \textbf{Powers.} \code{pow(x, p)} for a double \var{p} no longer truncates
  \var{p}; \code{pow(x, y)} keeps the negative part of \var{x} for an integer
  exponent only, and takes the \ieee{pow} of IEEE 1788-2015 otherwise
  (\pxref{sec:powers}). \code{nth\_root} of an odd order is defined on the whole
  real line, and takes a negative order.
\item \textbf{The comparisons} follow IEEE 1788-2015 for the empty set and the
  infinite bounds: \code{certainly\_leq()} and its companions are true when
  either interval is empty, and \code{set\_strictly\_contains()} takes an
  infinite bound as beyond itself (\pxref{sec:interval-relations}).
\item \textbf{Input and output.} Numbers are read exactly, and the literals of
  IEEE 1788-2015 are read, in any case of letters; \code{"[inf]"} is the empty
  set. Every string is an expression, where an interval may stand wherever a
  number may (\code{"1+[1,2]"}), an uncertain number being no bound
  (\code{"[5?1]"} is refused), and an error anywhere stops the reading. The
  reader knows every function of GAOL, and \code{gaol\_ieee1788::textToInterval}
  reads the names of IEEE 1788-2015 (\code{pown}, \code{rootn}\ldots). The format \code{hexa} writes interval literals, which are read back bit
  for bit; the decimal formats round the bounds outward whatever the C
  library, and the format \code{agreeing} keeps the digits of each bound
  (Chapter~\ref{chap:io}).
\item \textbf{New operations}: \code{atan2()} is implemented, and the functions
  of IEEE 1788-2015 GAOL did not provide are added (\code{exp2}, \code{hypot},
  \code{sign}, \code{trunc}\ldots, Chapter~\ref{chap:arithmetic}).
\end{itemize}


\schapter{Library Copying}\label{chap:copying}
%%========================
\input{lgpl}

\section*{The license of CORE-MATH}
%%---------------------------------

The sources of CORE-MATH, in \file{3rd/math-core}, which GAOL compiles into its
library, are distributed under the following license (MIT license):

\begin{verbatim}
Permission is hereby granted, free of charge, to any person
obtaining a copy of this software and associated documentation
files (the "Software"), to deal in the Software without
restriction, including without limitation the rights to use,
copy, modify, merge, publish, distribute, sublicense, and/or
sell copies of the Software, and to permit persons to whom the
Software is furnished to do so, subject to the following
conditions:

The above copyright notice and this permission notice shall be
included in all copies or substantial portions of the Software.

THE SOFTWARE IS PROVIDED "AS IS", WITHOUT WARRANTY OF ANY KIND,
EXPRESS OR IMPLIED, INCLUDING BUT NOT LIMITED TO THE WARRANTIES
OF MERCHANTABILITY, FITNESS FOR A PARTICULAR PURPOSE AND
NONINFRINGEMENT. IN NO EVENT SHALL THE AUTHORS OR COPYRIGHT
HOLDERS BE LIABLE FOR ANY CLAIM, DAMAGES OR OTHER LIABILITY,
WHETHER IN AN ACTION OF CONTRACT, TORT OR OTHERWISE, ARISING
FROM, OUT OF OR IN CONNECTION WITH THE SOFTWARE OR THE USE OR
OTHER DEALINGS IN THE SOFTWARE.
\end{verbatim}


\bibliographystyle{plain}
\bibliography{gaol}

\printindex[idx]

\end{document}

%%% Local Variables:
%%% mode: latex
%%% TeX-master: t
%%% End:
